% concatenated from main.tex, Chapter_Introduction/Main.tex, Chapter_Problem_Construction/Main.tex, Chapter_Main_Theorems/Main.tex, Chapter_Preparatory_Results/Main.tex, Chapter_Main_Proofs/Main.tex, Chapter_Characterisation_and_Local_uniqueness/Main.tex, Chapter_Conclusion/Main.tex, Chapter_Acknowledgements/Main.tex, Chapter_Alexander_Basis_Exposition/Main.tex
\documentclass[11pt]{article}
\usepackage{graphicx} 
\usepackage{amssymb}
\usepackage{amsmath}
\usepackage{bm,cancel}
\usepackage{amsthm}
\usepackage{parskip}
\usepackage{geometry} 
\geometry{a4paper, total={170mm, 257mm}, left=20mm, top=20mm}
\usepackage{wrapfig, framed}
\usepackage{overpic}
\usepackage{comment}
\usepackage{mathtools, mathrsfs}
\usepackage{scalerel}
\usepackage{relsize}
\usepackage{hyperref}
\usepackage{graphbox}
\usepackage{authblk}
\usepackage{stmaryrd}
\usepackage[T1]{fontenc}
\usepackage{enumitem}
\usepackage{tikz-cd} 

\DeclareFontEncoding{LS1}{}{}
\DeclareFontSubstitution{LS1}{stix}{m}{n}
\newcommand*\kay{%
  \text{%
  \fontencoding{LS1}%
  \fontfamily{stixscr}%
  \fontseries{\textmathversion}%
  \fontshape{n}%
  \selectfont\symbol{"6B}}}
\makeatletter
  % the current math version is saved in \math@version
  \newcommand*\textmathversion{\csname textmv@\math@version\endcsname}
  \newcommand*\textmv@normal{m}
  \newcommand*\textmv@bold{b}
\makeatother

\newcounter{descriptcount}
\newlist{enumdescript}{description}{1}
\setlist[enumdescript,1]{%
  before={\setcounter{descriptcount}{0}%
          \renewcommand*\thedescriptcount{\boldsymbol{\alpha}rabic{descriptcount}}},
        font={\bfseries\stepcounter{descriptcount}Aim \thedescriptcount:}
}
\usepackage{tikz}

\newcommand{\foo}[1]{%
\begin{tikzpicture}[#1]%
\draw (0,0) -- (1ex,1ex);%
\draw (0.5ex,0) -- (1.3ex,0.8ex);%
\end{tikzpicture}%
}

% Theorem Environments
\theoremstyle{plain}
\newtheorem{theorem}{Theorem}[section]
\newtheorem*{theorem*}{Theorem}
\newtheorem{proposition}[theorem]{Proposition}
\newtheorem*{proposition*}{Proposition}
\newtheorem{corollary}[theorem]{Corollary}
\newtheorem{lemma}[theorem]{Lemma}
\newtheorem{conjecture}[theorem]{Conjecture}
\theoremstyle{definition}
\newtheorem{definition}[theorem]{Definition}
\newtheorem{question}[theorem]{Question}
% \newtheorem{example}{Example}
\theoremstyle{remark}
%\newtheorem{example}[theorem]{Example}
\newtheorem{remark}[theorem]{Remark}

\usepackage{chngcntr}
\newtheorem{exampleinner}{Example}[section] 
\counterwithin{exampleinner}{section} 
\newenvironment{example}
  {\begin{list}{}{\setlength{\leftmargin}{2em} \setlength{\rightmargin}{0em}}
   \item\relax\begin{exampleinner}}
  {\end{exampleinner}\end{list}}

\newenvironment{rightbox}[1]
 {\itemize[
    %nosep,
    leftmargin=\dimexpr\textwidth-#1\relax,
    rightmargin=0pt,
    itemindent=\parindent,
    listparindent=\parindent,
  ]\item[]\relax}
 {\enditemize}

 
\newcommand{\TightBracketB}[1]{\raisebox{0.00ex}{\scalebox{1}{${#1}$}}}%
\newcommand*{\eq}[1]{\TightBracketB[#1\TightBracketB]}

\newcommand{\R}{\mathbb{R}}
\newcommand{\ssubset}{\subset\joinrel \! \!\subset}
\newcommand{\B}{\mathcal{B}}
\newcommand{\lie}{\ensuremath{\mathsterling}}
\newcommand{\Harm}{\text{Harm}}
\newcommand{\invcurl}{\text{curl}^{-1}}
\newcommand{\icfree}{\invcurl_{\text{free}}}
\newcommand{\icfreen}{\invcurl_{\text{free,N}}}
\newcommand{\proj}{\mathcal{F}}
\newcommand{\curl}{\text{curl}\,}
\newcommand{\Div}{\text{div}\,}
\newcommand{\FF}{\mathcal{F}_{\text{free}}}
\newcommand{\ec}{\text{ç}}
\newcommand{\bb}{\mathbf{b}}
\newcommand{\ST}{\Lambda}
\newcommand{\HT}{\mathrm{HT}}
\newcommand{\D}{\text{Emb}(\cup_{i=1}^n \ST_i)}
\newcommand{\C}{\mathcal{C} (\ST_i)}
\newcommand{\I}{\mathcal{I}(\ST)}
\newcommand{\h}{\mathcal{H}_D^2(\ST_i)}
\renewcommand{\hbar}{h}
\newcommand{\Hfree}{\mathcal{H}_{\text{free}}}
\renewcommand{\j}{\text{\reflectbox{$\iota$}}}
\makeatletter         
\renewcommand\maketitle{
{\raggedright 
\begin{center}
{\vspace{-5em} \huge \@title }\\[2ex] 
{\Large  \@author}\\[2ex]

\end{center}}} 
\makeatother

\title{A Variational Shape Optimisation Approach to Multi-region Relaxed Magnetohydrodynamic Equilibria}
\author[1]{K. de Lacy}
\author[1]{L. Noakes}
\author[1]{D. Pfefferl\'e}
\affil[1]{ \small The University of Western Australia, 35 Stirling Highway, Crawley WA 6009, Australia}

\date{\today}

\begin{document}
\maketitle
\newsavebox\mybox
\savebox{\mybox}{$\stretchrel*{|}{\rule[-.6ex]{0ex}{2.5ex}}$}
\def\myabs#1{\usebox\mybox#1\usebox\mybox}
\vspace{2em}


\begin{abstract}
\noindent
Let \(\ST\subset\mathbb{R}^3\) be a region admitting a partition into \(n\) compact, connected subregions \(\ST_1,\dots,\ST_n\), each with smooth boundary. Consider a vector field \(B\) on \(\ST\) where \(B|_{\ST_i}\) is smooth, divergence free, and tangent to \(\partial \ST_i\) for all \(i\). We show that the multi-region relaxed magnetohydrodynamics (MRxMHD) equilibrium equations are necessary and sufficient conditions for \( B \) and a metric to yield a stationary point of the magnetic energy under appropriate constraints. We constrain the pressure, relative helicity, and magnetic flux of \(B\) through all smooth surfaces in \(\ST_i\) whose boundary lies on \(\partial \ST_i\).

\noindent
We identify a previously overlooked gauge condition. A definition for relative helicity is introduced, its gauge invariance is proved, and the existence of a gauge where relative helicity reduces to conventional helicity is demonstrated. In the case of a single region an additional condition is introduced that is sufficient to ensure a critical point of the magnetic energy is also a minimiser.
\end{abstract}

\section{Introduction}\label{section:Introduction}




Magnetohydrodynamics (MHD) equilibrium equations describe the steady-state behaviour of a globally neutral plasma under the influence of a magnetic field. Formally, MHD is derived from the Vlasov–Maxwell equations by assuming quasi-neutrality and taking low-order velocity moments to produce a single-fluid model coupled to the magnetic field. This model neglects effects like collisionless damping, viscosity, and high-frequency waves \cite{Freidberg_2014}. The MHD equilibrium equations characterise steady-state MHD solutions and take the form
\begin{align*}
  (\nabla \times B) \times B = \nabla p \, , \quad \nabla \cdot B = 0 \, , \quad N \cdot B \big|_{\partial \ST} = 0 \, ,
\end{align*}
where the \emph{magnetic field} \(B\) is a smooth vector field on a compact, oriented domain \(\ST\subset\mathbb{R}^3\) with smooth boundary \(\partial \ST\). The \emph{pressure} \(p : \ST \to \R\) is a smooth function that is constant on each connected component of the boundary \(\partial \ST\). 


Let \(f:\ST \to f(\ST)\) be a diffeomorphism isotopic to the identity with \(f|_{\partial \ST} = \mathrm{id}\). Rather than work directly on a moving domain \(f(\ST)\) with a Euclidean volume form \(\varpi\), we pull back to a fixed reference domain \(\ST\). We equip \(\ST\) with a volume form \(f^*\varpi\), and interpret differential operators and inner products, such as those in the MHD equation, with respect to the pulled-back Euclidean metric. The domain \(\ST\) is fixed while \(f\) and \(B\) are allowed to vary. The outward unit normal \(N\) on \(\partial \ST\) is defined as the pullback of the Euclidean outward normal on \(\partial f(\ST)\). This keeps the reference domain fixed while allowing the geometry and the location of interior conditions to vary through \(f\).


The MHD equilibrium equations constitute a nonlinear boundary‐value problem. MHD equilibria are used in the early stages of designing and optimising magnetic field configurations for plasma confinement devices \cite{chen1984introduction, hazeltine2003plasma}. By exploring domain geometries one can identify promising reactor shapes and parameters before undertaking more realistic but computationally expensive simulations. Both iterative solvers and variational methods are in common use for computing or approximating MHD equilibria (assuming such solutions exist) with variational formulations typically exhibiting better convergence properties \cite{loizu2016verification, REIMAN1986157, hirshman2011siesta, helander2013ideal, dudt2020desc, hirshman1986three, chodura19813d, hayashi1995evolution}. Restricting to axisymmetric fields invariant under continuous rotations about the toroidal axis yields the 2D elliptic Grad–Shafranov equation, for which many numerical solvers exist \cite{PATAKI201328, LEE201572, zakharov1999theory, takeda1991computation}.



This work presents a variational formulation of the MHD equilibrium equations on relatively general domains \(\ST\). Specifically, we show that a variational method developed for nested toroidal domains \cite{dennis2014multi}---which can approximate MHD equilibria under certain assumptions \cite{10.1063/1.4795739}---can be extended to more general domains \(\ST\). To properly discuss the motivating ideas we introduce the necessary notation and provide a brief overview of the relevant literature for the case of non-degenerate pressure profiles.


\subsection{Setup}\label{sub:setup}

Let \(\Omega^k(\ST)\) be the space of all smooth \(k\)-forms on \(\ST\). The magnetic field $B$ and the volume form $f^*\varpi$ yield a smooth 2-form,
\begin{align*}
   \beta = i_B ( f^*\varpi ) \, ,
\end{align*}
where \(i_B : \Omega^k(\ST) \to \Omega^{k-1}(\ST)\) denotes interior multiplication by \(B\). The MHD equilibrium equations are equivalent to:
\begin{align} \label{eq:mhe_original}
  \delta \beta  \wedge \star \beta  &= \star dp \, , \quad
  d \beta  = 0 \, , \quad \mathcal{J}^* \beta  = 0 \, ,
\end{align}
where $d:\Omega^k(\ST)\to \Omega^{k+1}(\ST)$ is the \emph{exterior derivative}, $\star:\Omega^k(\ST)\to\Omega^{3-k}(\ST)$ is the \emph{Hodge star} operator, and the \emph{co-derivative} \(\delta\) applied to a \(k\)-form on \(\ST\) is defined as \(\delta \coloneqq (-1)^k \star d \star\) \cite{lee}. Additionally, \(\wedge\) denotes a wedge product. The Neumann condition on \(B\) is replaced by \(\mathcal{J}^* \beta  = 0\) where \(\mathcal{J}\colon \partial\ST \hookrightarrow \ST\) is the inclusion map and \(\mathcal{J}^*\) its pull‐back. We reserve the symbol \(\mathcal{J}\) for inclusion maps, specifying its domain and codomain only when the context does not make them clear. Likewise in integrals we often omit the explicit insertion of \(\mathcal{J}^*\) whenever the integration domain alone suffices to indicate the pull‐back. 





For \(\omega_1, \omega_2 \in \Omega^k (\ST)\) the \(L^2\) inner product and norm are given by:
\begin{align*}
\begin{gathered}
  \langle \omega_1, \omega_2 \rangle_{L^2} \coloneqq \int_{\ST} \langle \omega_1 ,\omega_2 \rangle \, \varpi = \int_{\ST} \omega_1 \wedge \star \omega_2 \, , \qquad
  \| \omega_1 \|_{L^2}^2 \coloneqq \langle \omega_1, \omega_1 \rangle_{L^2} \, ,
\end{gathered}
\end{align*}
where \(\langle \cdot , \cdot \rangle\) denotes the pointwise inner product induced by the Riemannian metric on \(\ST\).




As \(\ST\) is a compact, oriented subset of \(\R^3\) with a smooth boundary, a closed 2-form \(\beta  \in  \Omega^2(\ST)\) that is zero when pulled back by inclusion to the boundary, is automatically exact. Namely, there exists an \(a \in  \Omega^1(\ST) \) where \(\beta  = da\), as shown by Cantarella and Parsley (Proposition C.4 \cite{CANTARELLA20101127}). As we use this result throughout the paper, an exposition of this proof can be found in the appendix, Lemma \ref{Lemma:Potential}. 



Consider any closed \( s \in \Omega^1(\ST) \) and fix a potential \(a\) of \(\beta\). Since \( ds = 0 \), we have \( \beta  = d(a + s) \), so \( a + s \) is likewise a valid primitive for \( \beta  \). Similarly, the reverse argument shows that: for any two primitives of \(\beta \) their difference is a closed 1-form, demonstrating that \(\{a+s \in \Omega^1(\ST) : \; da = \beta  ,\; ds=0\} \) is the set of every potential for \(\beta \). 


Consider a set \(G \subseteq \Omega^1(\ST)\) of potentials with the following property: for every closed \(2\)-form \(\beta \) whose pullback to \(\partial \ST\) vanishes, there exists at least one \(a \in G\) such that \(\mathrm{d}a = \beta \), and whenever \(a_1,a_2 \in G\) are both potentials for \(\beta \), one has
\[
\int_{\ST} a_1 \wedge \beta  = \int_{\ST} a_2 \wedge \beta  \, .
\]
The \emph{helicity} of \(\beta \) on \(\ST\) with respect to \(G\) is defined by
\begin{align*}
  H_G(\beta ,\ST) \coloneqq \int_{\ST} a \wedge \beta  \, ,
\end{align*}
where \(a \in G\) is a potential of \(\beta \).

Without a restriction on the admissible potentials, the helicity need not be well-defined: if the flux of \(\beta \) through some smooth surface in \(\ST\) whose boundary lies in \(\partial \ST\) is non-trivial, then different choices of potential \(a\) lead to arbitrarily different values of the helicity (see Subsection \ref{sub:Well-Posed}). In Subsection \ref{subsection:Relative Helicity} we present a replacement for helicity that does not require restricting potentials to \(G\). Another approach, not adopted here but common in the literature, is to use a Biot--Savart type operator to select a gauge in which helicity becomes a well-defined function of the magnetic field and the underlying geometry \cite{parsley2001taylor, Arnold_Topological_Methods, soperimetric_problems_helicity, CANTARELLA20101127, MacTaggart_2023}.





Consider a partition of \(\ST\) into \(n<\infty\) compact, connected subregions \(\ST_i\) (\(i=1,\dots,n\)), each with smooth boundary, so that
\[
    \bigcup_{i=1}^n \ST_i = \ST \, ,
    \qquad \text{and} \qquad
    \mathrm{int}(\ST_i)\cap \mathrm{int}(\ST_j)=\emptyset
    \quad (i\neq j) \, .
\]
Let the set \(\I\) be the set of all diffeomorphisms \(f: \ST\to \ST\) such that \(f(\partial \ST) = \partial \ST\) pointwise and there exists a smooth isotopy \(f_t : \ST\to \ST\) for \(t \in [0,1]\) that is an orientation preserving diffeomorphism for all \(t\) and \(f_0 = id\), \(f_1 = f\). Therefore, the family of sets \(f_t (\ST_i)\) gives a smooth deformation between the partitions \(\{\ST_i\}_{i=1}^n \) and \( \{ f( \ST_i)\}_{i=1}^n\), with the subregions remaining pairwise disjoint and covering \(\ST\) for every \(t\).



For each \(f \in \I\), the image \(f(\ST_i)\) inherits a Euclidean metric and a corresponding volume form \(\varpi\). We then equip \(\ST_i\) with the pullback metric induced by \(f\), so that its volume form is \(f^{*}\varpi\). The associated Hodge star operator \(\star\) on \(\ST_i\) therefore depends on \(f\), although we suppress this dependence in the notation.


From this point onward, we permit \(B\) to be discontinuous across partition boundaries, requiring only that its restriction to each \(\mathrm{int}(\ST_i)\) is smooth. We define a smooth field \(B_i \in \Gamma \ST_i\) that matches \(B\) on \(\mathrm{int}(\ST_i)\). The process of replacing a smooth magnetic field with a piecewise smooth field is called \emph{relaxation}. Similarly, we allow the pressure \(p\) to be discontinuous across partition boundaries. 



For \(n>1\), a partition \((\ST_1,\dots,\ST_n)\) specifies the locations where \(B\) may be discontinuous. In each region \(\ST_i\) we express the vector field \(B_i\) in the language of exterior calculus by defining
\[
  \beta_i \coloneqq i_{B_i}\bigl(f^*\varpi\bigr)
  \in \Omega^2(\ST_i) \, ,
\]
for \(i=1,\dots,n\). If \(n=1\) then \(B\) is a smooth vector field and \(\beta _1\) reduces to the definition of \(\beta \).


Each non-empty intersection \(\ST_i \cap \ST_j\) for \(i \neq j\) is called an \textit{interface}. There are a finite number of pairwise distinct interfaces labelled \(I_q\) for \(q = 1, 2, \dots, m\). For a given \(I_q\) there exists, by definition, an \(i\) and \(j\) where \(i > j\) such that \(I_q = \ST_i \cap \ST_j\). For \(k\)-forms \(\omega_i \in \Omega^k(\ST_i)\), \(\omega_j \in \Omega^k(\ST_j)\) define the operator \(\llbracket \omega \rrbracket |_{I_q} = \llbracket \omega \rrbracket |_{\ST_i \cap \ST_j} \coloneqq \mathcal{J}^* ( \omega_{i}) - \mathcal{J}^* (\omega_{j} ) \). Here each \(\mathcal{J}\) is an inclusion map \(I_q \hookrightarrow \ST_i\) and \(I_q \hookrightarrow \ST_j\) respectively. 


With an inclusion map \(\mathcal{J} : \partial  \ST_i \hookrightarrow  \ST_i\) we define the \textit{tangential trace} of \(\omega \in \Omega^k(\ST_i) \) on \(\partial  \ST_i\) by \(t \omega = \mathcal{J}^* \omega\). Let \(N\) be a vector field on \(\ST_i\) that is a unit normal field to the surface \( \partial  \ST_i \). The \textit{normal trace} is defined to be the \(k-1\) form \(n \omega = \mathcal{J}^* ( i_N \omega )\).


We define the following sets of \(k\)-forms in \(\ST_i\):
\begin{align*}
   \Omega^k_D( \ST_i ) &\coloneqq \{  \beta  \in   \Omega^k( \ST_i )  : \; t\beta  = 0 \} \, , \\
   \Omega^k_N( \ST_i ) &\coloneqq \{  \beta  \in   \Omega^k( \ST_i )  : \; n \beta  = 0 \} \, , \\
  \C &\coloneqq  \{  \beta  \in   \Omega^2( \ST_i )  : \; t \beta  = 0 \, , \; d \beta  = 0 \} \, .
\end{align*}
Elements of \(\Omega^k_D( \ST_i )\) and \( \Omega^k_N( \ST_i )\) are called Dirichlet and Neumann \(k\)-forms respectively. 





For a given partition \((\ST_1, \dots,\ST_n)\) and piecewise-constant pressure \(p : \ST \to \R\) equal to \(p_i\) in \(\ST_i\), a solution to the multi-region relaxed magnetohydrodynamics (MRxMHD) equilibrium equations is given by \(\beta _i\in \C \), \(i = 1, \dots, n\) and \(f \in \I\) solving:
\begin{align}
\label{prob:MRxMHD} 
  d \star \beta _i = \mu_i \beta _i \, , \qquad 
  \big\llbracket \| \beta  \|^2 + 2p \big\rrbracket \big|_{I_j} = 0 \, , 
\end{align}
for each \(i\), \(j\) and constants \(\mu_i \in \R\). These conditions on \(f\) and each \(\beta _i\) are equivalent to first order necessary and sufficient conditions for a solution to a variational problem (Theorem \ref{thr:iff_simple}). In equations \eqref{prob:MRxMHD} the contribution from \(f\) lies in the Hodge star and norm. The first equation is a Beltrami equation, and the second is referred to as a \textit{jump condition}. In vector notation:
\begin{align*}
  \nabla \times B_i = \mu_i B_i \, , \qquad 
  \big\llbracket \|B\|^2 + 2 p  \big\rrbracket \big|_{I_j} = 0 \, .
\end{align*}
If MRxMHD equilibria approximate MHD equilibria the fact that each pressure discontinuity in equation \eqref{prob:MRxMHD} can contribute to plasma confinement would make MRxMHD a natural framework for studying magnetic confinement \cite{10.1063/1.4765691}.




\subsection{Background: Toroidal Geometry}\label{sub:back}


\setlength{\FrameSep}{3pt}

\begin{wrapfigure}{r}{0.3\linewidth}
    \vspace{-1em}
  \begin{framed}\raggedleft
  \begin{overpic}[width=\linewidth]{Chapter_Introduction/image1.png}
    \put(63,23){\(\ST_1\)}
    \put(63,36.5){\(\ST_2\)}
    \put(63,45.5){\(\ST_3\)}
    \put(64.5,53){\(\vdots\)}
    \put(63,63){\(\ST_n\)}
  \end{overpic}
  \caption{A cross-section and partition of a solid torus embedded in \(\R^3\).}\label{fig:torus} 
  \end{framed}
  \vspace{-2em}
\end{wrapfigure}

In this subsection, we restrict ourselves to a special geometric setting and review previous work. The existing literature addresses minimisation problems, but more generally we study stationary points.


Toroidal geometries are common when solving for MHD equilibria. Consider a smooth pressure \(p\) that is non‐degenerate (\(\nabla p\neq0\)) on \(\partial \ST\). Then an MHD solution satisfies \(B \cdot \nabla p = 0\), namely, \(B\) is tangential to the corresponding pressure level set.



On level sets of \(p\) with \(\nabla p \neq 0\) everywhere, we have \(\nabla p = (\nabla \times B) \times B\), so \(B\) and \(\nabla \times B\) are both non-zero and \(B\) is nowhere parallel to \(\nabla \times B\). The Poincaré--Hopf theorem implies that any compact manifold admitting a nowhere-zero continuous tangent vector field has Euler characteristic zero. The only compact, connected, smooth surface with Euler characteristic zero is diffeomorphic to a 2-torus \cite{Arnold_Topological_Methods}. Consequently, if the pressure is non-degenerate then each level set is a disjoint union of 2-tori.


In this subsection we consider regions \(\ST_i\) for \(i = 1, \dots, n\) that are solid and hollow tori embedded in \(\R^3\), as in Figure \ref{fig:torus}. This implies that interfaces are given by: \(I_k = \partial \cup_{i=1}^k\ST_i = \partial \ST_{k} \cap \partial \ST_{k+1}  \) for \(k=1,2, \dots, n-1\), which are pairwise-disjoint toroidal surfaces.




For a fixed set \(G\) as defined previously, consider the following minimisation problem over closed Dirichlet 2-forms \(\beta _i \in \C \) and diffeomorphisms \(f \in \I\) given by:
\begin{align} \label{prob:min1}
\begin{gathered}
    \min_{\substack{f \in \I \\ \beta _i \in \C \, , \; \forall i }} \sum_{i=1}^n \Big( \frac{1}{2} \| \beta _i\|_{L^2}^2 - p_i \myabs{\ST_i} \Big) \, , \\ H_G ( \beta_i , \ST_i ) = h_i  \, , \text{\footnotemark}
  \quad \int_{T_{i,j}} \beta _i = \psi_{i,j} \, , 
\end{gathered}
\end{align}
\footnotetext{Although a formulation in terms of the potential exists, we state a problem with \(\beta _i\) to match the literature. With adjustments the variational problem is well-posed in \(\beta _i\) and the two formulations are equivalent.}
for \(i = 1, \dots, n\), \(j = 1,2\), \((i,j) \neq (1,2)\), where \(f\) affects each term in the objective function through the metric of \(\Lambda_i\), which is the Euclidean metric pulled back by \(f^*\). Namely, given constants \((p_i, h_i, \psi_{i,j})\) and regions \(\ST_i\) for each \(i,j\), our objective is to determine \((\beta _1, \dots, b_n, f)\) that solve \eqref{prob:min1}. A pair \((\beta _i, f)\) determine a magnetic field via \(B_i = (\star  \beta _i)^\flat\), where \(\flat\) is the usual flat musical isomorphism. Additionally, note that \(\myabs{\ST_i}\) is the volume of \(\ST_i\).


In problem \eqref{prob:min1}, \(T_{i,1}\) is a poloidal surface through the solid or hollow tori \(\ST_i\) for all \(i = 1,2, \dots, n\). Similarly, \(T_{i,2}\) is a toroidal surface through each hollow torus \(\ST_i\) for \(i = 2, \dots, n\). The generalisation of such surfaces is explained in subsection \ref{sub:flux}. 


Pushing forward each integral in problem \eqref{prob:min1} and using a result that 
\[
H_G ( \beta_i , \ST_i ) = H_{(f^{-1})^* G} (  (f^{-1})^*  \beta_i , f(\ST_i) ) \, ,
\]
from Cantarella \cite{CANTARELLA20101127}, we have a minimisation problem with a Euclidean metric that is equivalent to \eqref{prob:min1}. For 2-forms \(\omega_i \in \mathcal{C} ( f(\ST_i) )\), with an appropriate \(G\), problem \eqref{prob:min1} is rewritten as:
\begin{align*} 
\begin{gathered}
\min_{\substack{f \in \I \\  \omega_i \in \mathcal{C} ( f(\ST_i) ) \, , \; \forall i }} \sum_{i=1}^n \Big( \frac{1}{2} \|  \omega_i\|_{L^2}^2 - p_i \myabs{f(\ST_i)} \Big) \, , \\ H_{G} ( \omega_i , f(\ST_i) ) = h_i \, , 
  \quad \int_{f(T_{i,j})}  \omega_i = \psi_{i,j} \, .
\end{gathered}
\end{align*}
This formulation is equivalent to the minimisation problem of Dewar, Yoshida, Bhattacharjee, and Hudson \cite{Dewar_Yoshida_Bhattacharjee_Hudson_2015}. These authors also consider the case in which the vector field \(B\) is a Neumann harmonic vector field, in the sense of Schwarz \cite{schwarz-1995}, on regions near the domain boundary. They derive the MRxMHD equilibrium equations \eqref{prob:MRxMHD} from \eqref{prob:min1} in a fixed gauge via a Lagrange multiplier argument \cite{Dewar_Yoshida_Bhattacharjee_Hudson_2015}.





The Stepped Pressure Equilibrium Code (SPEC) is a numerical solver for the MRxMHD variational problem, capable of reconstructing stellarator magnetic fields with arbitrarily high accuracy. SPEC can compute MRxMHD equilibria that incorporate cross-helicity conditions, field-aligned flow and angular-momentum constraints \cite{hudson2020free, qu2020stepped}. The program is used to generate and analyse stellarator configurations that exhibit low neoclassical damping relative to other reactor designs \cite{loizu2016verification, helander2008intrinsic}.



Numerical studies using SPEC provide evidence that solutions to the MRxMHD variational problem \eqref{prob:min1} can approximate MHD equilibria \eqref{eq:mhe_original}. However, because the existence of MHD equilibria is not known in general, there is no a priori guarantee that one can choose constants \(\psi_{i,j}, h_i, p_i\) for the MRxMHD problem which correspond to an MHD equilibrium \cite{10.1063/1.4795739}. More generally, passing to the infinite-interface limit in the MRxMHD variational problem \ref{prob:min1} then yields Euler–Lagrange conditions that coincide with the ideal MHD equilibrium equations \cite{10.1063/1.4795739}. 



\subsection{Background: Existence of MRxMHD Equilibria}



A longstanding conjecture on the existence of MHD equilibria was proposed by Grad \cite{grad_conjecture}. Grad conjectured that, under suitable regularity assumptions, any MHD equilibrium in \(R^3\) must satisfy one of two alternatives: either the pressure is constant, or the magnetic field is invariant under a nontrivial one-parameter group of Euclidean isometries.




An early study on the existence of MHD equilibria with stepped pressure, namely MRxMHD equilibria was conducted by Bruno and Laurence who showed the existence of solutions on domains that are a perturbation away from axisymmetry \cite{bruno_laurence_1996}. In Section \ref{section:Preparatory Results} we show an equivalence between the boundary conditions they provide and the jump condition given in the MRxMHD equilibrium equations when \(\llbracket p \rrbracket \neq 0\).


More recently, Enciso, Luque, and Peralta-Salas demonstrate the existence of solutions to \eqref{prob:MRxMHD} on nondegenerate, analytic, toroidal domains using a Cauchy-Kovalevskaya theorem. They also show that these results are generic \cite{enciso2023mhd}.

The existence of energy minimisers for \(n=1\) on any given domain \(\ST\), which may not be toroidal, has been studied. When the magnetic field is flux-free, existence is guaranteed for any helicity \cite{exist_gerner}. Similarly, Laurence and Avellaneda provide the corresponding Euler-Lagrange equations and an existence proof \cite{laurence1991woltjer}. It was also shown by Parsley that if there exist magnetic fields with a given flux, helicity, and gauge then there exists a corresponding energy minimiser (Theorem 6.1 \cite{parsley2001taylor}). Similarly, a study by Yoshida and Dewar showed existence of solutions to the Beltrami equations with helicity and flux constraints \cite{yoshida2012helical}. 



\subsection{Paper Layout}



One of our principal results is that the MRxMHD equations \eqref{prob:MRxMHD} are necessary and sufficient conditions for a non-zero magnetic field to be a stationary point of the associated variational problem. In particular, we replace the earlier non-degeneracy hypothesis \(\nabla p\neq0\) — which forces boundary components to be unions of tori — by the substantially weaker assumption that each partition element \(\ST_i\) is a compact, oriented, connected volume with smooth boundary. This result permits computation of MRxMHD equilibria via variational methods on a much broader class of geometries than previously considered, including domains with knotted or linked toroidal boundaries (see Figure \ref{fig:trefoil}). 



In Section \ref{section:Problem Construction} we introduce the definitions required to state our principal theorems. Section \ref{section:Main Results} then states and discusses these theorems. Our formulation is expressed in terms of a relative helicity which, under a suitable gauge choice, reduces to the conventional helicity. 



Section \ref{section:Preparatory Results} provides useful results and a reformulation of our variational problem constraining relative helicity. Section \ref{section:Main Proofs} presents proofs for the main results in Section \ref{section:Main Results}. Lastly, Section \ref{section:Characterisation and Local uniqueness} gives a sufficient condition for a critical point of the magnetic energy to also be a local minimum in the \(n=1 \) case. 





\section{Notation}\label{section:Problem Construction}





\subsection{De Rham Cohomology and Singular Homology}

Let \((\ST_1, \dots, \ST_n)\) be a partition of $\ST$ with corresponding interfaces $I_1,\dots,I_m$ (see Section \ref{section:Introduction}). For each $i$ and each integer $k\ge 0$ we write
\[
H^k_{\mathrm{dR}}(\ST_i)\quad\text{and}\quad H_k(\ST_i)
\]
for the $k$th de Rham cohomology group and the $k$th singular homology group of $\ST_i$, respectively, both taken with real coefficients. The homology and cohomology groups of the boundary are denoted analogously for $\partial\ST_i$. The corresponding groups for \(\ST_i\) relative to the boundary \(\partial \ST_i\) are written
\[
H^k_{\mathrm{dR}}(\ST_i,\partial\ST_i)\quad\text{and}\quad H_k(\ST_i,\partial\ST_i) \, .
\]
For the constructions of singular homology and de Rham cohomology we refer to Appendix \ref{sub:construct}. Within the current section we utilise de Rham's theorem and Lefschetz duality, but prove most other results. The outcomes of this section are known and presented here for readers convenience.


For each subdomain $\ST_i$ and each integer $k \in \mathbb{Z}$ there is a bilinear pairing
\begin{align}\label{eq:pairing}
    \big\langle \eq{\omega}, \eq{z} \big\rangle = \int_{z}\omega \, ,
\qquad \eq \omega \in H^k_{\mathrm{dR}}(\ST_i) \, ,\ \ \eq z \in H_k(\ST_i) \, .
\end{align}
The pairing is well-defined by the divergence theorem: it depends only on the cohomology classes \(\eq \omega\) and \(\eq z\), not on the particular representatives \(\omega\) and \(z\). Namely, if $\omega$ is replaced by $\omega+d\nu$ with $\nu$ a smooth $(k-1)$-form, then
\[
\int_z(\omega+d\nu)-\int_z\omega=\int_z d\nu=\int_{\partial z}\nu=0,
\]
since $z$ is a cycle ($\partial z=0$). Likewise, if $z$ is replaced by $z+\partial y$ for some $(k+1)$-chain $y$, then
\[
\int_{z+\partial y}\omega-\int_z\omega=\int_{\partial y}\omega=\int_y d\omega=0,
\]
because $\omega$ is closed (\(d \omega = 0\)). The same Stokes-type argument shows the pairing descends to the relative groups. Namely, there are well defined pairings \(\langle \cdot , \, \cdot \rangle\) by integration for \(H^k_{\mathrm{dR}}(\ST_i, \partial \ST_i) \times H_k(\ST_i)\) and \(H^k_{\mathrm{dR}}(\ST_i, \partial \ST_i) \times H_k(\ST_i, \partial \ST_i)\). 


We write the 3-sphere as \(\mathbb S^3\). Let \(\overline{\ST}_i^c\) denote the closure of \( (\R^3\cup\{\infty\} ) \setminus \ST_i \cong   \mathbb{S}^3 \setminus \ST_i \). Since \(\mathbb S^3\) is compact and \(\overline{\ST}_i^c\) is closed in \(\mathbb S^3\), it follows that \(\overline{\ST}_i^c\) is compact. Moreover, the first homology group \(H_1(\overline{\ST}_i^c)\) is isomorphic to the first homology of the closure of \(\R^3\setminus\ST_i\), see Appendix \ref{sub:complement} for details.
\begin{lemma}\label{Lemma:finite}
    The homologies \(H_k (\ST_i)\) and \(H_k (\overline{\ST}_i^c)\) are finite-dimensional real vector spaces for all \(k,i\).
\end{lemma}

\begin{proof}
    By construction the homologies are vector spaces. To show that these are also finite dimensional, we follow a proof by Bott and Tu (proposition 5.3.1 \cite{bott2013differential} or Appendix \ref{sub:finite}).
\end{proof}

Given Lemma \ref{Lemma:finite}, let the dimension of \(H_1 (\ST_i)\) be denoted \(\ell_i < \infty\) and the dimension of \(H_1 (\overline{\ST}_i^c)\) be denoted \(\kay_i < \infty\). If \(\ell_i, \kay_i > 0 \) then select any bases
\begin{align*}
    \big\{  \eq{s_{i,j} } \big\}_{j=1}^{\ell_i} \ \text{ for } \ H_1(\ST_i) \, , \qquad \text{and} \qquad \big\{  \eq{ t_{i,j} } \big\}_{j=1}^{\kay_i} \ \text{ for } \ H_1 (\overline{\ST}^c_i ) \, .
\end{align*}
If \(\ell_i = 0\) or \(\kay_i = 0\) then the corresponding basis is an empty set. 


Assuming that \(\ell_i, \kay_i > 0 \) and bases for \( H_1(\ST_i) \) and \(H_1 (\overline{\ST}^c_i )\) are given, we may define dual bases for the corresponding de Rham cohomologies because the de Rham theorem tells us that equation \ref{eq:pairing} is a perfect pairing \cite{lee}. Namely, there are unique basis elements \(\eq{s_{i,j}^*} \in H^1_{\mathrm{dR}}(\ST_i)\) and \(\eq{ t_{i,j}^*} \in  H^1_{\mathrm{dR}}(\overline{\ST}^c_i)\) given by 
\begin{align*}
    \begin{gathered}
        \big\langle \eq{ s_{i,k}^* }, \eq{ s_{i,j} } \big\rangle = 
        \delta_{jk}
        \, , \ \text{ for } \ k,j = 1,2, \dots, \ell_i \, ,
        \qquad 
        \big\langle \eq{ t_{i,k}^* }, \eq{ t_{i,j} } \big\rangle = \delta_{jk} \, , \ \text{ for } \ k,j = 1,2, \dots, \kay_i  \, .
    \end{gathered}
\end{align*}
Lefschetz duality combined with de Rham duality provides an isomorphism over compact, oriented 3-manifolds \cite{bott2013differential}
\begin{align*}
    H^1_{\mathrm{dR}}(\ST_i) \cong H_2(\ST_i, \partial \ST_i)
\end{align*}
given by
\begin{align*}
     [\omega] \mapsto [z] \qquad \text{for} \qquad  \int_{\ST_i} \omega \wedge \omega_0
    =   \int_{z}   \omega_0 \, , \quad \text{for all } [ \omega_0 ] \in H^2_{\mathrm{dR}}(\ST_i, \partial \ST_i) \, ,
\end{align*}
where \([\omega] \in H^1_{\mathrm{dR}}(\ST_i)\) and \([z] \in H_2(\ST_i, \partial \ST_i) \). Similarly, for \(H^1_{\mathrm{dR}}(\overline{\ST}^c_i) \cong H_2(\overline{\ST}^c_i, \partial \overline{\ST}^c_i)\). 

With this isomorphism we may utilise the bases for \(H^1_{\mathrm{dR}}(\ST_i)\) and \(H^1_{\mathrm{dR}}(\overline{\ST}^c_i) \) to generate bases for \(H_2(\ST_i, \partial \ST_i)\) and \(H_2(\overline{\ST}^c_i, \partial \overline{\ST}^c_i)\) respectively. These basis elements are denoted \([T_{i,j}]  \in H_2(\ST_i, \partial \ST_i)\), \([S_{i,j}]\in H_2(\overline{\ST}^c_i, \partial \overline{\ST}^c_i)\), and they are given uniquely by the relations,
\begin{align*}
    \int_{\ST_i} s_{i,j}^* \wedge \omega_1  =   \int_{T_{i,j}}  \omega_1 \, , \ \text{ for } \ j = 1,2, \dots, \ell_i \, , 
    \qquad 
    \int_{\overline{\ST}^c_i} t_{i,j}^* \wedge \omega_2  =   \int_{S_{i,j}}  \omega_2\, , \ \text{ for } \ j = 1,2, \dots, \kay_i \, , 
\end{align*}
for all \(\omega_1 \in H^2_{\mathrm{dR}}(\ST_i, \partial \ST_i)\) and \(\omega_2 \in H^2_{\mathrm{dR}}(\overline{\ST}^c_i, \partial \overline{\ST}^c_i)\). 

Utilising de Rham's theorem, the corresponding algebraic duals are computed via the perfect pairing in equation \eqref{eq:pairing}, acting on \( H^2_{\mathrm{dR}}({\ST}_i, \partial {\ST}_i) \times H_2( {\ST}_i, \partial \ST_i)\) or \( H^2_{\mathrm{dR}}(\overline{\ST}^c_i, \partial {\ST}_i) \times H_2( \overline{\ST}^c_i, \partial \ST_i)\). Hence, we define the corresponding basis elements \([T_{i,j}^*] \in H^2_{\mathrm{dR}}( \ST_i, \partial \ST_i) \) and \([S_{i,j}^*]\in H^2_{\mathrm{dR}}( \overline{\ST}^c_i, \partial \ST_i) \) via the relationships
\begin{align*}
    \begin{gathered}
        \big\langle [T_{i,k}^*], [T_{i,j}] \big\rangle = 
        \delta_{jk}
        \, , \ \text{ for } \ k,j = 1,2, \dots, \ell_i \, ,
        \qquad 
        \big\langle [S_{i,k}^*], [S_{i,j}] \big\rangle = \delta_{jk} \, , \ \text{ for } \ k,j = 1,2, \dots, \kay_i  \, .
    \end{gathered}
\end{align*}
We note that combining these algebraic duals with the Lefschetz duality isomorphism, gives the usual non-degenerate Lefschetz pairing,
\begin{align*}
    \int_{\ST_i} s_{i,j}^* \wedge T_{i,k}^* & =   \int_{T_{i,j}} T_{i,k}^* = \delta_{jk}\, , \ \text{ for } \ j,k = 1,2, \dots, \ell_i \, , 
    \\ 
    \int_{\overline{\ST}^c_i} t_{i,j}^* \wedge S_{i,k}^*  &=   \int_{S_{i,j}}  S_{i,k}^*  = \delta_{jk} \, , \ \text{ for } \ j,k = 1,2, \dots, \kay_i \, .
\end{align*}
We are interested in one final collection of variables, again given through an isomorphism. This isomorphism we provide in the following Lemma.
\begin{lemma}\label{Lemma:isomorphism_inclusion}
    Consider the inclusion maps \(\mathcal{J} : \partial \ST_i \hookrightarrow \ST_i\) and \(\mathcal{J}^c : \partial \ST_i \hookrightarrow \overline{\ST}^c_i\). Then \((\mathcal{J}_*, \mathcal{J}^c_*) : H_1(\partial \ST_i) \to H_1(\ST_i) \oplus H_1(\overline{\ST}^c_i)\) is an isomorphism.
\end{lemma}
\begin{proof} 
    By a Mayer–Vietoris sequence for the decomposition \( \mathbb S^3 = \ST_i \cup \overline{\ST}^c_i\), we obtain for \(k>0\) the exact sequence
    \begin{align*}
        H_{k+1} (\mathbb{S}^3) \to H_k (\partial \ST_i) \xrightarrow{ \, (\mathcal{J}_* , - \mathcal{J}^c_* ) \, } H_k (\ST_i) \oplus H_k (\overline{\ST}_i^c) \to H_{k} (\mathbb{S}^3) \, .
    \end{align*}
    We have that \( H_{k} (\mathbb{S}^3) \cong \{0\}\) for \(k=1,2\), hence the map \((\mathcal{J}_* , -\mathcal{J}^c_* )\) is an isomorphism 
    \begin{align}\label{eq:split_boundary}
        H_1 (\partial \ST_i) \cong H_1 (\ST_i) \oplus H_1 (\overline{\ST}_i^c) \, ,
    \end{align}
    as is the map \((\mathcal{J}_* , \mathcal{J}^c_* )\). 
\end{proof}
As \((\mathcal{J}_*, \mathcal{J}^c_*) : H_1(\partial \ST_i) \to H_1(\ST_i) \oplus H_1(\overline{\ST}^c_i)\) is an isomorphism, it has an inverse \((\mathcal{J}_*, \mathcal{J}^c_*)^{-1}\) giving a unique element \([\sigma_{i,j}] \coloneqq (\mathcal{J}_*, \mathcal{J}^c_*)^{-1} ( [s_{i,j}] , 0) \in H_1 (\partial \ST_i)\) for each \(i,j\). Similarly, we define a unique element \([\tau_{i,j}] \coloneqq (\mathcal{J}_*, \mathcal{J}^c_*)^{-1} ( 0, [t_{i,j}] ) \in H_1 (\partial \ST_i)\) for each \(i,j\).


\begin{lemma}\label{Lemma:form_basis}
Given non-empty bases \( \{[s_{i,j}]\}_{j=1}^{\ell_i}\), \(\{[t_{i,j}]\}_{j=1}^{\kay_i}\), of \(H_1(\ST_i)\) and \(H_1(\overline{\ST}^c_i)\) respectively, then
\[
\big\{[\sigma_{i,1}], \dots , [\sigma_{i,\ell_i}],[\tau_{i,1}], \dots, [\tau_{i,\kay_i}]\big\}
\]
is a basis of \(H_1(\partial\ST_i)\).
\end{lemma}
\begin{proof}
The collection \(\{([s_{i,1}],0), \dots, (0,[t_{i,1}]), \dots\}\) is a basis of the direct sum \(H_1(\ST_i)\oplus H_1(\overline{\ST}_i^c)\). Applying the isomorphism \((\mathcal J_*,\mathcal J^c_*)^{-1}\) carries that basis to the set \(\{[\sigma_{i,1}]\mathbin{,}\dots \mathbin{,} [\tau_{i,1}] \mathbin{,} \dots\}\) in \(H_1(\partial\ST_i)\). An isomorphism sends bases to bases, therefore giving a basis of \(H_1(\partial\ST_i)\).
\end{proof}



We can consider maps induced on de Rham cohomologies via the inclusions \(\mathcal J\) and \(\mathcal J^c\). In this case there is a Mayer-Vietoris sequence for de Rham classes
\begin{align*}
    \{0\} \cong H^{1}_{\mathrm{dR}} ( \mathbb{S}^3 ) \to H^{1}_{\mathrm{dR}} ( \ST_i ) \oplus H^{1}_{\mathrm{dR}} ( \overline{\ST}^c_i ) \xrightarrow{\mathcal{J}^* -(\mathcal{J}^c)^*} H^{1}_{\mathrm{dR}} ( \partial \ST_i ) \to H^{2}_{\mathrm{dR}} (  \mathbb{S}^3 ) \cong \{0\} \, , 
\end{align*} 
so the map \(\mathcal{J}^*- (\mathcal{J}^c)^*\) is an isomorphism, acting as
\begin{align*}
    (\mathcal{J}^*- (\mathcal{J}^c)^*) ([s],[t]) = \mathcal{J}^*[s]- (\mathcal{J}^c)^*[t] \qquad \text{for} \qquad ([s], [t]) \in H^{1}_{\mathrm{dR}} ( \ST_i ) \oplus H^{1}_{\mathrm{dR}} ( \overline{\ST}^c_i ) \, .
\end{align*}
Hence \(\mathcal{J}^* + (\mathcal{J}^c)^* \) is an isomorphism. We define the unique elements \([\sigma_{i,j}^*] \coloneqq (\mathcal{J}^* + (\mathcal{J}^c)^*) ( [s_{i,j}^*] , 0) \in H^1_{\mathrm{dR}} (\partial \ST_i)\) and \([\tau_{i,j}^*] \coloneqq (\mathcal{J}^* + (\mathcal{J}^c)^*)( 0, [t_{i,j}^*] ) \in H^1_{\mathrm{dR}} (\partial \ST_i)\). Note that \([\sigma_{i,j}^*], [\tau_{i,j}^*]\) are clearly not defined as the duals of \([\sigma_{i,j}], [\tau_{i,j}]\) respectively. But these are, in fact, the unique duals in \(H_1(\partial \ST_i)\), as
\begin{align*}
\big\langle [\sigma^*_{i,k}],[\sigma_{i,j}] \big\rangle
=\int_{\sigma_{i,j}} \sigma^*_{i,k}
=\int_{\sigma_{i,j}} \mathcal J^* (s^*_{i,k}) +  (\mathcal{J}^c)^* (0)
= \int_{\sigma_{i,j}} \mathcal J^* (s^*_{i,k})
= \int_{\mathcal J^* \sigma_{i,j}} s^*_{i,k}
= \int_{s_{i,j}} s^*_{i,k}
= \delta_{jk} \, , 
\end{align*}
similarly \(\delta_{jk} = \langle  [\tau_{i,k}^*]  , [\tau_{i,j}] \rangle\). Additionally, the pairings between the remaining elements,
\begin{align*}
\big\langle [\tau^*_{i,k}],[\sigma_{i,j}] \big\rangle
= \int_{\sigma_{i,j}} \tau^*_{i,k}
= \int_{\sigma_{i,j}}   \mathcal J^* (0) +  (\mathcal{J}^c)^* (t^*_{i,k})
= \int_{\mathcal{J}^c_* \sigma_{i,j}}   t^*_{i,k}
= 0\, , 
\end{align*}
similarly \(\langle [\sigma^*_{i,k}],[\tau_{i,j}]\rangle=0\).



\begin{theorem}\label{thr:Alexander}
    For \(\ell_i > 0\) consider any basis \( \{ [s_{i,1}], \dots , [s_{i,\ell_i}]\} \subset H_1 ( \ST_i) \). Then there exist a choice of basis \(\{[t_{i,1}], \dots , [t_{i,\ell_i}]\}\) for \(H_1 (\overline{\ST}^c_i )\) so that the following additional conditions hold:
    \begin{enumerate}
        \item For all \( j, k \in \{ 1, \dots, \ell_i \} \) and \(i = 1,2, \dots, n\): 
        \begin{align*}
        \begin{gathered}
             \int_{\partial \ST_i} \tau_{i,j}^* \wedge \sigma_{i,k}^* = \delta_{jk}  \,, \qquad
            \int_{\partial \ST_i} \sigma_{i,j}^* \wedge \sigma_{i,k}^* = 
             0 \,,    \qquad 
             \int_{\partial \ST_i} \tau_{i,j}^* \wedge \tau_{i,k}^*
            = 0 \, .
        \end{gathered}
        \end{align*}   
        \item The basis elements \( [S_{i,j}] \in H_2( \overline{\ST}^c_i, \partial \overline{\ST}^c_i) \) and \( [T_{i,j}] \in H_2( \ST_i, \partial \ST_i) \) satisfy \( \partial [S_{i,j}] =  [\sigma_{i,j}] \) and \( \partial [T_{i,j}] =  [\tau_{i,j}] \) for \( j = 1, \dots, \ell_i \). 
    \end{enumerate}
\end{theorem}


\begin{proof}
    An exposition of the original proof by Cantarella and Parsley is adapted into differential forms and presented in Appendix \ref{section:Alexander Basis} (theorems B.2, B.3 \cite{CANTARELLA20101127}). 
\end{proof}




For each domain \(\ST_i\) with \(\ell_i > 0\), any choice of basis for \(H_1(\ST_i)\), together with a basis for \(H_1(\overline{\ST}_i^{\,c})\) as in Theorem~\ref{thr:Alexander}, induces a basis for \(H_1(\partial \ST_i)\) via Lemma~\ref{Lemma:form_basis}. We refer to such a basis of \(H_1(\partial \ST_i)\) as an \emph{Alexander basis}. If \(\ell_i = 0\), then \(H_1(\partial \ST_i) \cong \{0\}\) (see appendix~\ref{section:Alexander Basis}), and we define the Alexander basis to be empty.



The constructions in this section are encoded by the following commutative diagram, which records the relevant homology classes associated to the reference basis element \([s_{i,j}] \in H_1(\ST_i)\). Taking duals with respect to the corresponding basis yields the associated de Rham cohomology basis elements:
\[
    \begin{tikzcd}[column sep=5em]
        {[T_{i,j}]\in H_2(\ST_i,\partial\ST_i)} \arrow[swap]{d}{\partial}
        & {[s_{i,j}]\in H_1(\ST_i)} \arrow[swap]{l}{\text{LD}^*}  
        & {[\sigma_{i,j}]\in H_1(\partial\ST_i)} \arrow[swap]{l}{\mathcal{J}_*} \\
        {[\tau_{i,j}]\in H_1(\partial\ST_i)} \arrow{r}{\mathcal{J}^c_*}
        & {[t_{i,j}]\in H_1(\overline{\ST}_i^c)} \arrow{r}{\text{LD}^*} 
        & {[S_{i,j}]\in H_2(\overline{\ST}^c_i,\partial\overline{\ST}^c_i)} \arrow[swap]{u}{\partial}
    \end{tikzcd}
\]
Here \(\mathrm{LD}^*\) denotes Lefschetz duality followed by passage to the dual basis. The induced map \([s_{i,j}] \mapsto [t_{i,j}]\), together with this dualisation, realises \emph{Alexander duality}. Moreover, the composition of all maps around the diagram is the identity.


\begin{figure} 
  \centering
  \begin{overpic}[width=0.5\linewidth]{Chapter_Problem_Construction/topology3.png}
  \end{overpic}
  \caption{A volume with 3 toroidal boundaries.}\label{fig:topology} 
\end{figure}


We define the \emph{genus} of \(\partial \ST_i\) to be the dimension \(\ell_i\) of \(H_1(\ST_i)\) (equivalently, of \(H_1(\overline{\ST}^{\,c}_i)\)). This equals the total genus of the boundary \(\partial \ST_i\). In MHD, as discussed in Subsection \ref{sub:back}, we typically work with toroidal level sets and boundaries. In this setting, the genus is obtained by counting the number of disconnected toroidal boundary components, since each torus has genus \(1\). For example, the domain shown in Figure \ref{fig:topology}, whose boundary consists of three tori, has total genus \(3\) \cite{cantarella2002vector}.



\subsection{Relative Helicity}\label{subsection:Relative Helicity}




As noted in the introduction \ref{sub:setup}, helicity \(H_{\Omega^1(\ST_i) }( \beta_i , \ST_i ) \) is not, in general, a well-defined function of \( \beta _i \), owing to its dependence on the choice of gauge for the associated potential. One way to obtain a well-defined quantity is to restrict the admissible class of potentials \(G\) so that the helicity \(H_{G }( \beta_i , \ST_i ) \) is uniquely determined for any given \( \beta _i \), as is done in problem \ref{prob:min1}. 

An alternative, which we adopt here, is to introduce a \emph{relative helicity}. This quantity is constructed to be independent of the gauge choice for a given \( \beta _i \), while coinciding with the standard helicity when evaluated in a fixed reference gauge. In this way, helicity retains a clear physical interpretation without the need to impose additional constraints on the potential.



Given an Alexander basis we define the \textit{relative helicity} of \(\beta_i\) over \(\ST_i\) for \(\ell_i > 0\) and any primitive \(a_i \in \Omega^1(\ST_i)\) of \(b_i\):
\begin{align*}
    \mathscr{H}_{a_i}(\beta_i, \ST_i) \coloneqq \int_{\ST_i} a_i \wedge da_i - \sum_{j=1}^{\ell_i} \int_{\sigma_{i,j}} a_i \int_{\tau_{i,j}} a_i \, ,
\end{align*}
and \(\mathscr{H}_{a_i}(\beta_i, \ST_i) \coloneqq H_G (\beta_i, \ST_i)\) for \(\ell_i = 0\), and \(G=\Omega^1(\ST_i)\). Refer to Lemma \ref{Lemma:well-pose} to see that the relative helicity is independent of the chosen primitive for a given \(\beta _i\). 



We refer to the following set as an Amperian gauge:
\begin{align}\label{eq:amperian}
    \mathscr{A}(\ST_i) = \bigg\{ a_i \in  \Omega^1 (\ST_i) \, : \; \int_{\sigma_{i,j}} a_i = 0 \, , \; \forall j \bigg\} \, .
\end{align}
An Amperian primitive for \(\beta _i\) is any \(a_i \in \mathscr{A}(\ST_i) \) such that \(d a_i = \beta _i\). See Lemma \ref{Lemma:gauge_exist} for a proof that any \(\beta _i \in \C\) has a primitive in \(\mathscr{A}(\ST_i)\). We compute the following change in helicity by applying the divergence theorem:
\begin{align*}
    H_{G}(\beta_i, \ST_i) - H_{G'}(\beta_i, \ST_i) = \int_{\partial \ST_i} s \wedge a_i \, ,
\end{align*}
for primitives \(a_i \in G, a_i + s \in G'\) of \(\beta _i\), where \(s\) is closed, for some \(G,G' \subseteq \Omega^1(\ST_i)\) that ensure helicity is independent of the chosen potential. This expression can be interpreted as a cup product on the boundary \( \partial \ST_i \), and its value depends on the interaction between cohomology classes for primitives. As the Amperian gauge fixes the primitive up to a de Rham cohomology class, the helicity is independent of the Amperian primitive choice (Lemma \ref{Lemma:well-pose}). Note that the definition of an Amperian gauge is similar to the definition of Amperian knots given by Cantarella \cite{plasma_flow}. 

Selecting an Amperian gauge reduces the relative helicity to helicity,
\begin{align*}
    H_{\mathscr{A}(\ST_i)}(\beta_i,\ST_i) = \mathscr{H}_{a_i + s}(\beta_i,\ST_i) \quad \text{for} \quad a_i \in \mathscr{A}(\ST_i)  \, ,    
\end{align*}
for all closed \(s \in \Omega^1(\ST)\). Recall that all primitives of \(\beta \) are of the form \(a_i + s\). Intuitively, the relative helicity encodes the helicity in an Amperian reference gauge. Therefore, any result using a relative helicity also holds for helicity under a restriction of the primitives to an Amperian gauge. 






\subsection{Hodge-Morrey-Friedrichs Decomposition} \label{sub:Hodge}

The space $\C$ decomposes via a Hodge-Morrey-Friedrichs (HMF) decomposition \cite{schwarz-1995, Gauge_freedom}. Firstly, we define some notation for function spaces. We reserve \(\mathcal{H}^k(\ST_i)\) for harmonic \(k\)-forms on \(\ST_i\), and a subscript \(D\) or \(N\) indicates the Dirichlet of Neumann boundary conditions, as in \( \Omega^k (\ST_i)\). More explicitly, we have
\begin{align*}
    \mathcal{H}^k_D (\ST_i)  = \{   \lambda \in  \Omega^k (\ST_i) \, : \; d \lambda =  \delta \lambda =  t \lambda = 0    \} \, , \qquad \mathcal{H}^k_N (\ST_i)  = \{   \lambda \in  \Omega^k (\ST_i) \, : \; d \lambda =  \delta \lambda =  n \lambda = 0    \} \, . 
\end{align*}
Consider \(\beta _i \in \C\), then we have a HMF decomposition (see Appendix \ref{append:Hodge}):
\begin{align*}
    \star \beta _i = \delta (\star \eta_i) + d \phi_i + dz_i + \star \lambda_i \, , 
\end{align*}
where \(\star \eta_i \in  \Omega_N^2 (\ST_i) \), \(\lambda_i \in \mathcal{H}_D^2(\ST_i)\), and \(\phi_i,z_i \in  \Omega^0 (\ST_i)\) for \(\phi_i|_{\partial \ST_i} = 0\), and \(z_i\) is a harmonic function \(\Delta z_i = 0\) where \(\Delta \) is the Laplacian. Via an application of a Hodge star isomorphism:
\begin{align*}
    \beta _i = d \eta_i + \star d\phi_i + \star dz_i + \lambda_i \, .
\end{align*}
Using \(d \beta _i = d \star d \phi_i = 0\) and \(\phi_i|_{\partial \ST_i} = 0\) with the divergence theorem:
\begin{align*}
    \| d \phi_i \|_{L^2}^2 = \int_{\ST_i}d \phi_i \wedge \star d \phi_i = -  \int_{\ST_i} \phi_i \wedge d \star d \phi_i + \int_{\partial \ST_i} \phi_i \wedge \star d \phi_i = 0 \, .
\end{align*}
With smoothness of \(\phi_i\), we have that \(d \phi_i = 0\). Similarly for the \(dz_i\) component, utilising that \(t \beta _i = t ( \star dz_i ) =0 \):
\begin{align*}
    \| d z_i \|_{L^2}^2 = \int_{\ST_i}d z_i \wedge \star d z_i =  -  \int_{\ST_i} z_i \wedge d \star d z_i + \int_{\partial \ST_i} t(z_i) \wedge t( \star d z_i) = 0 \, . 
\end{align*}
By smoothness of \(z_i\) we have \(dz_i = 0\). Therefore we are left with a decomposition \(\beta _i = d \eta_i + \lambda_i\). By the Hodge isomorphism and Poincaré–Lefschetz duality
\begin{align*}
    \h \cong H^2_{\mathrm{dR}} ( \ST_i, \partial \ST_i) \cong H^1_{\mathrm{dR}} ( \ST_i)  \, ,
\end{align*}
for coefficients in \(\R\). So \(\h\) is a \(\ell_i\)-dimensional vector space. The set \(\{ [T_{i,j}^*] \}_{j=1}^{\ell_i}\) (\(\ell_i > 0\)) form a basis for \( H^2_{\mathrm{dR}} (\ST_i, \partial \ST_i )\). The Hodge theorem says that there is a unique harmonic representative for each class \([T_{i,j}^*]\). So, without loss of generality, we select \(T_{i,j}^* \in \h \). As \(T_{i,j}^*\) are a basis for \(\h\), for constants \(\psi_{i,j}\) we have
\begin{align}
    \beta _i = d \eta_i + \lambda_i = d \eta_i + \sum_{j=1}^{\ell_i} \psi_{i,j} T_{i,j}^* \, . 
\end{align}
In the case of \(\ell_i = 0\):
\begin{align*}
    \beta _i = d \eta_i  \, . 
\end{align*}
Examples of these decompositions in the case of solid and hollow tori are given by Pfefferlé, Noakes, and Perrella \cite{Gauge_freedom}.




\subsection{Flux} \label{sub:flux}

Let \(\beta _i\in\C\) be a closed Dirichlet \(2\)-form on \(\ST_i\), and let \([T]\in H_2(\ST_i,\partial\ST_i)\) be a relative homology class. The \emph{flux} of \(\beta _i\) through \(T\) is by definition the pairing \(\langle[\beta _i],[T]\rangle\). Applying the divergence theorem to this pairing gives
\begin{align*}
\big\langle [\beta _i],[T]\big\rangle
&= \int_T \beta _i
= \int_{\partial T}\mathcal{J}^*\eta_i + \sum_{j=1}^{\ell_i}\psi_{i,j}\int_T T_{i,j}^*
= \sum_{j=1}^{\ell_i}\psi_{i,j}\int_T T_{i,j}^* =\sum_{j=1}^{\ell_i}\psi_{i,j}\big\langle [T_{i,j}^*],[T]\big\rangle \, ,
\end{align*}
since \(\eta_i\) vanishes on \(\partial\ST_i\). In particular, for the relative cycles \(T_{i,1},\dots,T_{i,\ell_i}\) we obtain
\[
\big\langle [\beta _i],[T_{i,j}]\big\rangle=\psi_{i,j}\qquad(1\le j\le\ell_i) \, ,
\]
so that specifying the values \(\psi_{i,1},\dots,\psi_{i,\ell_i}\) is equivalent to fixing the harmonic component \(h_i\) and therefore determines the flux of \(\beta _i\) through every class \([T]\). This is the set of constraints we impose when formulating the variational problem with fixed flux. The flux constraints are invariant under the push-forward action of any \(f\in\I\), since
\[
\big\langle[\beta _i],[T_{i,k}]\big\rangle
=\big\langle[(f^{-1})^* \beta _i],[f(T_{i,k})]\big\rangle \, .
\]
When \(\ell_i=0\) there is no harmonic component and every closed Dirichlet \(2\)-form is exact: \(\beta_i=d\eta_i\) with \(\eta_i\in\Omega^1_D(\ST_i)\). The divergence theorem then shows that the flux through any relative class vanishes,
\[
\langle[\beta _i],[T]\rangle=\int_T d\eta_i=\int_{\partial T}\mathcal{J}^*\eta_i=0 \, .
\]
Therefore \(\beta_i\) is an element of the flux-free subspace of closed Dirichlet \(2\)-forms when \(\ell_i=0\).






\subsection{Connected \texorpdfstring{\(\ST_i\)}{Domains}}

We assume that each partition element \(\ST_i\) is connected. This hypothesis removes a benign but inconvenient source of non-uniqueness, which we now explain.

If \(\ST_i\) is a connected, compact domain with smooth boundary, then a Beltrami field on \(\ST_i\) is determined uniquely by the Beltrami parameter \(\mu_i\) together with the flux constraints \(\psi_{i,1},\dots,\psi_{i,\ell_i}\) \cite{yoshida2012helical}. Replacing the prescription of \(\mu_i\) by a constraint on the relative helicity still yields existence of a solution to the Beltrami equations.

By contrast, if \(\ST_i\) has several connected components then a helicity condition imposed only on the whole of \(\ST_i\) fixes merely the sum of the helicities of the components. In that situation the individual component helicities may be varied continuously while keeping their total fixed, which produces a non-trivial continuous family of admissible Beltrami fields. Requiring each \(\ST_i\) to be connected therefore removes this extra degree of freedom.





\section{Main Results}\label{section:Main Results}



Consider the following objects as given:
\begin{enumerate}
  \item A finite collection of compact, oriented, connected, pairwise disjoint subdomains \(\ST_i\) \((i=1,\dots,n)\) that partition \(\ST\). 
  \item An Alexander basis of \(H_1(\partial\ST_i)\) for each \(i\).
  \item For each \(i\), if \(\ell_i=0\), then \(p_i\) and \(h_i\) are given. If \(\ell_i>0\), then \(p_i\), \(h_i\) and \(\psi_{i,1},\dots,\psi_{i,\ell_i}\) are given. 
\end{enumerate}
Consider the variational problem:
\begin{align} \label{eq:ext}
    \begin{gathered} 
    \underset{\substack{f \in \I \\ \beta _i \in \mathcal{C} ( \ST_i)} }{ \text{c. p.} } \left( \sum_{i=1}^n  \frac{1}{2} \big\|\beta _i\big\|_{L^2}^2 -  p_i \myabs{ \ST_i } \right) \, , \\
    \mathscr{H}_{a_i}( \beta_i,  \ST_i) = h_i\, , \qquad
    \int_{T_{i,j} } \beta _i = \psi_{i,j}\, ,
    \end{gathered}
\end{align}
where \(d a_i = \beta_i\) for all \(i\). The constraints hold for all \(i = 1, \dots, n\) and relevant \(j\). The notation “c. p.'' stands for critical points. If the pressure, helicity, and flux constants for any \(i\) are identically zero (\(p_i = h_i = \psi_{i,1} = \dots = \psi_{i,\ell_i}=0\)) then \(\beta _i=0\) is the unique global minimiser. 

We remind the reader that we are fixing \(\ST_i\) and allowing the metric to change through \(f \in \I\). This effectively determines the locations of interfaces within \(f(\ST)\). 

The following theorem details a connection between solutions to the variational problem \eqref{eq:ext} and solutions to the MRxMHD equations \eqref{prob:MRxMHD}. 

\begin{theorem}\label{thr:iff_simple}
    A non-zero element \(\beta_i \in \C\) is a solution to the variational problem \eqref{eq:ext} if and only if it satisfies the MRxMHD equations \eqref{prob:MRxMHD} with the corresponding helicity and flux constraints.
\end{theorem}


\begin{corollary}%\label{cor:iff_simple}
    If \((h_i, \psi_{i,1}, \dots, \psi_{i, \ell_i}) \neq 0\) for \(i = 1, \dots, n\) then the MRxMHD equations \eqref{prob:MRxMHD} with the corresponding helicity and flux constraints are necessary and sufficient conditions for all solutions to the variational problem \eqref{eq:ext}.
\end{corollary}






\begin{remark}\label{remark:MHD}
  The variational problem \eqref{eq:ext} is formulated in terms of critical points, but for physical applications one is usually interested in local \emph{minimisers} (which correspond to stable equilibria). Woltjer's 1958 analysis introduced a closely related variational principle and argued in favour of minimising solutions. In the proof of Theorem \ref{thr:iff_simple} we show that, for the case \(n=1\) and upon adopting an Amperian gauge, Woltjer's theorem is recovered.
\end{remark}

Theorem \ref{thr:iff_simple} implies that the MRxMHD equations \eqref{prob:MRxMHD} are necessary for a minimiser to problem \eqref{eq:ext} with nested toroidal surfaces. This result is known, as discussed in subsection \ref{sub:back}. 

On a technical note, we define a critical point of a map \( \I \times  \mathcal{C} ( \ST_i) \to \R\) as a point in the domain where the Fréchet derivative in a \(C^\infty\) topology for \(\mathcal{C} ( \ST_i)\) and a Gâteau derivative with respect to the \(\I\) coordinates are zero. We require a Fr\'echet derivative combine with a choice of topology for \(\mathcal{C} ( \ST_i)\) in order to determine a necessary and sufficient condition when applying the Lagrange multiplier theorem. 


  \begin{remark}
    A result establishing the existence of solutions to the Beltrami equation was proven by Yoshida and Dewar. These authors use the Fredholm alternative and a technique they refer to as \textit{helicity matching} (section 2.3 \cite{yoshida2012helical}). Therefore, there is always a solution to problem \eqref{eq:ext} when \(n=1\). In Appendix \ref{append:exist}, we prove an existence result that is weaker than the known result above, using the direct method from calculus of variations.
 \end{remark}
    
    


\section{Preparatory Results}\label{section:Preparatory Results}


\begin{lemma}\label{Lemma:compute}
    For any closed 1-forms \(s, \tilde{s} \in \Omega^{1} (\partial \ST_i)\) and an Alexander basis \(\{[\sigma_{i,j}],[\tau_{i,j}]\}_{j=1}^{\ell_i}\) of \(H_1(\partial \ST_i)\) we have a decomposition:
    \begin{align}\label{eq:compute_2}
    \int_{\partial \ST_i} \tilde{s} \wedge s = \sum_{j=1}^{\ell_i}  \left(\int_{\sigma_{i,j}} \! s \int_{\tau_{i,j}}  \! \tilde{s} -\int_{\sigma_{i,j}}  \! \tilde{s}    \int_{\tau_{i,j}}  \! s   \right) 
     \, .
\end{align}
\end{lemma}


\begin{proof}

Any 1-form \( s \in \Omega^{1} (\partial \ST_i)\) admits a Hodge-Morrey decomposition into an exact form, co-exact form, and a harmonic 1-form representing a cohomology class (\cite{schwarz-1995}, see Appendix \ref{append:Hodge} for this decomposition):
\begin{align*}
    s = \delta \eta + d \phi + \gamma  \, , %& &  \tilde{s} = d \tilde{\phi} + \tilde{\gamma} \, ,
\end{align*}
where \( \eta \in  \Omega^2 (\partial \ST_i) \), \(\phi \in  \Omega^0 (\partial  \ST_i)\), and a harmonic \(\gamma \in  \mathcal{H}^1(\partial \ST_i)\). Utilising the closure of \(s\) (\(0 = ds = d \delta  \eta\)) and the divergence theorem,
\begin{align*}
    \| \delta  \eta \|^2_{L^2} = \int_{\partial \ST_i} \delta  \eta \wedge \star \delta  \eta
    = -\int_{\partial \ST_i} 
    \delta  \eta \wedge d \star \eta
    = - \int_{\partial \ST_i} 
    d\delta  \eta \wedge  \star \eta
+
\int_{\partial^2 \ST_i} \delta  \eta \wedge  \star \eta = 0 \, . 
\end{align*}
Note that the Hodge star operators on \(\partial \ST_i\) acting on 1-forms satisfy \(\star \star = -1\). Therefore, the co-exact component in the Hodge-Morrey decomposition of a closed 1-form disappears. Taking the Hodge decomposition for two closed \(1\)-forms \(s, \tilde{s} \in \Omega^{1} (\partial \ST_i)\):
\begin{align*}
    s = d \phi + \gamma  \, , & &  \tilde{s} = d \tilde{\phi} + \tilde{\gamma} \, ,
\end{align*}
where \(\tilde{\phi} \in \Omega^0 (\ST_i)\) and \(\tilde{\gamma} \in  \mathcal{H}^1 (\partial \ST_i)\). 

Moving all terms in equation \eqref{eq:compute_2} to one side, we want to see if the result is zero. We begin this check by replacing the closed 1-forms with their Hodge-Morrey decomposition:
\begin{align*}
\begin{gathered}
    \int_{\partial \ST_i} \tilde{s} \wedge s - \sum_{j=1}^{\ell_i}  \left( \int_{\sigma_{i,j}} s \int_{\tau_{i,j}} \tilde{s} - \int_{\sigma_{i,j}} \tilde{s}    \int_{\tau_{i,j}} s \right) 
    \\ = \\   \int_{\partial \ST_i} (d \tilde{\phi} + \tilde{\gamma}) \wedge (d \phi + \gamma) -  \sum_{j=1}^{\ell_i}  \left( \int_{\sigma_{i,j}} (d \phi + \gamma) \int_{\tau_{i,j}} (d \tilde{\phi} + \tilde{\gamma}) - \int_{\sigma_{i,j}} (d \tilde{\phi} + \tilde{\gamma})    \int_{\tau_{i,j}} (d \phi + \gamma) \right)  \, . 
\end{gathered}
\end{align*}
Knowing that \(\partial \tau_{i,j} = \partial \sigma_{i,j} = 0\) we have, by the divergence theorem:
\begin{align*}
    \int_{\tau_{i,j}} d \phi  = \int_{\partial \tau_{i,j}} \phi = 0 \, ,
\end{align*}
similarly for the integrals over \(\sigma_{i,j}\). Additionally:
\begin{align*}
    \int_{\partial \ST_i} d \tilde{\phi}  \wedge s
    =
    - \int_{\partial \ST_i}  \tilde{\phi}  \wedge ds
    +
    \int_{\partial^2 \ST_i}  \tilde{\phi}  \wedge s = 0 \, .
\end{align*}
Similarly for the term involving \(\tilde s \wedge d \phi  \). This leaves us with the relation:
\begin{align}\label{eq:compute}
    \int_{\partial \ST_i} \tilde{s} \wedge s - \sum_{j=1}^{\ell_i}  \left( \int_{\sigma_{i,j}} \! s \int_{\tau_{i,j}}  \!  \tilde{s} - \int_{\sigma_{i,j}} \!  \tilde{s}    \int_{\tau_{i,j}} \!  s \right) 
    &=   \int_{\partial \ST_i} \tilde{\gamma} \wedge \gamma -  \sum_{j=1}^{\ell_i}  \left( \int_{\sigma_{i,j}} \!  \gamma \int_{\tau_{i,j}}  \! \tilde{\gamma} - \int_{\sigma_{i,j}}  \! \tilde{\gamma}    \int_{\tau_{i,j}}  \! \gamma \right)  \, .
\end{align}
Replacing \(\gamma\) and \(\tilde{\gamma}\) with their classes in \(H^1_{dR}(\partial \ST_i)\) given by \([\gamma]\) and \([\tilde{\gamma}]\) respectively, these integrals remain well-defined. Decomposing \([\gamma] \) and \([\tilde{\gamma}]\) in the dual to the Alexander basis we have that 
\begin{align*}
    [\gamma] = \sum_{j=1}^{\ell_i} c_{j}^1  [\sigma_{i,j}^*]+ c_{j}^2 [\tau_{i,j}^*] \, , & &
    [\tilde{\gamma}] = \sum_{j=1}^{\ell_i} \tilde{c}_{j}^1 [\sigma_{i,j}^*]+ \tilde{c}_{j}^2 [\tau_{i,j}^*] \, ,
\end{align*}
for constants \( c_{j}^1,  c_{j}^2, \tilde{c}_{j}^1\) and \(\tilde{c}_{j}^2\). Computing our wedge product between the gamma's:
\begin{align*}
    \int_{\partial \ST_i} \tilde{\gamma} \wedge \gamma  & = \scalebox{1.2}[1.4]{$\displaystyle\int$}_{ \! \! \! \! \! \partial \ST_i} \!  \Bigg(  \sum_{j=1}^{\ell_i} \tilde{c}_{j}^1 \sigma_{i,j}^*+ \tilde{c}_{j}^2 \tau_{i,j}^* \Bigg) \wedge \Bigg(   \sum_{j=1}^{\ell_i} c_{j}^1  \sigma_{i,j}^*+ c_{j}^2 \tau_{i,j}^*   \Bigg) \, .
\end{align*}
Recall from Theorem \ref{thr:Alexander} that: 
\begin{align*}
    \begin{gathered}
         \int_{\partial \ST_i} \tau_{i,j}^* \wedge \sigma_{i,k}^* = \delta_{jk}  \,, \qquad
        \int_{\partial \ST_i} \sigma_{i,j}^* \wedge \sigma_{i,k}^* = 
         0 \,,    \qquad 
         \int_{\partial \ST_i} \tau_{i,j}^* \wedge \tau_{i,k}^*
        = 0 \, ,
    \end{gathered}
\end{align*}   
so
\begin{align*}
\int_{\partial \ST_i} \tilde{\gamma} \wedge \gamma & 
  =  \sum_{j=1}^{\ell_i} 
  \tilde{c}_{j}^2  c_{j}^1  - \tilde{c}_{j}^1  c_{j}^2 
  \, . 
\end{align*}
And:
\begin{align*}
     \sum_{j=1}^{\ell_i}  \left( \int_{\sigma_{i,j}}  \!  \gamma \int_{\tau_{i,j}}  \!  \tilde{\gamma} - \int_{\sigma_{i,j}}  \!  \tilde{\gamma}    \int_{\tau_{i,j}}  \!  \gamma \right) 
      %&= 
     = \sum_{j=1}^{\ell_i}  \left( c^1_{j}  \tilde{c}^2_{j} - \tilde{c}^1_{j}  c^2_{i,j}  \right) 
     & = \int_{\partial \ST_i} \tilde{\gamma} \wedge \gamma \, . 
\end{align*}
Substituting this relation into equation \eqref{eq:compute} we get that:
\begin{align*}
    \int_{\partial \ST_i} \tilde{s} \wedge s - \sum_{j=1}^{\ell_i}  \left( \int_{\sigma_{i,j}}  \!  s \int_{\tau_{i,j}}  \!  \tilde{s} - \int_{\sigma_{i,j}}  \!  \tilde{s}    \int_{\tau_{i,j}}  \!  s \right) 
    &=   0 \, .
\end{align*}
\end{proof}





\subsection{Well-Posed}\label{sub:Well-Posed}


\begin{lemma}\label{Lemma:well-pose}
    For a magnetic field represented by a closed, Dirichlet 2-form \(\beta _i \in \C\) with potential \(a_i \in  \Omega^1(\ST_i)\), the relative helicity \(\mathscr{H}_{a_i}( \beta_i, \ST_i)\) is independent of the choice of primitive \(a_i\). Additionally, if \(f \in \I\) then \(\mathscr{H}_{a_i}( \beta_i, \ST_i)= \mathscr{H}_{(f^{-1})^* a_i}( (f^{-1})^*  \beta_i, f(\ST_i ) )\). 
\end{lemma}

\begin{proof} 
Consider \(a_i, \tilde{a}_i\in  \Omega^1 (\ST_i)\) such that \(da_i = d \tilde{a}_i = \beta_i \in \C\). Then the difference in relative helicities applied to each primitive is:
\begin{align*}
    \mathscr{H}_{a_i}(\beta_i, \ST_i) - \mathscr{H}_{\tilde{a}_i}(\beta_i, \ST_i) = \int_{\partial \ST_i} \mathcal{J}^* \tilde{a}_i \wedge \mathcal{J}^* a_i
    - \left( \sum_{j=1}^{\ell_i} \int_{\sigma_{i,j}}  \!  \mathcal{J}^* a_i\int_{\tau_{i,j}}   \!  \mathcal{J}^* a_i
    - \int_{\sigma_{i,j}}  \!  \mathcal{J}^* \tilde{a}_i \int_{\tau_{i,j}}  \!  \mathcal{J}^* \tilde{a}_i \right)
    \, .
\end{align*}
The forms \(\mathcal{J}^* (\tilde{a}_i)\) and \(\mathcal{J}^* (a_i)\) are closed 1-forms in \(\Omega^1(\partial \ST_i)\), so we may apply equation \eqref{eq:compute_2}:
\begin{align*}
    \mathscr{H}_{a_i}(\beta_i, \ST_i) - \mathscr{H}_{\tilde{a}_i}(\beta_i, \ST_i) = -
    \sum_{j=1}^{\ell_i}  \Bigg( & \int_{\sigma_{i,j}}  \! \mathcal{J}^* \tilde{a}_i    \int_{\tau_{i,j}}  \! \mathcal{J}^*a_i
    -
    \int_{\sigma_{i,j}}  \!  \mathcal{J}^*a_i \int_{\tau_{i,j}}  \! \mathcal{J}^*\tilde{a}_i 
    \\
    & \! \! 
    + \int_{\sigma_{i,j}}  \! \mathcal{J}^*a_i \int_{\tau_{i,j}}  \! \mathcal{J}^*a_i
    - \int_{\sigma_{i,j}}  \! \mathcal{J}^*\tilde{a}_i \int_{\tau_{i,j}}  \! \mathcal{J}^*\tilde{a}_i \Bigg) \, .
\end{align*}
By construction, for each \(j\), there exists a \(T_{i,j}\) where \(\partial T_{i,j} = \tau_{i,j}\), therefore:
\begin{align*}
    \int_{\tau_{i,j}} \mathcal{J}^* \tilde{a}_i = \int_{T_{i,j}} \mathcal{J}^* b_i = \int_{\tau_{i,j}} \mathcal{J}^* a_i \, .
\end{align*}
Applying this to the difference between relative helicities computed in different gauges gives:
\begin{align*}
    \mathscr{H}_{a_i}(\beta_i, \ST_i) - \mathscr{H}_{\tilde{a}_i}(\beta_i, \ST_i) 
    = 0 \, .
\end{align*}
Hence, the relative helicity is independent of the gauge. Due to this, we now drop the subscript of \(\mathscr{H}\) that specifies the potential. 





Additionally, for \(f \in \I\) and an appropriate \(G \subseteq \Omega^1(\ST_i)\): 
\begin{align*}
\mathscr{H}( \beta_i, \ST_i ) - \mathscr{H}( (f^{-1})^* \beta_i, f(\ST_i ) ) = &H_G( \beta_i, \ST_i ) - H_{(f^{-1})^* G}( (f^{-1})^* \beta_i, f(\ST_i ) )  \\
&- \sum_{j=1}^{\ell_i} \int_{\sigma_{i,j}} a_i \int_{\tau_{i,j}}  a_i
 +
 \sum_{j=1}^{\ell_i} \int_{f(\sigma_{i,j})} (f^{-1})^*a_i \int_{f(\tau_{i,j})} (f^{-1})^* a_i
 \\
 = &H_G( \beta_i, \ST_i ) - H_{(f^{-1})^* G}( (f^{-1})^* \beta_i, f(\ST_i ) ) 
 \,  .
\end{align*}
We have the freedom to select a \(G\) without altering the relative helicity. Hence, we restrict primitives to a Biot-Savart gauge, as defined by Cantarella \cite{CANTARELLA20101127}. With this gauge Cantarella shows that \(H_G( \beta_i, \ST_i ) = H_{(f^{-1})^* G}( (f^{-1})^* \beta_i, f(\ST_i ) ) \). Hence, \(\mathscr{H}( \beta_i, \ST_i ) = \mathscr{H}( (f^{-1})^* \beta_i, f(\ST_i ) )\) must hold in any gauge.
\end{proof}



\begin{corollary}
    The objective and constraints in the variational problem \eqref{eq:ext} are independent of the chosen gauge. Namely, the problem is well-defined in terms of the magnetic field \(\beta_i\) and \(f \in \I\). 
\end{corollary}





\begin{lemma} \label{Lemma:gauge_exist}
    There exists at least one element \(a_i \in \mathscr{A}(\ST_i)\) such that \(da_i = \beta_i\) for any \(\beta_i \in \C\). 
\end{lemma}
\begin{proof}
There exists at least one primitive \(a_i \in  \Omega^1 (\ST_i)\) such that \(da_i = \beta_i\) (Lemma \ref{Lemma:Potential}). Additionally, \(a_i + s\) is also a primitive for \(\beta_i\) for any closed 1-form \(s \in  \Omega^1 (\ST_i)\). Hence, \([s]\) is any element of \(H^1_{\mathrm{dR}} (\ST_i)\) and therefore has a decomposition:
\begin{align*}
    [s] = \sum_{j=1}^{\ell_i} c_{i,j} [ s_{i,j}^* ] \, . 
\end{align*}
Pulling back to the boundary via inclusion gives
\begin{align*}
    \mathcal{J}^* [s] = \sum_j c_{i,j}  \mathcal{J}^* [s_{i,j}^*]  = \sum_j c_{i,j} [\sigma_{i,j}^*] \, . 
\end{align*}
Giving
\begin{align*}
    \int_{\sigma_{i,j}} \mathcal{J}^* (a_i + s) =  \int_{\sigma_{i,j}} \mathcal{J}^* a_i + \int_{\sigma_{i,j}}  \sum_{k=1}^{\ell_i} c_{i,k}  \sigma_{i,k}^*   = 
    c_{i,j} + \int_{\sigma_{i,j}} \mathcal{J}^* a_i    \, .
\end{align*}
Setting \(c_{i,j}\) so that the right-hand side is zero uniquely specifies \([s]\) and we have that \(a_i + s\) is an Amperian primitive for \(\beta_i\).
\end{proof}



This lemma is important because the relative helicity given any gauge always aligns with the helicity for an Amperian gauge. 



\begin{remark}\label{remark:helicity}
    Given a \(\ST_i\) with non-trivial genus and \(\beta_i \in \C\), the helicity \(H_{G} ( \beta_i, \ST_i)\) can take any value via a choice of \(G \subseteq \Omega^1(\ST_i) \) if \(\ell_i > 0\) and \(\psi_{i,1}, \dots, \psi_{i,\ell_i} \) are not all identically zero.
\end{remark}

We consider this remark in more detail below. Under the conditions in remark \ref{remark:helicity}, \(\beta_i \neq 0\) and \(\ell_i >0\). Consider two primitives of \(\beta_i\) given by \(a_i \in G\) and \(a_i + s \in G'\) for a closed 1-form \(s \in  \Omega^1(\ST_i)\). Both \(s\) and \(\mathcal{J}^* a_i\) are closed on the boundary so formula \eqref{eq:compute_2} is applicable:
\begin{align*}
    H_G (\beta_i, \ST_i) - H_{G'} (\beta_i, \ST_i)  =  \int_{\partial \ST_i}  \!  s \wedge a_i &=  - \sum_{j=1}^{\ell_i}  \left(\int_{\sigma_{i,j}} \! s    \int_{\tau_{i,j}}  \! a_i -\int_{\sigma_{i,j}}  \! a_i \int_{\tau_{i,j}}  \! s  \right) 
    %\\
    %&
    =  - \sum_{j=1}^{\ell_i}  \left( \psi_{i,j} \int_{\sigma_{i,j}}  \!  s    \right) 
     \, .
\end{align*}
Therefore, if any of the fluxes \(\psi_{i,j}\) are non-trivial, we may select \(s\) to make the difference in helicity any value (we saw a comparable situation in the proof of Lemma \ref{Lemma:gauge_exist}). 


Therefore, if we consider problem \eqref{prob:min1} without constraining the potential for non-trivial flux: (1) the constraint \(H_{\Omega^1(\ST_i)} (\beta_i , \ST_i) = h_i\) is not well-defined as a function of \(\beta_i\) and (2) the helicity can take any value, so this constraint does not restrict \(\beta_i\) in any way. Similarly, for problem \eqref{eq:ext} with a helicity constraint rather than a relative helicity constraint for non-trivial genus and \(\psi_{i,1}, \dots, \psi_{i,\ell_i} \) not identically zero. 





\subsection{Jump Conditions}

The boundary and jump conditions that we are using in equations \eqref{prob:MRxMHD} are reported as a special case of the following conditions given by Bruno and Laurence \cite{bruno_laurence_1996}:
\begin{align*}
    0 & = \llbracket B \cdot N \rrbracket \, ,  \\
    0 & = \llbracket B^2 + 2p \rrbracket N - 2 (B \cdot N) \llbracket B \rrbracket \, \, ,
\end{align*}
on \(\partial \ST_i\) where \(N\) is the outward unit normal vector field to the boundary. The two terms in the second equation here are orthogonal, namely, the first term is parallel to \(N\) and the second term lies in the tangent space of \(\partial \ST_i\). Therefore, both terms must be \(0\) independent of each-other. So either \(B\) is continuous across the boundary and \(\llbracket p \rrbracket = 0\), or:
\begin{align*}
    0 & = B \cdot N \, ,  \\
    0 & = \llbracket B^2 + 2p \rrbracket \, .
\end{align*}
Under the assumption that \(\llbracket p \rrbracket \neq 0\) across our interfaces, the original conditions from Bruno and Laurence are equivalent to the boundary conditions in \eqref{prob:MRxMHD}. 



\subsection{Lagrange Multipliers} \label{sub:LM}



We know that any \(\beta_i \in \C\) for \(\ell_i>0\) can be written as:
\begin{align} \label{eq:b_form}
    \beta_i = d \eta_i + \sum_{j=1}^{\ell_i}  \psi_{i,j}  T_{i,j}^*  
    = d \Bigg( \eta_i + \sum_{j=1}^{\ell_i}  \psi_{i,j} \Gamma_{i,j}  \Bigg)\, ,
\end{align}
for \(\eta_i \in  \Omega_D^1(\ST_i)\), \( T_{i,j}^*  \in \mathcal{H}_D^2 (\ST_i) \), and \(\Gamma_{i,j} \in  \Omega^1 (\ST_i)\). The second equality holds because there always exists a primitive \(\Gamma_{i,j}\) for \( T_{i,j}^* \) (Lemma \ref{Lemma:Potential}). 



In the variational formulation \eqref{eq:ext} we incorporate the flux constraints directly into the objective. We see from equation \eqref{eq:b_form} that the harmonic component of \(\beta_i\) (which determine the flux of \(\beta_i\) through any relative \(2\)-chain) is fixed once the coefficients \(\psi_{i,1},\dots,\psi_{i,\ell_i}\) are prescribed, provided \(\ell_i>0\). Consequently, the objective and helicity are functions of \(f\in\I\) and \(d\eta_i\). The potential \(\eta_i\) is not always unique, different \(\eta_i\) may have the same exterior derivative \(\beta_i\), hence the same value for the objective and relative helicity. To incorporate variations of the potentials we formulate the optimisation over \(f\) and \(\eta_i\), giving
\begin{align}\label{prob:start_point}
\begin{gathered}
    \underset{\substack{  \eta_i \in  \Omega^1_D(\ST_i)  \\  f\in\I} }{ \text{c. p.} } \sum_{i=1}^n \Bigg( \frac{1}{2}\Big\| d\eta_i + \sum_{j=1}^{\ell_i} \psi_{i,j}\,T_{i,j}^* \Big\|_{L^2}^2 - p_i\,|\ST_i| \Bigg) \\
  \text{subject to}\qquad 
   \mathscr{H}\!\Big( d\eta_i + \sum_{j=1}^{\ell_i} \psi_{i,j}\,T_{i,j}^* ,\; \ST_i\Big) = h_i,\qquad i=1,\dots,n \, . 
\end{gathered}
\end{align}
If \(\ell_i=0\) there is no harmonic contribution in the displayed expressions.

For the remainder of this section we freeze the deformation \(f\). The variational problem reduces to finding admissible potentials \(\eta_i\). A standard way to treat the constraint in \eqref{prob:start_point} is to introduce Lagrange multipliers. If the constraint map is a submersion, the regular value theorem implies that the preimage of the prescribed value is a submanifold whose tangent space at each point is characterised as the kernel of the linearised constraint. Imposing that admissible variations lie in this tangent space and the Lagrange multipliers theorem yields Euler--Lagrange equations for the constrained optimisation problem.

Let's consider a regular value theorem for the relative helicity constraints in problem \eqref{prob:start_point}. Our definition of a critical point requires us to determine the Fr\'echet derivative in a \(C^\infty\) topology, hence, we consider \(\Omega^1_D(\ST_i)\) with this topology. 


We can view \(\Omega_D^1(\ST_i)\) as the kernel of a continuous linear map given by the pullback \(\mathcal{J}^* : \Omega^1(\ST_i) \to \Omega^1(\partial \ST_i)  \). So \(\Omega_D^1(\ST_i)\) is a Fr\'echet space because it is a closed subspace of the Fr\'echet space \(\Omega^1(\ST_i)\). 


Our constraint function is \(g_i:  \Omega^1_D(\ST_i) \to \R\), given by \(g_i : \eta_i \mapsto \mathscr{H} ( \beta_i, \ST_i) \). For any \(\eta_i \in \Omega^1_D(\ST_i)\), the \textit{tangent spaces} \(T_{\eta_i}  \Omega^1_D(\ST_i) \) and \( T_{g_i(\eta_i)} \R\) may be canonically identified with \( \Omega^1_D(\ST_i) \) and \( \R\) respectively \cite{kriegl1997convenient}. A standard computation gives us that \(g_i\) is continuously Fréchet differentiable, and therefore has a directional derivative (proof in Appendix \ref{app:differentiable}). The derivative of \(g_i\) with respect to \(\eta_i\) in the direction \(\eta_i' \in  \Omega^1_D(\ST_i)\) is denoted \(d_{\eta_i} g_i(\eta_i')\). Given this, we may adapt a regular value theorem given by Gl\"ockner (theorem D \cite{glockner2015fundamentals}) to our setting. 




\begin{theorem}[Gl\"ockner, theorem D adaptation \cite{glockner2015fundamentals}]\label{theorem:Glock}
    Consider a smooth submersion \(g_i : \Omega_D^1(\ST_i) \to \R\) that has a Fr\'echet derivative \(d_{\eta_i} g_i : \Omega_D^1(\ST_i) \to  \R \) at \(\eta_i \in \Omega_D^1(\ST_i)\). If \(d_{\eta_i} g_i\) is surjective, then \(\ker (d_{\eta_i} g_i) = T_{\eta_i} g_i^{-1}(y)\) for any \(y = g_i(\eta_i)\). 
\end{theorem}
Note that Gl\"ockner's result relies on \(d_{\eta_i} g_i\) having a \emph{complemented} kernel at \(\eta_i \in \Omega_D^1(\ST_i)\), which is true when \(g_i\) has a one dimensional codomain \cite{glockner2015fundamentals} (Appendix \ref{app:Complemented}). 




Since \(g_i\) is Fréchet differentiable, \(d_{\eta_i} g_i\) coincides with the Gâteaux derivative \cite{zeidler1986nonlinear}. Hence, we proceed with the Gâteaux derivative.


Consider \(\eta_{i,t} = \eta_i + t \eta_i' + o(t)  \in  \Omega^1_D(\ST_i)\) that is differentiable in \(t \in [-\epsilon, \epsilon]\), \(\epsilon > 0\). A partial derivative of \(\eta_{i,t}\) with respect to \(t\) is denoted \(\partial_t\) and \(\eta_i' = \partial_t \eta_{i,t}|_{t=0}\) is an element of the tangent space \(T_{\eta_i}  \Omega^1_D (\ST_i)\). So:
\begin{align} 
    d_{\eta_i} g_i (\eta_i')  &=  \partial_t g_i (\eta_{i,t}) |_{t=0} 
     \label{eq:deriv}
    %& 
    = 2  \int_{\ST_i}  \eta_i' \wedge \beta_i 
    -
    \sum_{j=1}^{\ell_i} \psi_{i,j} \int_{\sigma_{i,j}}  \eta_i' \, .
\end{align}
We have used the divergence theorem in our simplification. 




Assume that \(\beta_i\) is not identically zero on \(\ST_i\). Then there exists a point \(x_0 \in \ST_i\) such that \(\beta_i(x_0) \neq 0\). By continuity, \(\beta_i\) remains non-zero on some ball \(\B(r) \ssubset \ST_i\) with radius \(r>0\). Since \(\partial \ST_i\) is smooth, \(\ST_i\) satisfies the \(\epsilon\)-cone property \cite{henrot2018shape}; hence there exists a ball \(\B(r_0) \ssubset \B(r) \cap \ST_i\) for some \(r_0>0\). In particular, \(\beta_i\) is non-zero throughout a compact subset of \(\ST_i\). It follows that
\[
\|\beta_i\|_{L^2(\ST_i)}^2 \geq \|\beta_i|_{\B(r_0)}\|_{L^2(\B(r_0))}^2 > 0 \, .
\]
Let \(\chi_{\B}\) denote the characteristic function of \(\B(r_0)\) on \(\ST_i\). For each \(k \in \mathbb{N}\), let \(\chi_k \colon \ST_i \to \mathbb{R}\) be a smooth function in \(\Omega^0(\ST_i)\) satisfying \(\chi_k|_{\partial \ST_i}=0\), and choose \(\chi_k \to \chi_{\B}\) pointwise as \(k \to \infty\). Consider the variation \(
\eta_i' = \chi_k \,\star \beta_i \) in the directional derivative \eqref{eq:deriv}. Then
\begin{align*}
    d_{\eta_i} g_i(\eta_i')
    &= 2 \! \int_{\ST_i} \! \! \chi_k \,\star \beta_i \wedge \beta_i
       + \sum_{j=1}^{\ell_i} \psi_{i,j} \int_{\sigma_{i,j}} \! \! \mathcal{J}^*(\chi_k \star \beta_i) 
    \xrightarrow{k \to \infty}
       2 \! \int_{\B(r_0)} \! \! \star \beta_i \wedge \beta_i
       = 2\|\beta_i|_{\B(r_0)}\|_{L^2(\B(r_0))}^2
       > 0 \, .
\end{align*}
Therefore, for \(k\) sufficiently large, \(
d_{\eta_i} g_i(\chi_k \,\star \beta_i) > 0 \). By linearity of the derivative, \(d_{\eta_i} g_i\) is surjective onto \(\mathbb{R}\). Consequently, every \(h_i\) is a regular value for the relative helicity constraint \(g_i(\eta_i) = h_i\), provided \(\beta_i \neq 0\) (Theorem \ref{theorem:Glock}). Hence the regular Lagrange multiplier condition gives necessary and sufficient conditions for a solution of the original problem (Appendix \ref{app:LM}).

The same argument applies verbatim to conventional helicity, up to any additional constraints on the potential.

Applying the method of Lagrange multipliers therefore yields the following necessary and sufficient conditions for a solution:
\begin{align*}
    \underset{\substack{\eta_i \in \Omega^1_D(\ST_i) \\ \mu_i \in \R}}{\mathrm{c.p.}}
    \sum_{i=1}^n
    \left(
        \frac{1}{2}\Big\| d\eta_i + \sum_j \psi_{i,j} T_{i,j}^* \Big\|_{L^2}^2
        - p_i \myabs{\ST_i}
        - \frac{\mu_i}{2}\big(\mathscr{H}(\beta_i,\ST_i)-h_i\big)
    \right),
\end{align*}
for Lagrange multipliers \(\mu_1,\dots,\mu_n \in \R\), assuming \(\beta_i \neq 0\). As each domain may be considered independently we obtain the following lemma.

\begin{lemma}\label{Lemma:replace}
    Let \((\ST_1,\dots,\ST_n)\) be a partition of \(\ST\), and assume that \(\beta \neq 0\). Fix the constants \(\psi_{i,j}\), \(p_i\), and \(h_i\). Given \(f \in \I\) and the volume form \(f^*\varpi\), problem \eqref{eq:ext} reduces, for each \(i\) with \(\ell_i>0\), to
    \begin{align}\label{eq:ext_lagrange}
        \underset{\substack{\eta_i \in \Omega^1_D(\ST_i) \\ \mu_i \in \R}}{\mathrm{c.p.}}
        \left(
            \frac{1}{2}\Big\| d\eta_i + \sum_{j=1}^{\ell_i}\psi_{i,j} T_{i,j}^* \Big\|_{L^2}^2
            - \frac{\mu_i}{2}\big(\mathscr{H}(\beta_i,\ST_i)-h_i\big)
        \right).
    \end{align}
    For \(i\) where \(\ell_i=0\), the harmonic terms vanish.
\end{lemma}





\section{Proof of Main Results}\label{section:Main Proofs}




\subsection{Proof of Theorem \ref{thr:iff_simple}}


We start from the original formulation for our variational problem \eqref{eq:ext} with \(\beta_i \neq 0\). Consider an \(f_t \in \I\) for each \(t \in [0,1]\) that is differentiable in \(t\) where \(f_1\) is the identity, \(f_0 = f \). The objective function, or magnetic energy is:
\begin{align*}
    \begin{gathered} 
    E(f_t, \beta_{1,t}, \dots, \beta _{n,t})  =  \sum_{i=1}^n \Big( \frac{1}{2} \big\| \beta_{i,t}\big\|_{L^2}^2 -  p_i \myabs{ \ST_i } \Big) \, , 
    \end{gathered}
\end{align*}
where \( \beta_{i,t} \in \C\) for each \(t \in [0,1]\) and \(a_{i,t}\) is continuously differentiable in \(t\), where \(da_{i,t} = \beta_{i,t}\). Let \(a_{i,0} = a_i\) and \(\beta_{i,0}= \beta _i\). 


We compute the derivative of \(E\) that is needed to define critical points: namely, a derivative that is Fr\'echet with respect to \(\eta\) and G\^ateaux with respect to \(f\). Since \(E\) is Fr\'echet differentiable in \(\eta\), this derivative may be obtained by differentiating \(E\) along all paths \(P_t \coloneqq (f_t, \beta_{1,t}, \dots, \beta_{n,t})\) that satisfy the flux and relative helicity constraints in Problem~\eqref{eq:ext}. We refer to such paths as \textit{allowable paths}. A point is critical for the magnetic energy precisely when the first variation of the energy vanishes at \(t=0\) along every allowable path. Hence, a necessary and sufficient condition for \(P_0\) to be a critical point is 
\begin{align*}
    0 = \partial_t E(P_t) |_{t=0}  =
     \left(  \partial_{f_t} E(f')  + \sum_{i=1}^n \partial_{ \beta_{i,t}} E( \beta_i') \right) \Bigg|_{t=0}  \, ,
\end{align*}
where \((f', \beta_{1}', \dots, \beta_{n}') = ( \partial_t f_{t}|_{t=0}, \partial_t \beta_{1,t}|_{t=0}      , \dots, \partial_t \beta_{n,t}|_{t=0} )
\) as in Section \ref{sub:LM}. Namely, \(\partial_{f_t} E(f')\) is the partial derivative with respect to \(f_t\), in the direction \(f'\) at the point \(P_t\) for each \(t\). So the conditions for \(P_0\) to be a critical point are exactly:
\begin{align} \label{eq:crit_cond_1}
    0 = d_{f_t} E(f') |_{t=0} \, , \qquad 
    0 = d_{\beta_{i,t}} E( \beta_i') |_{t=0}  \, , \; \;  i = 1, \dots, n\, ,
\end{align}
for all allowable paths passing through \(P_0\). As \(\beta_i \neq 0\), the second condition is equivalent to finding the critical points of \(E\) at a fixed \(f \in \I\). Recall from the previous section that this is equivalent to a Lagrange multipliers problem (Lemma \ref{Lemma:replace}). Namely, for each \(i = 1, \dots, n\) we intend to determine critical points of 
\begin{align*}
    L(\eta_i, \mu_i) \coloneqq \frac{1}{2} \big\| d a_i  \big\|_{L^2}^2 - \frac{\mu_i}{2} \left(   \mathscr{H}( d a_i,  \ST_i) -  h_i      \right) \, , & &
    a_i = \eta_i + \sum_{j=1}^{\ell_i}  \psi_{i,j} \Gamma_{i,j} \, ,
\end{align*}
where the harmonic terms are zero whenever \(\ell_i = 0\). So
\begin{align*}
    0 = d_{\beta_{i}} E(\beta_i') &\iff 0 = \partial_t L(\eta_{i,t}, \mu_{i,t}) |_{t=0} = d_{\eta_{i,t}} L(\eta_i') |_{t=0} +  d_{\mu_{i,t}} L(\mu_i') |_{t=0} \\ &\iff  0= d_{\eta_{i,t}} L(\eta_i') |_{t=0} \quad \text{and} \quad  0= d_{\mu_{i,t}} L(\mu_i') |_{t=0} \, ,
\end{align*}
for \(\mu_{i,t} \in \R\) that are differentiable in \(t\) where \(\mu_i' = \partial_t \mu_{i,t}|_{t=0}\). The latter equation here is given by:
\begin{align*}
     0= d_{\mu_{i,t}} L(\mu_i') |_{t=0} = \partial_t L(\eta_{i}, \mu_{i,t}) |_{t=0}  =   - \frac{\mu_i'}{2} \left(   \mathscr{H}( \beta_i,  \ST_i) -  h_i      \right)  \, ,
\end{align*}
for all \(i\) and \(\mu_i' \in \R\). Hence, this is equivalent to \(\mathscr{H}( \beta_i,  \ST_i) =  h_i\) for all \(i\). For the other derivative of \(L\) we have that:
\begin{align*}
d_{\eta_{i,t}} L(\eta_i') |_{t=0} = \partial_t L(\eta_{i,t}, \mu_{i}) |_{t=0}   = \langle d\eta_i',  \beta_i\rangle_{L^2} -  \frac{\mu_i}{2} \int_{\ST_i} (\eta'_i \wedge \beta_i + a_i \wedge d\eta_i')    = \langle \eta_i' , \star (d \star \beta_i - \mu_i \beta_i) \rangle_{L^2} \, .
\end{align*}
The vanishing of \(d_{\eta_{i,t}} L(\eta_i') |_{t=0}\) for all \(\eta'_i \in  \Omega^1_D(\ST_i)\) implies that \(\star (d \star \beta_i - \mu_i \beta_i) = 0\) on the interior of \(\ST_i\). Since \(\beta_i\) is smooth, the quantity \(\star (d \star\beta_i - \mu_i\beta_i)\) is continuous. As this vanishes on the interior of \(\ST_i\), it must vanish everywhere on \(\ST_i\). Therefore, we conclude that \(d \star\beta_i = \mu_i\beta_i\), which is a Beltrami equation in \(\beta_i\) for each \(i\). Note that the Hodge star operator encodes the volume form \(f^* \varpi\) on \(\ST_i\). Upon pushing forward via \(f\), we obtain a Beltrami equation for the elements of \(\mathcal{C}(f(\ST_i))\) with a Euclidean metric. This result is highlighted in Remark~\ref{remark:MHD}.






The conditions for a critical point in \eqref{eq:crit_cond_1} are equivalently written as:
\begin{align} \label{eq:crit_cond_2}
    0 = d_{f_t} E(f_t) |_{t=0} = \partial_t \left(   \sum_{i=1}^n \Big( \frac{1}{2} \big\| \beta_{i,t}\big\|_{L^2}^2 -  p_i \myabs{ \ST_i } \Big)      \right) \Bigg|_{t=0} \, , \quad 
    d \star \beta_{i} = \mu_{i} \beta_{i} \, , \quad \mathscr{H}( \beta_{i,t},  \ST_i) =  h_i \, , 
\end{align}
for all \(i\), where the harmonic terms in the HMF decomposition of \(\beta_{i,t}\) are given, as in equation \eqref{eq:b_form}. Note that the relative helicity holds for all \(t\) by definition of the admissible paths. We interpret this problem as: finding critical points of the magnetic energy under variations of the boundary through all paths \(P_t\) satisfying the prescribed flux and helicity constraints while \(P_0\) satisfies the Beltrami equation. 


Consider a derivative \(d/dt f_t : \ST \to T f_t(\ST) \). Composing this with a directional derivative \(d_t f_{t}^{-1} :T f_t(\ST) \to T (\ST) \) gives us a vector field \(V(x) \coloneqq d_t f_{t}^{-1} \ (d/dt (f_t(x))  )|_{t=0}\), for \(x \in \ST\), namely \(V:\ST \to T(\ST)\). With this we have the following formula, essentially Cartan's formula,
\begin{align*}
    \partial_t (f^*_t \omega_t) |_{t=0} = f^*_0 (\mathcal{L}_V (\omega_0) +  \partial_t \omega_t |_{t=0}) =  i_V (d\omega_0) + d i_V (\omega_0) + \partial_t \omega_t |_{t=0} \, ,
\end{align*}
for any \(\omega_t \in  \Omega^k (f_t(\ST_i) ) \) where \(\mathcal{L}_V\) is a Lie derivative with respect to \(V\). 

It is useful to perform some operations on \(f_t(\ST_i)\) because the metric here is the Euclidean metric and independent of \(t\). Let \(\epsilon_i \coloneqq f_t^* \partial_t ((f_t^{-1})^* \eta_{i,t}) |_{t=0}\) where \( (f_t^{-1})^* \eta_{i,t} \in  \Omega^1 (f_t(\ST_i))\) and \( (f_t^{-1})^*\) is a pushforward by \(f_t\). Then we have commutation of the Hodge star with a derivative: 
\begin{align*}
    f^*_t \partial_t (f_t^{-1})^* \star \eta_{i,t} |_{t=0} = \star f^*_t  \partial_t (f_t^{-1})^* \eta_{i,t} |_{t=0} = \star \epsilon_i \, ,
\end{align*}
where \(\star\) is the Hodge star in \(\ST_i\) with volume form \(f_t^* \varpi\). We use the notation \(a_i' = \partial_t \eta_{i,t}|_{t=0}\). This gives the relations:
\begin{align*}
    a_i' &= \partial_t \eta_{i,t}|_{t=0} = \partial_t  f_t^*  ((f_t^{-1})^* \eta_{i,t}) |_{t=0} = 
    i_V (d  \eta_{i}  ) + d i_V (  \eta_{i}  ) + \epsilon_i \, , \\
    (\star a_i)' &= \partial_t \star \eta_{i,t}|_{t=0} = \partial_t   f_t^*  (  (f_t^{-1})^*    \star \eta_{i,t} ) |_{t=0} = i_V ( d  \star \eta_{i}  ) + d i_V (  \star \eta_{i}  ) + \star \epsilon_i \, .
\end{align*}
Let's return to our problem, given by equation \eqref{eq:crit_cond_2}. For simplicity, we deal with each term in the derivative with respect to \(f_t\) independently until the end of our computation. First, let's consider \(\| \beta_{i,t}\|_{L^2}^2\) for any \(i \in \{1, 2, \dots, n\} \):
\begin{align*}
    \partial_t    \big\| \beta_{i,t}\big\|_{L^2}^2    \Big|_{t=0}   = 2  \langle da_i, d \epsilon_i \rangle + \! \int_{\ST_i} \! \! d i_V (da_i \wedge \star da_i) = 2 \mu_i \! \int_{\ST_i} \! \! \epsilon_i \wedge\beta_i + \! \int_{\partial \ST_i}  \! \! \mathcal{J}^* ( 2 \epsilon_i \wedge \star\beta_i +  i_V (da_i \wedge \star da_i ) ) \, .
\end{align*}
We have used the Beltrami equation and divergence theorem in this computation. Note that these computations hold for any \(\mu_i \). Additionally, the relative helicity is constant in \(t\), so:
\begin{align*}
    0 & = \partial_t  \left( \int_{\ST_i} a_{i,t} \wedge \beta_{i,t}  - \sum_{j=1}^{\ell_i} \int_{\sigma_{i,j}} \mathcal{J}^* a_{i,t} \int_{\tau_{i,j}} \mathcal{J}^* a_{i,t}  \right) \Bigg|_{t=0}\\
    & =  \int_{\ST_i} \epsilon_i \wedge\beta_i + a_i \wedge d \epsilon_i + \mathcal{L}_V(a_i \wedge\beta_i) - \sum_{j=1}^{\ell_i} \psi_{i,j}  \int_{\sigma_{i,j}} \mathcal{J}^* a_i'    \\ 
     2 \int_{\ST_i}   \epsilon_i \wedge\beta_i & =  \int_{\partial \ST_i} \mathcal{J}^* (  a_i \wedge \epsilon_i -   i_V(a_i \wedge\beta_i) ) + \sum_{j=1}^{\ell_i} \psi_{i,j}  \int_{\sigma_{i,j}} \mathcal{J}^* a_i' \, , 
\end{align*}
where we have used the divergence theorem. Using this equality in the derivative of \(\| \beta_{i,t}\|_{L^2}^2\):
\begin{align*}
    \partial_t    \big\| \beta_{i,t}\big\|_{L^2}^2    \Big|_{t=0}  \!  =   \int_{\partial \ST_i}  \! \! \mathcal{J}^* ( \mu_i a_i \wedge \epsilon_i -   \mu_i i_V(a_i \wedge\beta_i)   +    2 \epsilon_i \wedge \star\beta_i +  i_V (da_i \wedge \star da_i) ) + \mu_i \sum_{j=1}^{\ell_i} \psi_{i,j}  \int_{\sigma_{i,j}} \! \! \mathcal{J}^* a_i' \, .
\end{align*}
We see that this variation with respect to \(t\) reduces to integrals along the boundary \(\partial \ST_i\). With some rearrangement of terms we arrive at:
\begin{align*}
    \partial_t    \big\| \beta_{i,t}\big\|_{L^2}^2    \Big|_{t=0}  & =   \int_{\partial \ST_i} \mathcal{J}^* (   (i_V\beta_i + \epsilon_i ) \wedge ( 2 \star\beta_i - \mu_i a_i) +\beta_i (i_V \star\beta_i - \mu i_V a_i) - i_V( \beta_i \wedge \star \beta_i) ) \\ & \qquad  + \mu_i \sum_{j=1}^{\ell_i} \psi_{i,j}  \int_{\sigma_{i,j}} \mathcal{J}^* a_i' \, .
\end{align*}
Note that \(\mathcal{J}^* \beta_i = t \beta_i = 0\), so the second term disappears. This normal condition on \(\beta_i\) also implies the following:
\begin{align*}
    0 = \mathcal{J}^* \beta_{i,t} \implies 0 = \mathcal{J}^* (d \epsilon_i  + \mathcal{L}_V(\beta_i) ) = d \mathcal{J}^* ( \epsilon_i  + i_V(\beta_i) )  \, . 
\end{align*}
Additionally, for the first term we have that both terms in the wedge product are closed in \( \Omega^1 (\partial \ST_i)\). Applying equation \eqref{eq:compute_2}, gives
\begin{align*}
    \partial_t    \big\| \beta_{i,t}\big\|_{L^2}^2  &= - \sum_{j=1}^{\ell_i}  \left(   \int_{\sigma_{i,j}} \! (i_V\beta_i + \epsilon_i )    \int_{\tau_{i,j}} \!  ( 2 \star\beta_i - \mu_i a_i) -\int_{\sigma_{i,j}} \!  ( 2 \star\beta_i - \mu a_i) \int_{\tau_{i,j}} \!  (i_V\beta_i + \epsilon_i ) \right) \\ & \qquad -
    \int_{\partial \ST_i}  (i_V( \beta_i) \wedge \star\beta_i)  + \mu_i \sum_{j=1}^{\ell_i} \psi_{i,j} \int_{\sigma_{i,j}}   a_i'  \, .
\end{align*}
Using the divergence theorem 
\begin{align*}
    \int_{\tau_{i,j}} \mathcal{J}^* (i_V \beta_i + \epsilon_i ) = \int_{\tau_{i,j}} \mathcal{J}^* a_i' =  \int_{T_{i,j}} \mathcal{J}^*\beta_i' = \partial_t \psi_{i,j} |_{t=0} = 0 \, .
\end{align*}
So this term vanishes as flux is held constant. And
\begin{align*}
    \sum_{j=1}^{\ell_i}    \int_{\sigma_{i,j}} \mathcal{J}^*(i_V\beta_i + \epsilon_i )    \int_{\tau_{i,j}} \mathcal{J}^*( 2 \star\beta_i - \mu_i a_i)  = 
    \mu_i \sum_{j=1}^{\ell_i}   \psi_{i,j}   \int_{\sigma_{i,j}} \mathcal{J}^* a_i'   \, .
\end{align*}
So we are left with:
\begin{align*}
     \partial_t    \big\| \beta_{i,t}\big\|_{L^2}^2    \Big|_{t=0} & =   -
    \int_{\partial \ST_i} \mathcal{J}^* (i_V( \beta_i \wedge \star\beta_i) )  \, .
\end{align*}
Let's also simplify the \( \myabs{ \ST_i } \) terms in \(d_{f_t} E(f') |_{t=0}\):
\begin{align*}
    \partial_t \myabs{ \ST_i } |_{t=0} = \int_{\ST_i} \partial_t f^*_t \varpi |_{t=0} 
    = \int_{\ST_i} f^* \mathcal{L}_V \varpi =  \int_{\ST_i} d i_V \varpi =  \int_{\partial \ST_i} i_V \varpi \, . 
\end{align*}
Therefore the derivative \(d_{f_t} E(f') |_{t=0}\) is:
\begin{align*}
    0 = d_{f_t} E(f') |_{t=0} &= \partial_t \left(   \sum_{i=1}^n \Big( \frac{1}{2} \big\| \beta_{i,t}\big\|_{L^2}^2 -  p_i \myabs{ \ST_i } \Big)      \right) \Bigg|_{t=0} 
    =  -  \sum_{i=1}^n   
    \int_{\partial \ST_i}  \mathcal{J}^* (i_V \varpi) \Big( \frac{1}{2}  \| \beta_i\|^2_E \,  +  p_i  \Big) \, .
\end{align*}
This is required for all \(V\) with \(V|_{\partial \ST}=0\). We see that only the normal component in the image of \(V\) is relevant in this integral. Noticing that the normals are reversed on either side on any interface we have an equivalent condition:
\begin{align*}
    0 = \Big\llbracket \| \beta \|^2_E + 2 p \Big\rrbracket \Big|_{I_j} \, ,
\end{align*}
for all interfaces \(I_j\). So the necessary and sufficient conditions for a critical point with \(\beta_i \neq 0\) are equivalent to 
\begin{align*} 
    0 = \Big\llbracket \|\beta \|^2_E + 2 p \Big\rrbracket \Big|_{I_j} \, , \qquad 
    d \star\beta_i = \mu_i\beta_i \, , \qquad \mathscr{H}( \beta_i,  \ST_i) =  h_i \, ,
\end{align*}
where the flux of \(\beta_i\) through any element of \(H_2(\ST_i, \partial \ST_i)\) is specified. This concludes the proof of Theorem \ref{thr:iff_simple}. \qed




It is worth noting that the term:
\begin{align}\label{eq:floating}
    \mu_i \sum_{j=1}^{\ell_i} \psi_{i,j}  \int_{\sigma_{i,j}} \mathcal{J}^* a_i' \,  ,
\end{align}
in the derivative of \(\|\beta_{i,t}\|^2_{L^2}\) is cancelled by the corresponding term in the derivative of the relative helicity. If we were to replace the relative helicity with helicity in problem \eqref{eq:ext}, without any gauge restriction, and phrasing the problem in \(a_i\) to ensure the corresponding functions are well-defined, then the term \eqref{eq:floating} would remain in the derivative of \(\|\beta_{i,t}\|^2_{L^2}\). In this case we expect that a gauge choice will ensure that helicity is a well-defined function on \(\beta_i\) and make expression \eqref{eq:floating} equal to zero, allowing us to recover a similar result to Theorem \ref{thr:iff_simple}. As we know, this is always possible with an Amperian gauge. 






\section{Characterisation for small \texorpdfstring{$|\mu_i|$}{Lagrange Multiplier}}\label{section:Characterisation and Local uniqueness}


Consider problem \eqref{problem:a_only} and fix some \(i \in \{1,2,\dots, n\}\). Our goal here is to show that for small enough \(|\mu_i|\), a solution to problem \eqref{problem:a_only} is a local minimum of the magnetic energy. We already know from the previous sections that extrema to this problem exist and the Beltrami equations provide necessary and sufficient conditions for any extrema for problem \eqref{problem:a_only} for \(\beta_i \neq 0\).

Let's look to see if \(\| \beta_i  \|_{L^2}^2\) is minimised, maximised, or a saddle point. We consider the equivalent setup provided in problem \eqref{eq:ext_lagrange}. We perform a second derivative test for all paths \(\eta_{i,t} = \eta_i + t \eta_i' + t^2 \eta_i'' + o(t^2) \in \Omega^1_D( \ST_i )\) that satisfy the helicity constraint. Let \((\eta_i, \mu_i)\) denote a critical point of problem \eqref{eq:ext_lagrange}, we have that (utilising a result in \cite{exist_gerner}):
\begin{align*}
    \partial_t^2  \left( \frac{1}{2} \big\| d a_{i,t} \big\|_{L^2}^2 - \frac{\mu_i}{2} \left(   H_G (\beta_{i,t},  \ST_i) -  h_i      \right) \right) \Bigg|_{t=0} 
    & = \partial_t  \left( \langle da_{i,t}' , da_{i,t} \rangle
    - \mu_i \int_{\ST_i } a_{i,t}' \wedge \beta _{i,t} 
    \right)\Bigg|_{t=0} \\
    & =    \langle 2 d \eta_i'' , \beta_i  \rangle + \langle d \eta_i' , d \eta_i' \rangle
    - \mu_i \int_{\ST } 2 \eta_i'' \wedge \beta_i  + \eta_i' \wedge d \eta_i' 
    \\
    & = \| d \eta_i' \|^2_{L^2 }  -  \mu_i H_G (\beta_i', \ST_i) \\ & \geq \mu_{\pm} |H_G (\beta_i', \ST_i)|  -  \mu H_G (\beta_i', \ST_i) \, ,
\end{align*}
which is positive if \(\mu_{\pm} > |\mu_i|\) where \(\mu_{\pm} = \min \{ |\mu_+|,|\mu_- |\} \). Here \(\mu_+\) and \(\mu_-\) as the smallest positive and largest negative respective eigenvalues for problem \eqref{problem:a_only} with zero flux for a given \(i\). Hence, this condition guarantees that (for \(\mu_i\) small enough) the 2-form \(\beta_i \) is a local minimiser of the magnetic energy.







\section{Conclusion}\label{section:Conclusion}



In this work, we developed a variational framework for solutions to the multi-region relaxed magnetohydrodynamics (MRxMHD) equations on compact, oriented domains embedded in \(\mathbb{R}^3\) with smooth boundary. We introduced a definition of relative helicity and proved its gauge invariance. We demonstrated the existence of an Amperian gauge, noting that the relative helicity aligns with conventional helicity in this gauge. Within each subdomain, we imposed physically motivated conditions, including divergence-free fields, boundary tangency, and prescribed relative helicity, flux, and pressure constraints.

We showed that under these constraints non-zero critical points of the magnetic energy correspond precisely to solutions of the MRxMHD equilibrium equations. We identified a gauge condition that had not previously been noted in the literature. In the presence of a pressure jump, we established the equivalence between the MRxMHD formulation and a stronger formulation by Bruno and Laurence.


Our results extend the MRxMHD model beyond previously considered toroidal geometries. 




\section{Acknowledgements}\label{section:Acknowledgements}
This research was supported by an Australian Government Research Training Program (RTP) Scholarship at The University of Western Australia. The authors gratefully acknowledge this funding assistance.




%\section{Notes for Write-up}
%Below this point are notes that the above write-up is based on.
%\input{Chapter All Notes/Main}


\bibliographystyle{abbrv}
\bibliography{ref}


\newpage
\section*{Appendix}
\appendix



\section{De Rham Cohomology and Singular Homology}\label{section:Homology cohomology construction}


\subsection{Construction} \label{sub:construct}

For each \(k \in \mathbb{Z}\), the set of all smooth \(k\)-forms on \(\ST_i\) is \(\Omega^k(\ST_i)\). This set is an abelian group with a homomorphism given by the exterior derivative \(d :\Omega^k(\ST_i)\longrightarrow\Omega^{k+1}(\ST_i)\) satisfying \(d^2 = 0\). For \(k < 0\) or \(k > 3\) we have that \(\Omega^k(\ST_i) \cong 0\). The sequence \(\dots , \, \Omega^0(\ST_i), \, \Omega^1(\ST_i), \, \dots\) with operations \(d\) is commonly called the \emph{de Rham complex} of \(\ST_i\). The \(k\)-th \emph{de Rham cohomology} group of \(\ST_i\) is the \(k\)-th cohomology of the de Rham complex,
\[
H^k_{\mathrm{dR}}(\ST_i) \coloneqq
\frac{\ker\big(d:\Omega^k(\ST_i)\to\Omega^{k+1}(\ST_i)\big)}
     {\operatorname{im}\big(d:\Omega^{k-1}(\ST_i)\to\Omega^{k}(\ST_i)\big)} \, .
\]
If \(\omega\in\Omega^k(\ST_i)\) is closed (\(d\omega=0\)), its cohomology class \([\omega]\in H^k_{\mathrm{dR}}(\ST_i)\) is the coset
\[
[\omega]=\{\omega + d\nu:\nu\in\Omega^{k-1}(\ST_i)\} \, ,
\]
so two closed \(k\)-forms represent the same class precisely when their difference is exact. Similarly, one defines the \(k\)-th cohomology \(H^k_{\mathrm{dR}} (\partial \ST_i) \). 


The pullback \(\mathcal{J}^*:\Omega^k(\ST_i)\to\Omega^k(\partial\ST_i)\) of an inclusion map \(\mathcal{J}:\partial\ST_i\hookrightarrow\ST_i\) commutes with the exterior derivative. Therefore, the kernel of \(\mathcal{J}^*\), denoted \(\Omega^k_D(\ST_i) \subset \Omega^k(\ST_i)\), forms a sequence 
\[
\dots, \, \Omega^0_D(\ST_i), \, \Omega^1_D(\ST_i), \, \dots
\] 
of abelian groups with connecting homomorphisms \(d\). This is referred to as the \emph{de Rham complex} of \(\ST_i\) \emph{relative} to \(\partial \ST_i\). The corresponding \(k\)-th \emph{de Rham cohomology} of \(\ST_i\) \emph{relative} to \(\partial \ST_i\) is
\[
H^k_{\mathrm{dR}}(\ST_i,\partial\ST_i)
\coloneqq
\frac{\ker\big(d:\Omega^k_D(\ST_i)\to\Omega^{k+1}_D(\ST_i)\big)}
     {\operatorname{im}\big(d:\Omega^{k-1}_D(\ST_i)\to\Omega^{k}_D(\ST_i)\big)} \, .
\]
Relative de Rham cohomology classes are represented by closed \(k\)-forms that vanish under \(\mathcal{J}^*\) modulo exact \(k\)-forms with primitives vanishing under \(\mathcal{J}^*\).




For each integer \(k\ge 0\) and points \(v_0,\dots,v_k \in \R^k\), the \emph{geometric \(k\)-simplex} spanned by \(v_0,\dots,v_k\) is the set,
\[
[v_0,\dots,v_k]
\coloneqq
\Bigg\{\sum_{j=0}^k t_j v_j \, :\; t_j\in[0,1]\text{ for all }j,\ \sum_{j=0}^k t_j=1\Bigg\} \, .
\]
If the vertices \(v_0,\dots,v_k\) are affinely independent this set is a \(k\)-dimensional simplex, otherwise it lies in a lower-dimensional affine subspace. Let \(e_0 \coloneqq 0\in\mathbb{R}^k\) and, for \(i=1,\dots,k\), let \(e_i\) be the \(i\)-th standard basis vector of \(\mathbb{R}^k\). The \emph{standard \(k\)-simplex} is \([e_0,e_1,\dots,e_k] \subset \R^k\). 



A \emph{singular \(k\)-simplex} in \(\ST_i\) is a continuous map
\[
    \sigma: [e_0,e_1,\dots,e_k] \longrightarrow \ST_i \, .
\]
The \(k\)th \emph{singular chain group} of \(\ST_i\) with coefficients in the real numbers is the real vector space of finite, formal linear combinations of singular \(k\)-simplices:
\[
C_k(\ST_i)
\coloneqq
\Bigg\{\sum_{j=1}^N t_j \sigma_j \, : \; 0 < N \in\mathbb{N},\   \ t_j\in\mathbb{R} \text{ and } \sigma_j \text{ is a singular \(k\)-simplex for all } j \Bigg\} \, .
\]
Elements of \(C_k(\ST_i)\) are called \emph{singular \(k\)-chains}. For \(k < 0\), we set \(C_k(\ST_i) = \{0\}\).

Consider a boundary operator \( \partial\) that acts on any singular \(k\)-simplex \(\sigma\) and returns a singular \((k-1)\)-simplex via
\[
 \partial \sigma
=\sum_{j=0}^k (-1)^j\, \sigma \circ \mathfrak{d}_j \, ,
\]
where, in terms of any vector \((t_0 , \dots, t_{k-1}) \in [e_0, \dots, e_{k-1}]\)
\[
    \mathfrak{d}_j : (t_0 , \dots, t_{k-1}) \mapsto (t_0 , \dots, t_{j-1}, 0, t_j , \dots,  t_{k-1}) \in  [e_0, \dots, e_{k}] \, .
\]
This boundary operator is uniquely extended to a boundary operator on singular \(k\)-chains via 
\[
\partial : C_k(\ST_i)\longrightarrow C_{k-1}(\ST_i) \, , \quad \text{via} \quad \partial_k w = \partial_k \sum_{j=1}^N t_j \sigma_j =  \sum_{i=1}^N t_j \partial_k \sigma_j \, .
\]
for a singular \(k\)-chain \(w \) where \(t_j \in \R\) and \(\sigma_j\) is a singular \(k\)-simplex for all \(j\). We will identify the domain for the boundary map \(\partial\) whenever it is unclear from context.





The operators \(\partial\) acting on \(k\)-chains are homomorphisms satisfying \(\partial\circ\partial=0\) \cite{lee}, hence a sequence \(\dots , \, C_1(\ST_i) \mathbin{,} \, C_0(\ST_i) \mathbin{,} \, \dots \) with connecting boundary operators is a chain complex called the \emph{singular chain complex} of \(\ST_i\). The \(k\)th \emph{singular homology} group of \(\ST_i\) with real coefficients is the quotient vector space 
\[
  H_k(\ST_i)\coloneqq  \frac{\ker(\partial: C_k(\ST_i)\longrightarrow C_{k-1}(\ST_i)) }{  \operatorname{im}(\partial: C_{k+1}(\ST_i)\longrightarrow C_{k}(\ST_i))} \, .
\]
Similarly, one defines a homology group for the boundary \(H_k(\partial \ST_i)\). An inclusion \(\mathcal{J}:\partial \ST_i \hookrightarrow \ST_i\) induces an inclusion of chain complexes \( \mathcal{J}_* : C_k(\partial \ST_i )\hookrightarrow C_k(\ST_i)\) via a pushforward \(\mathcal{J}_*\). The quotient
\[
  C_k(\ST_i,\partial \ST_i )\coloneqq C_k(\ST_i)/C_k(\partial \ST_i ) \, ,
\]
is therefore a vector space for each \(k\), and the boundary map \(\partial\) acting on \(C_k(\ST_i)\) descends to a well-defined linear map on the quotient (page 115 \cite{hatcher2005algebraic})
\[
   \partial: C_k(\ST_i,\partial \ST_i )\longrightarrow C_{k-1}(\ST_i,\partial \ST_i ) \, .
\]
So \(\dots , \, C_1(\ST_i, \partial \ST_i ) , \, C_0(\ST_i, \partial \ST_i ), \, \dots \) with the boundary operations is a chain complex. The \(k\)th homology group of \(\ST_i\) relative to \(\partial \ST_i\) is 
\[
  H_k(\ST_i,\partial \ST_i)\coloneqq  \frac{\ker( \partial: C_k(\ST_i, \partial \ST_i )\longrightarrow C_{k-1}(\ST_i, \partial \ST_i)) }{  \operatorname{im}(  \partial: C_{k+1}(\ST_i, \partial \ST_i)\longrightarrow C_{k}(\ST_i, \partial \ST_i))} \, .
\]




\subsection{Proof of Lemma \ref{Lemma:finite}}\label{sub:finite}


Recall Lemma \ref{Lemma:finite}:
\begin{center}
    \textit{The homologies \(H_k (\ST_i)\) and \(H_k (\overline{\ST}_i^c)\) are finite-dimensional vector spaces for all \(k,i\).}
\end{center}

 
These homologies are vector spaces over the reals by construction, so we prove the finite dimensionality. We follow a discussion given by Bott and Tu \cite{bott2013differential}.


\begin{proof}
   


A finite open cover \(\mathcal{U}=\{U_1,\dots,U_N\}\) of \(\ST_i\) is called a \emph{good cover} if every non-empty finite intersection of the sets \(U_j\) is diffeomorphic to \(\mathbb{R}^3\). We now show that every Riemannian \(3\)-manifold admits a finite good cover.

Fix a point \(x\in \operatorname{int}(\ST_i)\). By the existence and uniqueness theorem for ordinary differential equations, for each \(v\in T_x\ST_i\) there exists \(\varepsilon>0\) and a unique geodesic
\[
    \gamma_v:(-\varepsilon,\varepsilon)\to \ST_i
\]
such that \(\gamma_v(0)=x\) and \(\gamma_v'(0)=v\). The exponential map at \(x\) is therefore defined on a neighbourhood \(U\) of the origin in \(T_x\ST_i\) such that for every \(v \in U\), the geodesic \(\gamma_v\) is also defined on \([0,1]\). We then define the exponential map by
\[
    \exp_x : U \subseteq  T_x\ST_i \longrightarrow \ST_i \, , 
    \qquad 
    \exp_x(v)\coloneqq \gamma_v(1) \, .
\]
By the homogeneity of geodesics, \(\exp_x(sv)=\gamma_v(s)\) for \(s \in [0,1]\). Differentiating at \(s=0\) gives
\[
    d\exp_x(v)
    =
    \frac{d}{ds}\Big|_{s=0}\exp_x(sv)
    =
    \gamma_v'(0)
    =
    v \, ,
\]
and hence \(d(\exp_x)\) is the identity on \(U \subseteq T_x \ST_i\). The inverse function theorem implies that there exists a ball \(\B(r) \subset  T_x\ST_i\) with radius \(r>0\) about \(0\) such that 
\[
    \exp_x : U = \B(r)\longrightarrow V_x\coloneqq \exp_x\bigl(\B(r)\bigr)
\]
is a diffeomorphism onto an open neighbourhood \(V_x\) of \(x\) in \(\ST_i\).

By choosing \(r\) sufficiently small \(V_x\) is geodesically convex \cite{lee}. That is, for any two points \(y,z\in V_x\), there exists a unique geodesic segment joining \(y\) to \(z\), and this segment lies entirely in \(V_x\).

A collection \(\{V_x : x\in \ST_i\}\) is therefore an open cover of \(\ST_i\). Since \(\ST_i\) is compact, this cover admits a finite subcover, namely, there exist points \(x_1,\dots,x_N\in \ST_i\) such that
\[
    \ST_i = V_{x_1}\cup\cdots\cup V_{x_N} \, .
\]
Moreover, each \(V_{x_j}\) is geodesically convex.








Consider the geodesically convex neighbourhoods \(V_{x_1},V_{x_2}\) that are small enough so that the geodesic between any two points within the neighbourhoods is unique, as in the above discussion. Then in \(V_{x_1} \cap V_{x_2}\), the geodesic is unique in both sets separately, and contained in both sets separately, so \(V_{x_1} \cap V_{x_2}\) is geodesically convex. 

As we stated, the exponential map is a diffeomorphism on the neighbourhoods of any interior point. It therefore provides a diffeomorphism to the tangent space \(T_x \ST_i \cong \R^3\) for all \(x \in \text{int} (\ST_i)\). Therefore any compact Riemannian 3-manifold has a finite good cover.





Recall the Poincaré Lemma: \(H^k_{\mathrm{dR}}(\mathbb{R}^3)\cong 0 \) for \(k>0\), and \(H^0_{\mathrm{dR}}(\mathbb{R}^3)\cong\mathbb{R}\). We prove by induction on the number of geodesically convex sets in a union that the de Rham cohomology of such a union is finite dimensional. Applying this to a finite cover of \(\ST_i\) will gives the desired statement.

For a single geodesically convex neighbourhood \(V_{x_1}\) the Poincaré Lemma (and the fact that \(V_{x_1}\) is diffeomorphic to an open ball in \(\mathbb{R}^3\)) implies \(H^k_{\mathrm{dR}}(V_{x_1})\) is finite dimensional.

Now suppose, for the sake of induction, that \(H^k_{\mathrm{dR}}(V)\) is finite dimensional where
\[
V=V_{x_1}\cup\cdots\cup V_{x_j} \, .
\]
The Mayer–Vietoris long exact sequence for de Rham cohomology gives a segment
\[
\cdots \longrightarrow
H^{k-1}_{\mathrm{dR}}(V\cap V_{x_{j+1}})
\xrightarrow{\;d^*\;}
H^{k}_{\mathrm{dR}}(V\cup V_{x_{j+1}})
\xrightarrow{\;\mathcal{J}_*\;}
H^{k}_{\mathrm{dR}}(V)\oplus H^{k}_{\mathrm{dR}}(V_{x_{j+1}})
\longrightarrow\cdots \, .
\]
Exactness at \(H^{k}_{\mathrm{dR}}(V\cup V_{x_{j+1}})\) yields
\[
\ker(\mathcal{J}_*)=\operatorname{Im}(d^*) \, ,
\]
hence, a short exact sequence of vector spaces is
\[
0 \longrightarrow \operatorname{Im}(d^*) \longrightarrow
H^{k}_{\mathrm{dR}}(V\cup V_{x_{j+1}}) \longrightarrow
\operatorname{Im}(\mathcal{J}_*) \longrightarrow 0 \, .
\]
Consequently
\[
\dim H^{k}_{\mathrm{dR}}(V\cup V_{x_{j+1}})
= \dim\operatorname{Im}(d^*) + \dim\operatorname{Im}(\mathcal{J}_*) \, .
\]
But \(\operatorname{Im}(d^*)\) is a quotient (hence a homomorphic image) of \(H^{k-1}_{\mathrm{dR}}(V\cap V_{x_{j+1}})\), so it is finite dimensional whenever \(H^{k-1}_{\mathrm{dR}}(V\cap V_{x_{j+1}})\) is finite dimensional. Also \(\operatorname{Im}(\mathcal{J}_*)\) is a subspace of the finite-dimensional space \(H^{k}_{\mathrm{dR}}(V)\oplus H^{k}_{\mathrm{dR}}(V_{x_{j+1}})\), so it is finite dimensional. Therefore, if
\[
H^{k-1}_{\mathrm{dR}}(V\cap V_{x_{j+1}}) \, ,\qquad
H^{k}_{\mathrm{dR}}(V) \, ,\qquad
H^{k}_{\mathrm{dR}}(V_{x_{j+1}}) \, , 
\]
are finite dimensional, then \(H^{k}_{\mathrm{dR}}(V\cup V_{x_{j+1}})\) is finite dimensional. We still need to verify that \(H^{k-1}_{\mathrm{dR}}(V\cap V_{x_{j+1}})\) is finite. Note that
\[
V\cap V_{x_{j+1}}
=\bigcup_{i=1}^j \bigl(V_{x_i}\cap V_{x_{j+1}}\bigr) \, .
\]
Each intersection \(V_{x_i}\cap V_{x_{j+1}}\) is geodesically convex and any non-empty finite intersection of these is again a good cover of \(V\cap V_{x_{j+1}}\). Hence, \(H^{k-1}_{\mathrm{dR}}(V\cap V_{x_{j+1}})\) is finite dimensional. Therefore, \(\operatorname{Im}(d^*)\) is finite dimensional and the argument above shows \(H^{k}_{\mathrm{dR}}(V\cup V_{x_{j+1}})\) is finite dimensional.

By induction, we conclude that a finite union
\[
\ST_i = V_{x_1}\cup\cdots\cup V_{x_N}
\]
produces a finite dimensional vector space \(H^k_{\mathrm{dR}}(\ST_i)\).

Using the de Rham theorem we see that \(H^k_{\mathrm{dR}} (\ST)\) has the same dimension as \( H_{k} (\ST) \), so \(H_{k} (\ST)\) is finite dimensional. The same argument applies to \(\overline{\ST}_i^c\).
\end{proof}




\subsection{Homology of Complements}\label{sub:complement}



Consider a result by Cantarella \cite{CANTARELLA20101127}: Let \(\mathbb{S}^k\) and \(\mathbb{D}^k\) be the usual \(k\)-dimensional sphere and disk. Thinking of \(\R^3\) topologically as \(\mathbb{S}^3 \setminus \{x\}\) for \(x \in \overline{\ST}_i^c\) we have
\begin{align*}
    H_1( \R^3 \setminus \ST_i ) \cong H_1( \mathbb{S}^3 \setminus \ST_i ) \, . 
\end{align*}
\begin{proof}
    Let \(\mathbb{D}^3\) be a disk about \(x\) and not intersecting \(\ST_i\). The union and intersection of \(\mathbb{D}^3\) and \(\R^3 \setminus \ST_i \cong (\mathbb{S}^3 \setminus\{x\}) \setminus \ST_i\) is
    \begin{align*}
        \mathbb{D}^3 \cup ( (\mathbb{S}^3 \setminus \{x\}) \setminus \ST_i ) = \mathbb{S}^3 \setminus \ST_i \, ,
        \qquad 
        \text{and}
        \qquad 
        \mathbb{D}^3 \cap ( (\mathbb{S}^3 \setminus \{x\}) \setminus \ST_i ) = \mathbb{D}^3  \setminus \{x\}  \cong \mathbb{S}^2 \, .
    \end{align*}
    This provides an exact short sequence:
    \begin{align*}
        H_1(\mathbb{S}^2) \to H_1(\mathbb{D}^3) \oplus H_1( \R^3 \setminus \ST_i )  \to H_1( \mathbb{S}^3 \setminus \ST_i ) \to H_0(\mathbb{S}^2) \, .
    \end{align*}
    As \(\mathbb{D}^3\) is contractible and the first and last homology groups in the exact sequence are zero implying the result. 
\end{proof}



\section{Alexander Basis Exposition}\label{section:Alexander Basis}




This section provides some background on the construction of an Alexander basis. We present a general formulation in subsection \ref{append:general_formulation}, and provide  an explicit example for a hollow torus in subsection \ref{append:example}. The material here is drawn from work by Cantarella and Parsley (appendix B \cite{CANTARELLA20101127}).





Consider a basis \(\{[s_{i,1}], \dots, [s_{i,{\ell_i}}]\}\) of \(H_1 (\ST_i) \) where \(\ST_i\) is a compact, oriented, connected domain in \(\R^3\) with a smooth boundary \(\partial \ST_i\). The genus of \(\partial \ST_i\) is \(0 <{\ell_i} < \infty\). We keep the notation consistent with the remainder of the document, this makes the partition subscript in \(\ST_i\) and other variables superfluous. 



The goal in subsection \ref{append:general_formulation} is to prove Theorem \ref{thr:Alexander}. Additionally, we consider any orientation preserving homeomorphism \(f : \ST_i \to \ST_i'\). We show that pushing forward the elements in an Alexander basis for \(H_1(\partial \ST_i)\) gives an Alexander basis for \(H_1(\partial \ST_i')\). This is used throughout our paper in order to have a consistent definition of flux, relative helicity, and Amperian gauge on domains \(f(\ST_i)\) with \(f \in \mathcal{I}( \ST_i )\).


First we provide and prove two well known results.
\begin{lemma}[Proposition C.4 \cite{CANTARELLA20101127}]\label{Lemma:Potential}
    Any closed, normal \(2\)-form on \(\ST_i\) is exact.
\end{lemma}

\begin{proof}
    Let \(b_i \in \Omega^2_D( \ST_i )\) be a closed, normal differential \(2\)-form. Applying a Hodge-Morrey-Friedrichs (HMF) decomposition to the 1-form \(\star b\), then applying a \(\star\)
    \begin{align*}
        b = d \eta_1 + \delta \eta_2 + d \eta_3 + \eta_4^* \, , 
    \end{align*}
    where \(\eta_1 \in \Omega^1 _D (\ST_i)\), \(\eta_2 \in  \Omega^3_N (\ST_i)\), \( d \eta_3 \in \mathcal{H}^2 (\ST_i) \), where \(\eta_3 \in \Omega^1  (\ST_i)\), and \(\eta_4^* \in \mathcal{H}_N^2 (\ST_i) \). Following a similar approach as in section \ref{sub:Hodge}, we apply the closed, and normal conditions:
    \begin{align*}
        0 = db = d   \delta \eta_2  \, , \qquad 0 = tb =  t (  \delta \eta_2 + d \eta_3 + \eta_4^* ) \, .
    \end{align*}
    Taking note of the first equality, and using the divergence theorem:
    \begin{align*}
        \|  \delta \eta_2 \|_{L^2}^2 = \int_{\ST_i}  \delta \eta_2 \wedge  d \star \eta_2 = -\int_{\ST_i} d \delta \eta_2 \; \star \eta_2 + \int_{\partial \ST_i} \mathcal{J}^* (\delta \eta_2 \; \star \eta_2) = 0 \, . 
    \end{align*}
    And by continuity we have that \(\delta \eta_2 = 0\). 

    
    The normal condition is now \( 0 =  t ( d \eta_3 + \eta_4^* ) \). Using de Rham's theorem, \(\eta_4^*\) is a unique harmonic representative for \([\eta_4^*] \in H^2_{\mathrm{dR}} (\ST_i) \). This has a unique dual element \([\eta_4] \in H_2 (\ST_i)\). Namely, we take an algebraic dual relative to a given basis (of dimension \(\ell_i\) by the Hodge isomorphism). Computing
    \begin{align*}
        \langle [ \eta_4^* ], [ \eta_4 ] \rangle = \int_{ \eta_4 } \eta_4^* 
        = \int_{ \mathcal{J}_* (\eta_4 ) } \mathcal{J}^* ( \eta_4^* )
        = \int_{\mathcal{J}_* (\eta_4 )} \mathcal{J}^* (d \eta_3 + \eta_4^* )   = 0 \, . 
    \end{align*}
    Here \(\mathcal{J}\) is the usual boundary inclusion map. The only way to achieve this is when \([\eta_4^*]=0\). Noting that \(0 = [\star \eta_4^*] \in H^1_{\mathrm{dR}} (\ST_i, \partial \ST_i) \) has a unique harmonic representative, we have that \(\eta_4^* = 0\). Hence, \(b = d \eta_3\) is exact.
\end{proof}





\subsection{General Formulation}\label{append:general_formulation}



Given any choice of bases for \(H_1(\ST_i)\) and \(H_1 (\overline{\ST}^c_i )\) one can, of course, define a basis for \(H_k (\partial \ST_i)\) (Lemma \ref{Lemma:form_basis}). Theorem \ref{thr:Alexander} asserts that we already have a basis for \(H_1(\ST_i)\) given by \(\{ [s_{i,1}], \dots, [s_{i,{\ell_i}}] \}\), all that remains is to construct a suitable basis for \(H_1(\overline{\ST_i}^c)\) that satisfies the requirements of an Alexander basis. We proceed in two steps: First, by constructing \(\{[t_{i,1}], \dots, [t_{i,\ell_i}]\}\) (a basis for \(H_1 (\overline{\ST}^c_i )\)). Secondly, by verifying this basis satisfies the conditions in Theorem \ref{thr:Alexander}. 




\bigskip 

\textit{Step 1. Construct a Basis for \(H_1(\partial \ST_i)\) via \(\{[s_{i,1}], \dots, [s_{i,{\ell_i}}] \}\).}


As a basis for \(H_1(\ST_i)\) is given. We construct basis elements \(T_{i,j}, \sigma_{i,j}\) and the duals \([s_{i,j}^*]\), \([T_{i,j}^*]\), and \([\sigma_{i,j}^*]\) as in the main body. We will see an explicit example to compute these in the next subsection. 

We want to see how \(H_2(\ST_i , \partial \ST_i)\) relates to \(H_1(\overline{\ST_i}^c)\). Let \((\ST_i)_\epsilon\) be all points in \(\R^3\) within \(\epsilon >0\) of \(\ST_i\) and \(\overline{\ST_i}^c_\epsilon\) is the closure of the complement of \((\ST_i)_\epsilon\). We take \(\epsilon\) small enough that these spaces are compact and oriented in \(\R^3\) with smooth boundary. Note that \(\overline{\ST_i}^c_\epsilon \subset \overline{\ST_i}^c \subset \R^3\), so by the excision theorem, the inclusion \(\mathcal{J}^0 : ( \R^3 - \overline{\ST_i}^c_\epsilon, \overline{\ST_i}^c - \overline{\ST_i}^c_\epsilon) \hookrightarrow (\R^3, \overline{\ST_i}^c)\) induces an isomorphism \(H_k (\R^3 - \overline{\ST_i}^c_\epsilon, \overline{\ST_i}^c - \overline{\ST_i}^c_\epsilon) \cong H_k (\R^3 , \overline{\ST_i}^c) \) (see page 119, theorem 2.20 \cite{hatcher2005algebraic}). Taking \(\epsilon \to 0\) (deformation retraction) we have that \(H_k (\ST_i, \partial \ST_i) \cong H_k (\R^3 , \overline{\ST_i}^c) \).


Taking a relative Mayer-Vietoris sequence for \(k>0\):
\begin{align*}
      0 
    \to  H_{k+1} (\R^3, \overline{\ST_i}^c )
    \xrightarrow{\partial}   H_{k} (\overline{\ST_i}^c ) 
    \to  0 \, .
\end{align*}
Giving isomorphisms \(H_{k+1} (\ST_i, \partial \ST_i) \cong H_{k+1} (\R^3 , \overline{\ST_i}^c) \cong H_{k} (\overline{\ST_i}^c ) \) given by \( \mathcal{J}^0_*\) and \(\partial\). Note that inverting the domains produces the result \(H_{k+1} (\overline{\ST_i}^c, \partial \overline{\ST_i}^c) \cong H_{k} (\ST_i ) \) which we will utilise later. We define (through these isomorphisms) a basis for \(H_{1} (\overline{\ST_i}^c )\) with elements given by \([t_i] \coloneqq \partial \mathcal{J}^0_* [T_{i, j}]  \). 


From the \([t_{i,j}]\) we define \([t_{i,j}^*], [S_{i,j}], [S_{i,j}^*], [\tau_{i,j}], [\tau_{i,j}^*]\) as in the main body. We also see that the basis for \(H_{1} (\overline{\ST_i}^c )\) has dimension \(\ell_i\). 



\bigskip 

\textit{Step 2. Verifying that the basis constructed in step 1 is an Alexander basis for \(H_1(\partial \ST_i)\).}



Theorem \ref{thr:Alexander} claims that for all \( j, k \in \{ 1, \dots, {\ell_i}\} \):
\begin{align*}
    \begin{gathered}
        \int_{\partial \ST_i} \sigma_{i,j}^* \wedge \sigma_{i,k}^* = 
         0 \,,    \qquad 
         \int_{\partial \ST_i} \tau_{i,j}^* \wedge \tau_{i,k}^*
        = 0 \, ,
        \qquad 
         \int_{\partial \ST_i} \tau_{i,j}^* \wedge \sigma_{i,k}^* = \delta_{jk}  \,.
    \end{gathered}
\end{align*}  
The first 2 integrals we can compute directly, namely, via the divergence theorem we have that 
\begin{align*}
    \int_{\partial \ST_i} \sigma_{i,j}^* \wedge \sigma_{i,k}^* = \int_{ \ST_i} d( s_{i,j}^* \wedge s_{i,k}^* ) = \int_{ \ST_i} 0  = 0 \, .
\end{align*}
Similarly,
\begin{align*}
    \int_{\partial \ST_i} \tau_{i,j}^* \wedge \tau_{i,k}^* = - \int_{ \overline{\ST}^c_i} d( t_{i,j}^* \wedge t_{i,k}^* ) = - \int_{ \overline{\ST}^c_i } 0  = 0 \, .
\end{align*}
For the last integral, we note that \([T_{i,j}^*] \in H^2_{\mathrm{dR}} (\ST_i, \partial \ST_i) \). By the de Rham theorem, we have a unique harmonic representative \(T_{i,j}^* \in \mathcal{H}^2_D(\ST_i)\). Importantly, \(T_{i,j}^*\) is a closed, normal 2-form on \(\ST_i\), so by Lemma \ref{Lemma:Potential}, \(T_{i,j}^* = d \Gamma_{i,j} \) for some \(\Gamma_{i,j} \in \Omega^1(\ST)\). We know that the collection of \([\sigma_{i,j}^*], [\tau_{i,j}^*]\) form a basis for \(H^1_{\mathrm{dR}}(\ST_i)\) (Lemma \ref{Lemma:form_basis}), so given \([\mathcal{J}^* \Gamma_{i,j} ] \in H^1_{\mathrm{dR}} (\partial \ST_i)\), 
\begin{align*}
    [\mathcal{J}^* \Gamma_{i,j} ] = \sum_{r=1}^{\ell_i} \Big( k_{j,r} [\tau_{i,r}^*] +   c_{j,r} [\sigma_{i,r}^*] \Big) \, .
\end{align*}
By definition of \([T_{i,j}^*]\),
\begin{align*}
    \delta_{jk}  = \int_{T_{i,j}} T_{i,k}^*
    =   \int_{T_{i,j}} d \Gamma_{i,k} 
    =    \int_{\tau_{i,j}} \mathcal{J}^*  \Gamma_{i,k} 
    =
    \sum_{r=1}^{\ell_i} \left( k_{k,r} \int_{\tau_{i,j}}  \tau_{i,r}^* +  c_{k,r}  \int_{\tau_{i,j}} \sigma_{i,r}^* \right) = k_{k,j}
    \, .
\end{align*}
Recall that section \ref{section:Problem Construction} shows \(\tau_{i,j}^*\), \(\sigma_{i,j}^*\) are the duals of \(\tau_{i,j}\), \(\sigma_{i,j}\) respectively. Substituting this back into the non-degenerate Lefschetz pairing
\begin{align*}
    \delta_{jk} &= \int_{ \ST_i} \sigma_{i,j}^* \wedge d \Gamma_{i,k} 
    = - \int_{\partial \ST_i} \sigma_{i,j}^* \wedge \mathcal{J}^*\Gamma_{i,k}
    = - \int_{\partial \ST_i} \sigma_{i,j}^* \wedge \Big( \tau_{i,k}^* + \sum_{r=1}^{\ell_i} c_{k,r} \sigma_{i,r}^* \Big)
    %\\
    %&= - \int_{\partial \ST_i} \sigma_{i,j}^* \wedge  \tau_{i,k}^*
    %-
    %\sum_{r=1}^{\ell_i} c_{k,r} \int_{\partial \ST_i} \sigma_{i,j}^* \wedge   \sigma_{i,r}^*
    =  \int_{\partial \ST_i}  \tau_{i,k}^*  \wedge \sigma_{i,j}^*
    \, .
\end{align*}
This completes the proof for the integral relations. But it remains to be shown that \(\partial [T_{i,j}] = [\tau_{i,j}]\) and \(\partial [S_{i,j}] = [\sigma_{i,j}]\). 

Consider inclusion maps \(\mathcal{J}^0 : (\ST_i, \partial \ST_i) \xhookrightarrow{} (\R^3 , \overline{\ST}_i^c)\), \(\mathcal{J} :  \partial \ST_i \xhookrightarrow{}  \ST_i\) and \( \mathcal{J}^c : \partial \ST_i \xhookrightarrow{} \overline{ \ST }_i^c \). Recalling that \((\mathcal{J}_*, \mathcal{J}^c_*) : H_1(\partial \ST_i) \to H_1(\ST_i) \oplus H_1(\overline{\ST}^c_i)\) is an isomorphism (Lemma \ref{Lemma:isomorphism_inclusion}), for any \([y] \in H_2(\ST_i, \partial \ST_i)\):
\begin{align*}
    (\mathcal{J}_*, \mathcal{J}^c_*) \partial [y] = ( \mathcal{J}_*  \partial [y], \mathcal{J}^c_*  \partial [y]) 
    =  ( \mathcal{J}_*  \partial [y],   \partial \mathcal{J}^0_* [y])  \, . 
\end{align*}
Here the boundary acts as \(\partial : H_2(\ST_i, \partial \ST_i) \to H_{1}(\overline{\ST}_i^c)\) except for the last boundary operator, acting on spaces \(\partial : H_2(\R^3 , \overline{\ST}_i^c) \to H_{1}(\overline{\ST}_i^c)\). Recall a portion of the long exact sequence of the pair \((\ST_i,\partial\ST_i)\):
\[
\cdots \longrightarrow H_2(\ST_i) \longrightarrow H_2(\ST_i,\partial\ST_i)
\xrightarrow{\; \partial \;} H_1(\partial\ST_i)
\xrightarrow{\;\mathcal{J}_*\;} H_1(\ST_i)\longrightarrow\cdots .
\]
Exactness implies \(\operatorname{im}(\partial ) = \ker( \mathcal{J}_* )\). Hence \(\mathcal{J}_*(\partial [y])=0\),
showing that \((\mathcal{J}_*, \mathcal{J}^c_*)  \partial [y]  =  ( 0,   \partial \mathcal{J}^0_* [y])\). As \((\mathcal{J}_*, \mathcal{J}^c_*) \) is an isomorphism, it has an inverse, so, setting \([y] = [T_{i,j}]\) for any \(i,j\),
\begin{align*}
    \partial [T_{i,j}]  = (\mathcal{J}_*, \mathcal{J}^c_*)^{-1} ( 0, [t_{i,j}]) = [\tau_{i,j}] \, ,
\end{align*}
as required. 


Now for the last claim, that \([\sigma_{i,j}] = \partial [S_{i,j}]\). We just showed \(\mathcal{J}_*(\partial [y])=0\) for all \([y] \in H_2( \ST_i, \partial \ST_i) \). By exactly the same argument for all \([z] \in H_2 ( \overline{\ST}_i^c ,  \partial \overline{\ST}_i^c )\) we have \(\mathcal{J}^c_*(\partial [z])=0\). Therefore 
\[\partial [S_{i,j}] = \sum_{k=1}^{\ell_i} c_{i,j,k} [\sigma_{i,k}] \, . \]
Select some \(\eta \in \Omega^1(\overline{\ST}^c_i)\) such that \(\eta |_{\partial \ST_i} = \sigma_{i,k}^*\). On the boundary we see that \(d\eta |_{\partial \ST_i} = d \sigma_{i,k}^* = 0\), so \([d\eta] \in H_2(\overline{\ST}^c_i, \partial \overline{\ST}^c_i ) \). Therefore we have, using Lefschetz duality and the divergence theorem 
\begin{align*}
    \int_{\partial S_{i,j}} \sigma_{i,k}^* =  \int_{S_{i,j}}  d\eta = \int_{\overline{\ST}^c_i} t_{i,j}^* \wedge d\eta  = - \int_{\overline{\ST}^c_i} d (t_{i,j}^* \wedge \eta)  
    =    \int_{\partial \ST_i} \tau_{i,j}^* \wedge \sigma_{i,k}^* = \delta_{jk} \, . 
\end{align*}
By selecting an outward orientation for \(\partial \ST_i\) the divergence theorem provides a negative sign when moving to the boundary of \(\overline{\ST}^c_i\). This gives the relation when placing in basis elements
\begin{align*}
    \delta_{jk} = \int_{\partial S_{i,j} }  \sigma_{i,k}^*
    = \int_{ \sum_{r=1}^{\ell_i}  c_{i,j,r} \sigma_{i,r} }  \sigma_{i,k}^*
    = \sum_{r=1}^{\ell_i}  c_{i,j,r}  \int_{ \sigma_{i,r} }  \sigma_{i,k}^*
    =   c_{i,j,k} \, .
\end{align*}
Hence \( \partial [S_{i,j}] =  [\sigma_{i,j}]  \). This completes the proof of Theorem \ref{thr:Alexander}. \qed



Let \(f:\widehat{\ST}_i \to \ST_i\) be an orientation-preserving homeomorphism. Since \(f\) is a homeomorphism the induced maps
\[
(f^{-1})^*:H_k(\widehat{\ST}_i) \to H_k(\ST_i), \qquad
f^*:H^k_{\mathrm{dR}}(\ST_i) \to H^k_{\mathrm{dR}}(\widehat{\ST}_i)
\]
are isomorphisms (with inverses \((f^{-1})_*\) and \((f^{-1})^*\) respectively). Let \(f_\partial:=f\big|_{\partial\widehat{\ST}_i}:\partial\widehat{\ST}_i\to\partial\ST_i\) be the corresponding map on the boundaries. 

Let \(\mathcal{J}:\partial \ST_i \hookrightarrow \ST_i\) denote the boundary inclusion and write \(\widehat{\mathcal{J}}:\partial\widehat{\ST}_i\hookrightarrow\widehat{\ST}_i\) for the analogous inclusion on the hatted side. Recall that an Alexander basis
\(\{[\sigma_{i,1}],\dots,[\sigma_{i,\ell_i}],[\tau_{i,1}],\dots,[\tau_{i,\ell_i}]\}\subset H_1(\partial\ST_i)\) has the property that the inclusion-induced map \(\mathcal{J}_*:H_1(\partial\ST_i)\to H_1(\ST_i)\) sends \(\,[\sigma_{i,j}]\mapsto[s_{i,j}]\) and \(\mathcal{J}^c :\partial \ST_i \hookrightarrow \overline{\ST}^c_i\) induces a map \(\,[\tau_{i,j}]\mapsto[t_{i,j}]\).

Assume \(\{[\widehat{s}_{i,1}],\dots,[\widehat{s}_{i,\ell_i}]\}\) is a basis of \(H_1(\widehat{\ST}_i)\), where \((f^{-1})^*[\widehat{s}_{i,j}] = s_{i,j}\) is the corresponding basis element of \(H_1(\ST_i)\). Defining \([\widehat \sigma_{i, j}]\) as \([\widehat s_{i, j}]\) pushed forward by inclusion to the boundary we have the relation \( [\widehat \sigma_{i, j}] =  f_\partial^* [\sigma_{i, j}]\) on the boundary. Therefore, one may construct an Alexander basis on \(\partial\widehat{\ST}_i\) using the \(\widehat{s}_{i,j}\). Denote the resulting basis \(\{[\widehat{\sigma}_{i,1}],\dots,[\widehat{\sigma}_{i,\ell_i}], [\widehat{\tau}_{i,1}],\dots,[\widehat{\tau}_{i,\ell_i}]\}\). Similarly, define \(\widehat{S}_{i,j}, \widehat{T}_{i,j} \) and the corresponding duals in the same way as their non-hatted counterparts. So
\begin{align*}
    \int_{T_{i, j}} (f^{-1})^*( \widehat T_{i, k}^* )  = \int_{\ST} s_{i, j}^* \wedge (f^{-1})^*( \widehat T_{i, k}^* ) = \int_{\ST} (f^{-1})^*( \widehat s_{i, j}^* \wedge  \widehat T_{i, k}^* )
    =
    \int_{\widehat \ST_i}  \widehat s_{i, j}^* \wedge \widehat T_{i, k}^* = \delta_{jk}
    =
    \int_{T_{i, j}} T_{i, k}^* 
    \, .
\end{align*}
This implies that \((f^{-1})^*([ \widehat T_{i, j}^*] ) = [T_{i, j}^*]\), so \([ \widehat T_{i, j}^*]  = f^* [T_{i, j}^*]\), and on the boundary this induces \( [ \widehat t_i^*]  =  f_\partial^* [t_i^*]\). Hence, in the current context, the Alexander basis is also pushed forward by \(f\). Namely, \(\{  f_\partial^* [\sigma_{i, 1}] , \dots, f_\partial^* [\sigma_{i, \ell_i}]  ,    f_\partial^* [\tau_{i, 1}] , \dots, f_\partial^* [\tau_{i, \ell_i}]\}\) is an Alexander basis for \(H_1(\partial \widehat \ST) \). This is a helpful when considering an MRxMHD problem where interfaces are pushed forward by \(f \in \I\) because, as long as pushing forward gives a basis for \(H_1(\widehat{\ST}_i)\), the corresponding Alexander basis and related definitions are also pushed forward.

We also see that, only requiring a map \(f^{-1}\) to define a new basis allows us to introduce knots in the domain, we see this in the following example. As \([S_{i,j}] \in H_2 ( \overline{\ST}_i^c , \partial \ST_i )\), we do not know the effect of \(f^{-1}\), but we are aware that some \([S_{i,j}] \) exists on the new domain. 






\begin{figure}
    \centering
    \begin{overpic}[width=0.5\linewidth]{"Chapter_Alexander_Basis_Exposition/torus_clip2.png"}
    \end{overpic}
    \caption{A hollow torus embedded in \(\R^3\) with a trefoil knot.}
    \label{fig:trefoil} 
\end{figure}

\subsection{Hollow Torus Example} \label{append:example}




We apply our construction of an Alexander basis to a hollow torus \(\HT\subset\R^3\). Rather than contract \(\HT\) to a torus, we work with an explicit orientation preserving homeomorphism
\[
\iota:[1,2]\times \mathbb{T}^2 = U \hookrightarrow \HT\subset\R^3 \, ,
\]
where \(\mathbb{T} \coloneqq \R / \mathbb{Z}\) and the coordinate system \((r,\vartheta_P,\vartheta_T)\) with \(r\in[1,2]\) and \(\vartheta_P,\vartheta_T\in[0,1)\). The subscript \(i\) is retained for consistency with the rest of the paper. On the parameter space define singular \(1\)-cycles
\[
s_{i,1}(\omega )=(1,\omega ,0),\qquad s_{i,2}(\omega )=(2,0, \omega ),\qquad \omega \in\Delta_1=[0,1] \, .
\]
Clearly \(\partial s_{i,j}=0\), so \([s_{i,1}]\) and \([s_{i,2}]\) represent classes in \(H_1( U  )\). Since this homology is two-dimensional (the boundary has total genus \(2\)), it suffices to verify linear independence. This follows either from the K\"unneth theorem or from the geometric observation that no singular \(2\)-chain relates \(s_{i,1}\) to \(s_{i,2}\), hence \(\{[s_{i,1}],[s_{i,2}]\}\) is a basis of \(H_1(  U  )\). 

Note that \(s_{i,1}\) and \(s_{i,2}\) are trivially included in the boundary, giving generators for \([\sigma_{i,1}]\) and \([\sigma_{i,2}]\). 

From the previous subsection, so long as \( \iota_* \) pushes elements in \(\{[s_{i,1}],[s_{i,2}]\}\) forward to a basis for \(H_1( \HT)\), the following construction of an Alexander basis on the parameter space is applicable to \(H_1(\partial \HT)\). In the case of an unknotted torus the choice of \(s_{i,j}\) align with those chosen in the construction of a relative helicity given by Pfefferlé, Noakes, and Perrella \cite{Gauge_freedom}. As \(\iota\) is an orientation-preserving homeomorphism, its image \(\HT\) may be a knotted or unknotted torus, a knotted example is Figure \ref{fig:trefoil}. In some knotted cases constructing \(s_{i,j}\) via boundaries of Seifert surfaces may provide geometric intuition (see \cite{seifert1935geschlecht}). 





Our computation of \([T_{i, j}]\) relies on cohomology elements \([s_{i,j}^*]\). We may compute a representative closed 1-form \(s_{i,j}^* \in \Omega^1 (\ST_i) \) through the relation
\begin{align*}
    \delta_{jk}  = \int_{s_{i,j}} s_{i,k}^* \, . 
\end{align*}
Selecting \(s_{i,1}^* = d \vartheta_P \) and \(s_{i,2}^* = d \vartheta_T \) we verify that 
\begin{align*}
    \int_{s_{i,1}} s_{i,1}^* &= \int_{\omega \mapsto  (1,  \omega , 0)} d \vartheta_P  = \int_{0}^1 d  \omega  = 1 \, , \qquad
    \int_{s_{i,1}} s_{i,2}^* = \int_{\omega \mapsto  (1,  \omega , 0)} d \vartheta_T  =  0 \, , \\
    \int_{s_{i,2}} s_{i,2}^* &= \int_{\omega \mapsto  (2, 0,  \omega )} d \vartheta_T  = 1 \, , \qquad \hspace{4em}
    \int_{s_{i,2}} s_{i,1}^* = \int_{\omega \mapsto  (2, 0, \omega )} d \vartheta_P  =0 \, .
\end{align*}
There is some freedom, up to a choice of volume in the parameter space, which we set by selecting a volume form \(\varpi = dr \wedge d \vartheta_P \wedge d \vartheta_T\). Then the corresponding \(T_{i, j}^*\) are defined by
\begin{align*}
\begin{gathered}
        \delta_{jk} =  \int_{ U } s_{i,j}^* \wedge T_{i, k}^* \,  .
\end{gathered}
\end{align*}
As \(T_{i, j}^*\) is a closed, normal 2-form we select 
\begin{align*}
    T_{i, j}^* = f_1(r, \vartheta_P) \,  d r \wedge d \vartheta_P + f_2(r, \vartheta_T) \, d r \wedge d \vartheta_T + f_3(\vartheta_P, \vartheta_T) \, d \vartheta_P \wedge d \vartheta_T \, . 
\end{align*}
Let's first consider \({T_{i,1}}^*\):
\begin{align*}
    0 &= \int_U (d \vartheta_T)  \wedge f_1(r, \vartheta_P) d r \wedge d \vartheta_P = \int_U f_1(r, \vartheta_P)  d r \wedge d \vartheta_P \wedge d \vartheta_T  \, , \\
    1 &= \int_U (d \vartheta_P)  \wedge f_2(r, \vartheta_T) d r \wedge d \vartheta_T = - \int_U f_2(r, \vartheta_T) d r \wedge d \vartheta_P \wedge d \vartheta_T   \, .
\end{align*}
Selecting \(f_2(r, \vartheta_T) = - 1\), \(f_1(r, \vartheta_P) = 0\) and \(f_3(\vartheta_P, \vartheta_T) = 0\) gives us a closed 2-form \(T_{i,1}^* = -  d r \wedge d \vartheta_T \) that satisfies the above conditions. And pulling back to the boundary, \(\mathcal{J}^* T_{i, j}^* = \mathcal{J}^* ( - d r \wedge d \vartheta_T) = 0 \) as \(\mathcal{J}^* (dr)=0\). 

Performing the same for \(T_{i,2}^*\),
\begin{align*}
    1 &= \int_U (d \vartheta_T )  \wedge f_1(r, \vartheta_P) d r \wedge d \vartheta_P = \int_U f_1(r, \vartheta_P)  d r \wedge d \vartheta_P \wedge d \vartheta_T  \, , \\
    0 &= \int_U (d \vartheta_P )  \wedge f_2(r, \vartheta_T) d r \wedge d \vartheta_T = - \int_U f_2(r, \vartheta_T) d r \wedge d \vartheta_P \wedge d \vartheta_T  \, .
\end{align*}
Selecting \(f_1(r, \vartheta_P) =  1\), \(f_2(r, \vartheta_T) = 0\) and \(f_3(\vartheta_P, \vartheta_T) = 0\) gives us a closed 2-form \(T_{i,2}^* = d r \wedge d \vartheta_P  \). Additionally, the boundary condition \(\mathcal{J}^*T_{i,2}^* = 0  \) is satisfied. 

We note that this choice for \(T_{i,1}^*, T_{i,2}^*\) agrees with Pfefferl\'e \emph{et al.} \cite{Gauge_freedom} up to a sign. Clearly \(T_{i, j}^*\) are closed and 0 when pulled back to the boundary, and therefore generators for elements of \(H^2_{\mathrm{dR}}(U, \partial U)\). 


For the singular relative 2-chains \(T_{i, j}\) we require that,
\begin{align*}
    \delta_{jk} =  \int_{T_{i, j}} T_{i, k}^* \, . 
\end{align*}
For \((w_1, w_2) \in  \Delta_2\), let's try the 2-chain:
\begin{align*}
    T_{i,1} (w_1, w_2) = ( 1 + w_1, 0, -  w_2    ) + ( 2 - w_1, 0, -1+w_2    ) \, . 
\end{align*}
The image of \(T_{i, j}\) is a region in an unknotted hollow torus \(\HT\) is given by a toroidal ribbon. 
\begin{align*}
        \int_{T_{i,1}} T_{i,1}^* &= \int_{( 1 + w_1, 0, -     w_2    ) + ( 2 - w_1, 0, -    (1-w_2)    )} - dr \wedge d \vartheta_T 
        %\\
    %&=      \int_0^1 \int_0^{1-w_2}  - d (1 + w_1) \wedge d (-     w_2)
    %+
    % \int_0^1 \int_0^{1-w_2}  - d(2-w_1) \wedge d (-    (1-w_2) )
    % \\
     =    2 \int_0^1 \int_0^{1-w_2}   d w_1 \wedge d  w_2 %\\
     =     1     \, .  \\
        \int_{T_{i,1}} T_{i,2}^* &= \int_{( 1 + w_1, 0, -     w_2    ) + ( 2 - w_1, 0, -    (1-w_2)    )}  dr \wedge d \vartheta_P = 0 \, . %\\
    %&=      \int_0^1 \int_0^{1-w_2}   d (1 + w_1) \wedge d (0)
    %+
    % \int_0^1 \int_0^{1-w_2}   d(2-w_1) \wedge d (0)
    % \\
    % &=   0 \, . 
\end{align*}
Note that the \( T_{i,1} \in H_2( U, \partial U )\) only if \(\partial T_{i,1} \in C_1 (\partial U) \), which we will verify. 





Similarly, we may verify the following choice for \(T_{i,2}\):
\begin{align*}
    T_{i,2} (w_1, w_2) = ( 1 + w_1,       w_2  , 0  ) + ( 2 - w_1,      (1-w_2) , 0   ) \, . 
\end{align*}
To complete our construction of an Alexander basis, we take the boundary of \([T_{i, j}]\):
\begin{align*}
    \partial T_{i, j}  
     : w \mapsto  T_{i, j} ( 1-w, w ) - T_{i, j} ( 0, w ) + T_{i, j} ( w, 0 ) \, ,
\end{align*}
where \(w \in [0,1] =  \Delta_1\). Starting with the representative \(T_{i,1}\):
\begin{align*}
    \partial T_{i,1} : w &\mapsto   T_{i,1} ( 1-w, w ) - T_{i,1} ( 0, w ) + T_{i,1} ( w, 0 ) \\
     &=  ( 2 - w, 0, -     w    ) + ( 1  + w, 0, -    (1-w)    ) \\
     & \; \; \; \; - (( 1 , 0, -     w    ) + ( 2 , 0, -    (1-w)    ) ) \\
     & \; \; \; \; + ( 1 + w, 0, 0   ) + ( 2 - w, 0, - 1       ) \, .
\end{align*}
Noting that \(\vartheta_T\) is periodic, the surviving terms are:
\begin{align*}
    \tau_{i,1} = \partial T_{i,1} : w &\mapsto   - (( 1 , 0, -     w    ) + ( 2 , 0, -    (1-w)    )) = ( 1 , 0,     w    )  + ( 2 , 0, -    w    )\, .
\end{align*}
Similarly for \(T_{i,2}\):
\begin{align*}
    \partial T_{i,2} : w &\mapsto   T_{i,2} ( 1-w, w ) - T_{i,2} ( 0, w ) + T_{i,2} ( w, 0 ) \\
     &=  ( 2 - w,       w ,0   ) + ( 1  + w,      (1-w)  ,0  ) \\
     & \; \; \; \; - (( 1 ,       w  ,0  ) + ( 2 ,      (1-w)  ,0  ) ) \\
     & \; \; \; \; + ( 1 + w,  0  ,0 ) + ( 2 - w,       ,0   ) \, .
\end{align*}
Once again the formal sum of the first 2 and last 2 pairs is zero, and we get
\begin{align*}
    \tau_{i,2} = \partial T_{i,2} : w &\mapsto   ( 1 ,   -     w  ,0  ) + ( 2 ,      w  ,0  )        \, .
\end{align*}
Together the classes for \(\sigma_{i,j}\) and \(\tau_{i,j}\) form an Alexander basis for the parameter space of a hollow torus. We observe that our definition of relative helicity in subsection \ref{subsection:Relative Helicity} is given in this context by Pfefferl\'e \emph{et al.} \cite{Gauge_freedom}.
 
 
 




\subsection{Hodge Decomposition}\label{append:Hodge}

Throughout the paper we utilise both a Hodge decomposition, and a Hodge-Morrey-Friedrichs decomposition. Here we provide these decompositions without proof. A Hodge-Morrey decomposition for the space \( \Omega^k  (\ST_i)\) for \(k \geq 1\) is given by:
\begin{align*}
     \Omega^k  (\ST_i) = &  \{      d\omega \in  \Omega^k (\ST_i)  \; : \, t \omega = 0     \} %\\  
     \oplus% &
    \{      \delta \omega \in  \Omega^k (\ST_i)    \; : \, n \omega = 0       \} %\\  
    \oplus 
    %&
     \{      \omega \in \Omega^k (\ST_i)   \; : \, d \omega = \delta \omega = 0       \} \, .
\end{align*}
These spaces are \(L^2\)-orthogonal to each other, giving a unique algebraic decomposition for any \(\omega \in  \Omega^k  (\ST_i)\) (Lemma 2.4.3 \cite{schwarz-1995}):
\begin{align*}
    \omega = d \phi + \delta (\star \eta) + \gamma  \, , %& &  \tilde{s} = d \tilde{\phi} + \tilde{\gamma} \, ,
\end{align*}
where \( \eta \in  \Omega_D^{3-k-1} (\ST_i) \), \(\phi \in  \Omega^0 (\ST_i)\) where \(\phi|_{\partial \ST_i} = 0\), and \(\gamma \in \mathcal{H}^k( \ST_i)\). This is further decomposed to give a Hodge-Morrey-Friedrichs decomposition (Corollary 2.4.9 \cite{schwarz-1995}):
\begin{align*}
     \Omega^k  (\ST_i) = &  \{      d\omega \in  \Omega^k (\ST_i)  \; : \, t \omega = 0   \} \\  \oplus &
    \{      \delta \omega \in  \Omega^k  (\ST_i)    \; : \, n \omega = 0   \} \\  \oplus &
    \{      \omega \in  \Omega^k  (\ST_i)   \; : \, d \omega = \delta \omega = t \omega= 0       \}\\ \oplus&
    \{      d \omega \in  \Omega^k  (\ST_i)   \; : \, \delta d \omega  = 0        \} \, .
\end{align*}
These spaces are \(L^2\)-orthogonal to each other and give a unique algebraic decomposition for any \(\omega \in  \Omega^k  (\ST_i)\):
\begin{align*}
    \omega = d \phi + \delta (\star \eta) + \star \lambda + dz  \, , 
\end{align*}
where \( \eta \in  \Omega_D^{3-k-1} (\ST_i) \), \(\phi \in  \Omega^0(\ST_i)\), \(\star \lambda \in \mathcal{H}^k_D (\ST_i)\), and \(z \in  \Omega^0 (\ST_i)\) for \(\phi|_{\partial \ST_i} = 0\), and \(z\) is a harmonic function \(\Delta z = 0\) where \(\Delta \) is the Laplacian.



\subsection{Fréchet Differentiable}\label{app:differentiable}


We claim in the main body that certain Fréchet derivatives exist. As these computations are relatively standard, we show this for the relative helicity. The magnetic energy and flux are similar computations.



Consider \( \eta_i , \eta_i' \in \Omega^1 _D(\ST_i)\), we look for the Fr\'echet derivative of \(\mathscr{H} \) with respect to the first coordinate \(\Omega^1 _D(\ST_i)\) equipped with the \(C^\infty\) topology, and the codomain is given the absolute value norm. As in the main body, let
\begin{align*}
    a_i =  \eta_i +  \sum_{j=1}^{{\ell_i}}  \psi_{i,j} \Gamma_{i,j} \, , 
\end{align*}
and \(\beta_i = da_i\). So,
\begin{align*}
    \mathscr{H}( d\eta_i + d\eta_i', \ST_i ) = g_i( \eta_i + \eta_i') &= \int_{\ST_i} \Big( \eta_i + \eta_i' + \sum_{j=1}^{{\ell_i}}  \psi_{i,j} \Gamma_{i,j}  \Big) \wedge d\Big( \eta_i + \eta_i' + \sum_{j=1}^{{\ell_i}}  \psi_{i,j} \Gamma_{i,j}  \Big) \\ & \hspace{12em} - \sum_{j=1}^{{\ell_i}} \psi_{i,j} \int_{\sigma_{i,j}} \Big( \eta_i + \eta_i' + \sum_{j=1}^{{\ell_i}}  \psi_{i,j} \Gamma_{i,j}  \Big) \\
    &= g_i( \eta_i ) +
    2 \int_{\ST_i} \beta_i \wedge  \eta_i'   - \sum_{j=1}^{{\ell_i}} \psi_{i,j} \int_{\sigma_{i,j}}  \eta_i' 
    +
    \int_{\ST_i}  \eta_i'  \wedge d \eta_i' 
\end{align*}
We see that 
\begin{align*}
    \partial_{\eta_i} g_i(  \eta_i'  ) \coloneqq 2 \int_{\ST_i} \beta_i \wedge  \eta_i'   - \sum_{j=1}^{{\ell_i}} \psi_{i,j} \int_{\sigma_{i,j}}  \eta_i' 
\end{align*}
is a linear, continuous operator in \(\eta_i'\). For \(\partial_{\eta_i} g_i(  \eta_i'  )\) to be a Fr\'echet derivative requires
\begin{align*}
    \Big| \int_{\ST_i}  \eta_i'  \wedge d \eta_i' \Big|  / \| \eta_i' \|_{C^\infty} \to 0 \qquad \text{as} \qquad  \| \eta_i' \|_{C^\infty} \to 0 \, . 
\end{align*}
And we see that
\begin{align*}
    \Big| \int_{\ST_i}  \eta_i'  \wedge d \eta_i' \Big| = \langle \eta_i'  , \star d \eta_i' \rangle
    \leq
    \| \eta_i' \|_{L^2} \, \| \star d \eta_i'  \|_{L^2}
    =
    \| \eta_i' \|_{L^2} \, \|  d \eta_i'  \|_{L^2} \, .
\end{align*}
By definition of the \(C^\infty\) topology we see that \(\| \eta_i' \|_{L^2}\) is controlled by the \(C^0\) norm and the volume of \(\ST_i\), giving \( \| \eta_i' \|_{L^2} \leq C \| \eta_i' \|_{C^\infty}\) for some \(C\). Similarly, \(\|  d \eta_i'  \|_{L^2} \leq \|  d \eta_i'  \|_{C^0} | \ST_i |^{1/2} \leq C \|  \eta_i'  \|_{C^1} \leq C \|  \eta_i'  \|_{C^\infty} \) for some constant \(C\) at each step \cite{warner1983foundations}. Therefore, we see that 
\begin{align*}
    0 < \Big| \int_{\ST_i}  \eta_i'  \wedge d \eta_i' \Big|  / \| \eta_i' \|_{C^\infty} \leq \| \eta_i' \|_{L^2} \, \|  d \eta_i'  \|_{L^2}  / \| \eta_i' \|_{C^\infty}  \leq C \| \eta_i' \|_{C^\infty} \, .
\end{align*}
the central terms evidently go to 0 as \(\| \eta_i' \|_{C^\infty} \to 0 \). Hence, \(\partial_{\eta_i} g_i(  \eta_i'  )\) defines a Fr\'echet derivative. 





\subsection{Complementation of the Kernel}\label{app:Complemented}

Consider \(g_i(\eta_i)\) given by the relative helicity. Gl\"ockner's regular value result \ref{theorem:Glock} requires the derivative $d_{\eta_i} g_i$ to have a \emph{complemented} kernel at $\eta_i \in \Omega_D^1(\ST_i)$ \cite{glockner2015fundamentals}. We show that this condition is satisfied whenever $g_i$ has a one-dimensional codomain.


Recall that a closed subspace $K$ of a topological vector space $E$ is said to be \emph{complemented} if there exists a closed subspace $M \subset E$ such that the map 
$K \times M \to E$ given by addition is a topological isomorphism. This is equivalent to saying $E = K \oplus M$.

In our context, we consider the surjective derivative (for \(\beta_i \neq 0\))
\[ L := d_{\eta_i} g_i : T_{\eta_i} \Omega_D^1(\ST_i) \to \mathbb{R} \, . \]
To prove that $\ker(L)$ is complemented in the tangent space $E = T_{\eta_i} \Omega_D^1(\ST_i)$, we construct an explicit projection. Since $L$ is a surjective linear functional, there exists some vector $\nu \in E$ such that $L(\nu) = 1$. We define the continuous linear map $P: E \to E$ by:
\[ P(\eta) = L(\eta)\nu. \]
By construction, $P$ is a projection onto the one-dimensional subspace $M = \text{span}\{\nu\}$, since:
\[ P(P(\eta)) = L(L(\eta)\nu)\nu = L(\eta)L(\nu)\nu = L(\eta)\nu = P(\eta). \]
Any element $\eta \in E$ can be decomposed uniquely as:
\[ \eta = (\eta - P(\eta)) + P(\eta). \]
We observe the following:
\begin{enumerate}
    \item The first term $(\eta - P(\eta))$ lies in $\ker(L)$, because $L(\eta - L(\eta)\nu) = L(\eta) - L(\eta)L(\nu) = 0$.
    \item The second term $P(\eta)$ lies in $M = \text{span}\{\nu\}$.
    \item The intersection $\ker(L) \cap M$ is trivial, as $L(c\nu)  = 0$ implies $c=0$.
\end{enumerate}
This induces the topological direct sum $E = \ker(L) \oplus M$. Since $M \cong \mathbb{R}$, it is a closed subspace, and so $\ker(d_{\eta_i} g_i)$ is complemented in $T_{\eta_i} \Omega_D^1(\ST_i)$.




\subsection{Lagrange Multipliers}\label{app:LM}

The definition of a critical point $\eta_i$ for an objective function $F: \Omega_D^1(\ST_i) \to \mathbb{R}$, subject to the constraint $g_i(\eta_i) = h_i$, is given by the condition:
\begin{align*}
    d_{\eta_i} F(\nu) = 0 \quad \text{for all} \quad \nu \in T_{\eta_i} g_i^{-1}(h_i).
\end{align*}
Assuming $\eta_i$ is a regular value of the constraint $g_i$, the tangent space to the constraint manifold is precisely the kernel of the derivative: $T_{\eta_i} g_i^{-1}(h_i) = \ker(d_{\eta_i} g_i)$. Thus, the condition for a critical point is that $d_{\eta_i} F$ vanishes on $\ker(d_{\eta_i} g_i)$.

We wish to show that this is equivalent to the existence of a Lagrange multiplier $\lambda \in \mathbb{R}$ such that:
\begin{align*}
    d_{\eta_i} F + \lambda d_{\eta_i} g_i = 0 \quad \text{on the whole space } T_{\eta_i} \Omega_D^1(\ST_i).
\end{align*}
As established in Appendix \ref{app:Complemented}, since $d_{\eta_i} g_i$ is a surjective linear map to a one-dimensional codomain, its kernel is complemented. We can decompose the tangent space $E = T_{\eta_i} \Omega_D^1(\ST_i)$ into a direct sum:
\[ E = \ker(d_{\eta_i} g_i) \oplus \text{span}\{\nu_0\} \]
where $\nu_0 \in E$ is a vector chosen such that $d_{\eta_i} g_i(\nu_0) = 1$. 

Any arbitrary vector $\nu \in E$ can be uniquely decomposed as $\nu = \nu_k + c\nu_0$, where $\nu_k \in \ker(d_{\eta_i} g_i)$ and $c \in \mathbb{R}$. Note that by applying $d_{\eta_i} g_i$ to both sides, we find that the scalar $c$ is exactly $d_{\eta_i} g_i(\nu)$.

Now, consider the action of the derivative of the objective function on $\nu$:
\begin{align*}
    d_{\eta_i} F(\nu) &= d_{\eta_i} F(\nu_k + c\nu_0) \\
    &= d_{\eta_i} F(\nu_k) + c d_{\eta_i} F(\nu_0).
\end{align*}
By our definition of a critical point, $d_{\eta_i} F(\nu_k) = 0$. Substituting $c = d_{\eta_i} g_i(\nu)$, the expression simplifies to:
\[ d_{\eta_i} F(\nu) = d_{\eta_i} g_i(\nu) \cdot d_{\eta_i} F(\nu_0). \]

If we define the constant $\lambda = -d_{\eta_i} F(\nu_0)$, then for any $\nu \in E$, we have:
\[ d_{\eta_i} F(\nu) + \lambda d_{\eta_i} g_i(\nu) = 0. \]
This confirms the existence of a Lagrange multiplier at a critical point. Going the other way, we need to show that \(d_{\eta_i} F + \lambda d_{\eta_i} g_i = 0\) implies \(\eta_i \) is a critical point. This is trivial since \(\nu \in \ker (d_{\eta_i} g_i )\) implies that \(d_{\eta_i} F(\nu) = 0\). 





\section{Existence of Solutions} \label{append:exist}


With \(n=1\) there are no internal interfaces, so the diffeomorphism \(f\in\I\) plays no role. Hence, we take \(f=\mathrm{id}\) and drop partition subscripts. Otherwise, the notation in this section is consistent with the remainder of the paper. 



 We look for a solution \(\beta\) to
\begin{align}  \label{problem:a_only}
    \begin{gathered} 
    \underset{\beta  \in  \mathcal{C} ( \ST ) }{ \text{min} }  \| \beta   \|_{L^2}^2  \, , \\
    H_{\mathscr{A}(\ST)} (\beta,  \ST) = h \, , \qquad \int_{T_{j}} \beta  = \psi_{j} \quad j = 1,2, \dots, \ell \, ,
    \end{gathered}
\end{align}
for an Amperian primitive \(a \in \mathscr{A}(\ST)\), where \(\{[T_1], \dots, [T_\ell]\}\) is a basis of \(H_2(\ST,\partial\ST)\) whose boundaries coincide with the usual \([\tau_j]\). Fixing an Amperian gauge makes the helicity a function of \(\beta \) alone. The pressure is irrelevant and is omitted.


\begin{lemma}\label{lemma:c1}
    The following set is non-empty:
    \[
    \mathscr Q(\psi_1,\dots,\psi_\ell,h,\ST)
    :=\Big\{\beta \in\mathcal C(\ST)\ :\ H_{\mathscr{A}(\ST)}(\beta,\ST)=h,\ \langle[ \beta ],[T_j]\rangle=\psi_j\ \forall j\Big\},
    \]
    with the flux conditions absent when \(\ell=0\). 
\end{lemma}
\begin{proof}
As $\ST$ is a non-empty, compact 3-manifold there exists an open 3-dimensional ball $\B(r) \subset\ST$ of some radius $r>0$. Set \(\widehat{\ST}\coloneqq \ST\setminus \B(r)\). Consider \(\ell>0\) and a closed \(w \in \Omega^1(\partial \ST)\), so
\begin{align*}
    [w] = \sum_{j=1}^\ell c^1_j [\tau_j^*] + c^2_j [\sigma_j^*] \, .
\end{align*}
We select \(c^1_j = \psi_j\) and \(c^2_j = 0\). In Euclidean coordinates \((x_1, x_2, x_3)\), there are \(\widehat g_j \in \Omega^0(\partial \ST)\) for \(j=1,2,3\) such that
\begin{align*}
    w = \sum_{j=1}^3 \widehat g_{j}\,dx_j \, .
\end{align*}
Since $\partial {\ST}$ is compact in \(\R^3\) and contained in $\ST$, we may extend $\widehat g_j$ to a smooth function \(g_j\) on $\ST$ (\(j=1,2,3\)) by Whitney's extension theorem \cite{whitney1992analytic}. 

Let \(\chi\in \Omega^0 (\ST)\) be a bump function with \(\chi=0\) on \(\B(r_0)\) (with \(r_0 < r\)) and \(\chi=1\) on \(\widehat{\ST}\). Define
\[
a_{\ST}\coloneqq \chi \sum_{j=1}^3 g_{j}\,dx_j   \in\Omega^1(\ST),\qquad \text{and} \qquad \beta_{\ST}\coloneqq d a_{\ST} \in\Omega^2(\ST) \, ,
\]
giving \(a_{\ST}|_{\B(r_0)}= 0\). Additionally, \(\beta _{\ST}\) is exact by construction and vanishes when pulled back to \(\partial\ST\)
\begin{align*}
    t(\beta _{\ST}) = t(d a_{\ST}) = d t(a_{\ST}) = d w = 0 \, .
\end{align*}
Hence \(\beta _{\ST} \in \mathcal{C}(\ST)\) and satisfies the flux constraint:
\begin{align*}
    \int_{T_j} \beta _{\ST} = \int_{\tau_j} \mathcal{J}^* a_{\ST} = \int_{\tau_j} w = \psi_j \, .
\end{align*}
Additionally, 
\begin{align*}
    \int_{\sigma_j} \mathcal{J}^* a_{\ST} = \int_{\sigma_j} w = 0 \, ,
\end{align*}
indicating that \(a_{\ST} \in \mathscr{A}(\ST)\). 



 There exists a field \(\beta _\B \in \mathcal{C}(\ST)\) that is only non-zero within \(\B(r_0)\) and able to achieve any value for helicity (\cite{PhysRevLett.117.274501} supplementary material). Such \(\beta _\B\) were constructed explicitly by Kedia, Foster, Dennis, and Irvine. These authors show that a smooth potential \(a_\B\) may be constructed from two linked trefoil knots to achieve any helicity, and that \(a_\B, \beta _\B = 0\) outside of a bounding ball and therefore the flux of \(\beta _{\ST} + \beta _\B \in \mathcal{C}(\ST)\) is given by \(\beta _{\ST}\). As \(a_\B = 0\) on the boundary this implies \(a_\B \in \mathscr{A}(\ST)\). 

The helicity for \(\beta _{\ST} + \beta _\B\) is:
\begin{align*}
    H_{\mathscr{A}(\ST)}( \beta_\B + \beta_\ST , \ST) = H_{\mathscr{A}(\ST)} (\beta_\ST , \ST \setminus \B(r_0) ) +  H_{\mathscr{A}(\ST)} ( \beta_\B , \B(r_0) ) \, .
\end{align*}
Immediately we see that, after fixing a selection for \({\beta}_\ST\) with the desired flux, we may select a \(  \beta_\B \) to generate any value for \(H_{\mathscr{A}(\ST)}(\beta , \ST) \), without affecting the flux. We now have a way to construct any flux and helicity with a closed Dirichlet 2-form \(\beta  \in \mathcal{C} (\ST)\).



If \(\ell = 0\) then there is no flux and one only needs to reconstruct the helicity. The above arguments tell us that \(\mathscr{Q}( \psi_1, \dots, \psi_\ell , h, \ST )\) is non-empty for any \(\ell\). 
\end{proof}


We now apply the direct method in the calculus of variations to prove the existence of a minimiser of problem \ref{problem:a_only}, and hence the existence of a stationary point.

By Lemma \ref{lemma:c1}, we may choose
\[
\beta \in \mathscr{Q}(\psi_1,\dots,\psi_\ell,h,\ST).
\]
Moreover, there exists a potential for \(\beta\) given by
\[
a=\eta+\Gamma \in \mathscr{A}(\ST),
\]
where \(d\Gamma \in \mathcal{H}^2_D (\ST)\) is the unique harmonic term and \(\eta \in \Omega^1_D(\ST)\). Given that \(\eta\) is zero when pulled back to the boundary, only the \(\Gamma\) term contributes to \(\langle [a], [\sigma_j]\rangle\). Therefore, \(\Gamma\) is solely responsible for ensuring \(a \in \mathscr{A}(\ST)\). Without loss of generality, we fix \(\Gamma\). 


Orthogonality of the HMF decomposition (see Appendix \ref{append:Hodge}) implies 
\begin{align*}
    \| \beta  \|_{L^2}^2 = \|d \eta \|_{L^2}^2 + \| d \Gamma \|_{L^2}^2 \, .
\end{align*}
The last term is determined by a choice of \(\psi_1, \dots, \psi_\ell\). Hence, problem \ref{problem:a_only} descends to an equivalent problem in \(d \eta\), namely, to determine if \(\|d \eta \|_{L^2}^2\) has any critical points with the constraint \(H_{\mathscr{A}(\ST)}( \beta , \ST) = h\).





It is convenient to enlarge the admissible space of $2$-forms containing $d\eta$. We therefore formulate the variational problem in a Hilbert space \(L^2\Omega_D^2(\ST)\), obtained as the completion of $\Omega_D^2(\ST)$ with respect to the $L^2$ inner product.

For $1$-forms, we equip $\Omega_D^1(\ST)$ with the $H^1$ norm
\[
\|\omega\|_{H^1}^2 \coloneqq \|\omega\|_{L^2}^2 + \|D\omega\|_{L^2}^2,
\qquad \omega \in \Omega_D^1(\ST),
\]
where $D$ denotes the weak covariant derivative used by Schwarz \cite{schwarz-1995}. The corresponding $H^1$ inner product is
\[
\langle \omega,\omega_0\rangle_{H^1}
\coloneqq
\langle \omega,\omega_0\rangle_{L^2}
+
\langle D\omega, D\omega_0\rangle_{L^2},
\qquad \omega,\omega_0 \in \Omega_D^1(\ST).
\]

More generally, for the definitions of the spaces $H^k\Omega_D^1(\ST)$, $H^k_{\mathrm{loc}}\Omega_D^1( \text{int}( \ST ) )$, and $C^k\Omega_D^1(\ST)$ on a manifold with boundary, together with their associated norms and the tangential trace operator $t$, we refer the reader to \cite{schwarz-1995}. We now turn to the minimisation problem
\begin{align}\label{prob:inter} 
    \begin{gathered} 
    \underset{  d \eta \in L^2 \Omega^2_D(\ST)}{ \text{min}}   \| d \eta  \|_{L^2}^2  \, , \qquad 
    \hat H ( \eta) = h  \, ,
    \end{gathered}
\end{align}
where 
\begin{align*}
    \hat H  : H^1 \Omega^1_D(\ST) \to \R \qquad \text{is given by} \qquad \hat H  ( \eta) = \int_{\ST} (\eta + \Gamma) \wedge d (\eta + \Gamma)  \, ,
\end{align*}
where the helicity \(\hat H\) is dependent upon \( d \eta\) irrespective of the primitive \(\eta\) chosen. This is in precisely the same form as the problem discussed by Gerner who used the direct method in calculus of variations to show that there exists a minimiser \(d \eta \) (Theorem 2.1 \cite{exist_gerner}). Additionally, as the Euclidean metric is real analytic then for non-zero helicity \(d \eta\) is real analytic on \(\text{int}(\ST_i)\) (Corollary 2.1 \cite{exist_gerner}). We provide an exposition for these results below.





For a minimising sequence, \( d \eta_j \) in \(j\), we are aware that \(L^2 \Omega^2_D(\ST_i)\) is complete, so \(d \eta_j \to d \eta\) weakly in \(L^2\) as \(j \to \infty\) (up to a subsequence). Consider the map \(\star d \eta \mapsto \eta \in H^1 \Omega^1_D(\ST)\) as defined by Gerner which is a continuous linear operator (Lemma 3.4 in \cite{exist_gerner}). Hence, \( \eta_j \to  \eta\) converges weakly in \(H^1\). By the Rellich–Kondrachov theorem there is an embedding into  \(L^2 \Omega^1_D(\ST)\) where a subsequence of \(  \eta_j\) converges strongly to \(\eta\). So
\begin{align*}
    |\hat H  (   \eta_j ) - \hat H  (   \eta )| &= \bigg| \int_{\ST} \left(   \eta_j + \Gamma \right)  \wedge \beta _j
    - 
    \int_{\ST} \left(   \eta +  \Gamma \right)  \wedge \beta  \bigg| 
    \\
      &= \bigg| \int_{\ST} \left(   \eta_j - \eta \right)  \wedge \beta _j 
    - 
    \int_{\ST} \left(   \eta +  \Gamma \right)  \wedge  (\beta - \beta_j) \bigg| \, .
\end{align*}
The first term tends to zero because \( \eta_j - \eta \to 0\) strongly in \(L^2\) and \( \| \beta_j \|_{L^2}\) is bounded. The second term goes to zero because \(\beta - \beta_j\) goes to zero in \(L^2\) when tested against the test function \(\eta +  \Gamma  \) by weak convergence. We see that \(|\hat H  (   \eta_j ) \to \hat H  (   \eta )| \), namely \(\hat H\) is sequentially weakly continuous. This implies that \(\hat H^{-1}\) is weakly closed.


The objective \(\| d \eta  \|_{L^2}^2\) is weakly lower semicontinuous and trivially coercive. It follows from the direct method on Banach spaces that there exists a minimiser \(d \eta\).


We now determine the Euler-Lagrange equations for problem \eqref{prob:inter}. Note that a Banach space is automatically Fréchet, so Theorem \ref{theorem:Glock} immediately applies. First we show that the corresponding helicity is Fr\'echet differentiable. So,
\begin{align*}
    \hat H ( \eta + \eta') 
    &= \hat H ( \eta ) +
    2 \int_{\ST} \beta \wedge  \eta'  
    +
    \int_{\ST}  \eta'  \wedge d \eta' 
\end{align*}
We see that 
\begin{align*}
    \partial_{\eta} \hat H(  \eta'  ) \coloneqq 2 \int_{\ST} \beta \wedge  \eta'  
\end{align*}
is a linear, continuous operator in \(\eta'\). For \(\partial_{\eta} \hat H(  \eta'  )\) to be a Fr\'echet derivative requires
\begin{align*}
    \Big| \int_{\ST}  \eta'  \wedge d \eta' \Big|  / \| \eta' \|_{H^1} \to 0 \qquad \text{as} \qquad  \| \eta' \|_{H^1} \to 0 \, . 
\end{align*}
And we see that
\begin{align*}
    \Big| \int_{\ST}  \eta'  \wedge d \eta' \Big| = \langle \eta'  , \star d \eta' \rangle
    \leq
    \| \eta' \|_{L^2} \, \| \star d \eta'  \|_{L^2}
    =
    \| \eta' \|_{L^2} \, \|  d \eta'  \|_{L^2} \, .
\end{align*}
By definition \( \| \eta' \|_{L^2} \leq C \| \eta' \|_{H^1}\) for some \(C\). And the fact that \(d\) is a first-order differential operator, so on a compact smooth volume \( \| d \eta' \|_{L^2} \leq C \| \eta' \|_{H^1}\) for some \(C\). Giving,
\begin{align*}
    0 < \Big| \int_{\ST}  \eta'  \wedge d \eta' \Big|  / \| \eta' \|_{H^1} \leq \| \eta' \|_{L^2} \, \|  d \eta'  \|_{L^2}  / \| \eta' \|_{H^1}  \leq C \| \eta' \|_{H^1} \, .
\end{align*}
The central terms evidently go to 0 as \(\| \eta' \|_{H^1} \to 0 \). Hence, \(\partial_{\eta} \hat H(  \eta'  )\) defines a Fr\'echet derivative. 



Let \(\chi_{\B}\) be the characteristic function for a 3-dimensional ball \(\B(r_0)\) of radius \(r_0\) contained in the interior of \(\ST\). As \(\ST\) is a compact Riemannian 3-manifold, there exists such a ball. Additionally, let \(\chi_k : \ST \to \R\) for each \(k \in \mathbb{N}\) represent an element in \( \Omega^0 (\ST)\) with \(\chi_k |_{\partial \ST} = 0\). Select \(\chi_k\) such that \(\chi_k \to \chi_{\B}\) pointwise as \(k \to \infty\). 



Consider the direction \(\hat \eta'_k = \chi_k  \, \star \beta\). Then, as \(H^1 \Omega^1(\ST)\) is dense in \(L^2 \Omega^1(\ST)\) we can select a \(\eta'_{k,q} \to  \hat \eta'_k \) in an \(L^2\) topology as \(q \to \infty\) such that \(\eta'_{k,q} \in H^1 \Omega^1_D (\ST)\):
\begin{align*}
    d_{\eta} \hat H (\eta'_{k,q}) 
    \xrightarrow{q \to \infty}
    2\int_{\ST} \! \chi_k  \, \star \beta  \wedge \beta   \xrightarrow{k \to \infty} 2 \int_{\chi_{\B}} \! \star \beta  \wedge \beta 
    = 2 \| \beta |_{\B(r_0)} \|_{L^2}^2   
     \, .
\end{align*}
So long as \(\| \beta \|_{L^2} > 0\) there exists a choice of ball \(\B(r_0)\) so that \( \| \beta |_{\B(r_0)} \|_{L^2}^2  > 0 \), we select such a ball. 


So, by an epsilon-delta argument, there are \(k_0,q_0\) such that for \(k > k_0, q>q_0\) we have \(d_{\eta} \hat H (\eta'_{k,q}) > 0\) and surjectivity for any \( \| \beta \|_{L^2}  \neq 0\). This means that any \( h \) is a regular value for the relative helicity constraint \(\hat H (\eta) =  h \) (Theorem \ref{theorem:Glock}) if \(\| \beta \|_{L^2} > 0\). One usually checks the abnormal case \(\| \beta \|_{L^2} = 0\) using higher derivatives \cite{tret1984necessary}. Assuming that \(\| \beta \|_{L^2} > 0\) the regular Lagrange multiplier condition provides necessary and sufficient conditions for a critical point to the original problem \eqref{prob:inter} (see Theorem 43.C (Ljusternik (1934)) in \cite{Lys} for constraint qualifications),
\begin{align*}
    \underset{\substack{\eta \in  H^1 \Omega^1_D(\ST) \\  \mu \in \R} }{ \text{c. p.} }  \left(  \frac{1}{2} \Big\|d \eta + \sum_j  \psi_{j}  T_{j}^*  \Big\|_{L^2}^2   - \frac{\mu}{2} \big(   \hat H (\eta) -   h       \big) \right) \, .
\end{align*}
Proceeding with the Gâteaux derivative, namely, \(\beta_0 = \beta \), \(\partial_t \beta_t |_{t=0} =  \beta' \), similarly for \(\mu\) we have that stationary points of the above are the derivative with respect to \(\mu\) in the direction \(\mu'\),
\begin{align*}
    0 &=    - \frac{\mu'}{2} \big(   \hat H (\eta) -   h       \big) \, ,
\end{align*}
for all \(\mu'\) which implies that the helicity constraint holds. And the derivative with respect to \(\eta\) in the direction \(\eta'\) is
\begin{align*}
      0 &= \int_{\ST} d \eta' \wedge \star \beta   - \frac{\mu}{2} \int_{\ST} \eta'  \wedge \beta + a \wedge d \eta'  \\
      0 &= \int_{\ST} \eta' \wedge (d\star \beta  - \mu \beta ) \, ,
\end{align*}
for all \(\eta'\), where we have used the divergence theorem and that \(\beta = da\) to simplify this condition. The condition implies that
\begin{align*}
    \langle   (d\star \beta  - \mu \beta )   , \star \eta' \rangle = 0 \, . 
\end{align*}
On the interior of \(\ST\) we see that \( d\star \beta  - \mu \beta \) projected onto any element of \(\Omega^2( \text{int} (\ST ))\) is zero. Hence, in \(\text{int} (\ST)\) we have that \( d\star \beta  - \mu \beta  = 0\) in the sense of distributions with \(\beta \in H^0 \Omega^2 (\ST)\). 

We may apply \((d\star   - \mu   )\), as a distributional operator, to \( d\star \beta  - \mu \beta  = 0\). Namely,
\begin{align*}
    0 = (d\star   - \mu   )^2 \beta &= (d\star d \star   - 2 \mu d\star + \mu^2   ) \beta =  (d\star d \star   -  \mu^2   ) \beta = (\Delta - \mu^2 ) \beta  \, ,
\end{align*}
where \(\Delta \beta = (d \delta + \delta d) \beta = d \delta \beta= d \star d \star  \beta \) and \(\Delta\) is the Hodge Laplacian on 2-forms. A Hodge decomposition gives \(\beta = d ( \eta + \Gamma) \) where \(\Gamma\) is known and smooth, and \(d \Gamma \) is harmonic, so
\begin{align*}
    (\Delta - \mu^2 ) d \eta &=  \mu^2   d \Gamma \, . 
\end{align*}
Since the right-hand side of this equation is smooth, interior elliptic regularity theory implies that \(d \eta \in H^k_{loc} (  \text{int}(\ST)  )\) for all \(k\), namely, \(d \eta \) is smooth within the domain interior \cite{gilbarg1998elliptic}. 



Hence, there exists a minimiser \(\eta\) to problem \eqref{prob:inter}, giving a \(\beta = d (\eta + \Gamma)\) and if the minimiser satisfies \(\| \beta \|_{L^2} > 0 \) then \(d \eta\) and therefore \(\beta\) must also be smooth within the domain interior. We note that if at least one of the flux values \(\psi_1, \dots, \psi_\ell\) is non-zero this implies that (by \(L^2\) orthogonality of the Hodge decomposition) \(\| \beta \|_{L^2}^2 \geq \|d \Gamma \|_{L^2}^2 > 0\). 







\end{document}
